% La radioactivité et la loi de décroissance — Physique Tle C — Cours signé OnBuch+
% Programme officiel MINESEC. Compiler avec : tectonic P19-radioactivite.tex
\newcommand{\DOCMATIERE}{Physique}
\newcommand{\DOCNIVEAU}{Tle C}
\newcommand{\DOCMODULE}{Module 4 : Onde, matière et transformations nucléaires}
\newcommand{\DOCLECON}{P19}
\newcommand{\DOCTITRE}{La radioactivité et la loi de décroissance}
% OnBuch+ — préambule « cours signés » ENRICHI (compilé avec tectonic / XeLaTeX)
% Système visuel « Pop » d'OnBuch+ : fond crème (jamais blanc), bordures encre
% épaisses, ombres « dures » décalées, titres Archivo Black en capitales.
% Version MATHÉMATIQUES : graphes (pgfplots/tikz), boîtes pédagogiques
% (définition, propriété, méthode, attention, le savais-tu, exercice résolu,
% exercice, TP, bilan), page de garde et signature OnBuch+ ancrée en bas.
%
% Macros attendues dans main.tex AVANT \input{preamble} :
%   \DOCMATIERE  \DOCNIVEAU  \DOCMODULE  \DOCLECON (numéro)  \DOCTITRE
\documentclass[11pt]{article}

\usepackage{fontspec}
\usepackage[french]{babel}
\usepackage[a4paper,margin=2cm,top=3.4cm,bottom=2.8cm,headheight=1.4cm,footskip=1.3cm]{geometry}
\usepackage{amsmath,amssymb,amsfonts}
\usepackage{mathtools} % \xmapsto, \coloneqq... souvent utilisés par les modèles
\usepackage{cancel} % simplifications barrées (bilans de forces)
% \abs{z} (module, valeur absolue) et \norm{u} (norme), utilisés par les modèles
\providecommand{\abs}{}\let\abs\relax\DeclarePairedDelimiter{\abs}{\lvert}{\rvert}
\providecommand{\norm}{}\let\norm\relax\DeclarePairedDelimiter{\norm}{\lVert}{\rVert}
\usepackage[table]{xcolor} % [table] : fournit aussi \rowcolors (alternance de lignes)
\usepackage{pagecolor}
\usepackage{tikz}
\usetikzlibrary{arrows.meta,positioning,calc,shapes.geometric,shapes.misc,shapes.arrows,decorations.pathmorphing,decorations.markings,patterns,shadows,mindmap,trees,fit,backgrounds,matrix,intersections,angles,quotes}
\usepackage{pgfplots}
\usepackage[version=4]{mhchem} % équations nucléaires \ce{^{235}_{92}U}
\usepackage[european,siunitx]{circuitikz} % schémas électriques
\pgfplotsset{compat=1.18}
\usepgfplotslibrary{fillbetween} % aires entre courbes (calcul intégral)
\usepackage{enumitem}
\usepackage{fancyhdr}
% [most] charge toutes les bibliothèques tcolorbox (listings, minted, raster,
% documentation...) et gonfle inutilement la mémoire TeX sur un document long
% (~300 boîtes) : on ne charge que ce qui est réellement utilisé.
\usepackage[skins,breakable]{tcolorbox}
\usepackage{graphicx}
\usepackage{colortbl}
\usepackage{booktabs}
\usepackage{tabularx}
\usepackage{diagbox} % case d'en-tête diagonale des tableaux à double entrée
% Colonnes extensibles centrée / gauche / droite (C, L, R), souvent utilisées par les modèles
\newcolumntype{C}{>{\centering\arraybackslash}X}
\newcolumntype{L}{>{\raggedright\arraybackslash}X}
\newcolumntype{R}{>{\raggedleft\arraybackslash}X}
\usepackage{array}
\usepackage{multicol}
\usepackage{siunitx}
% parse-numbers=false : accepte aussi \SI{2,5 \times 10^{-3}}{...}
\sisetup{locale=FR,output-decimal-marker={,},group-minimum-digits=4,parse-numbers=false}
\DeclareSIUnit\ohm{\ensuremath{\symup{\Omega}}} % U+2126 absent de la police
\DeclareSIUnit\division{div} % réglages d'oscilloscope (ms/div, V/div)
\usepackage{qrcode}
% \degree : commande fréquemment utilisée par les modèles pour un angle ou une
% température en dehors de \SI{}{} (non fournie par les paquets chargés).
\providecommand{\degree}{^{\circ}}
% \qed (fin de démonstration), utilisé par les modèles sans amsthm
\providecommand{\qed}{\hfill$\square$}
% \barre : double barre des valeurs interdites dans un tableau de variations (tkz-tab)
\providecommand{\barre}{\ensuremath{\|}}
% environnement proof (amsthm non chargé) utilisé par les modèles
\ifdefined\proof\else\newenvironment{proof}[1][Démonstration]{\par\noindent\textit{#1.}\ }{\qed\par}\fi
\usepackage{lastpage}
\usepackage{hyperref}
\hypersetup{hidelinks}

% ============================================================
% Palette « Pop » OnBuch+ — mêmes hex que la classe P (plus_ui.dart)
% ============================================================
\definecolor{poporange}{HTML}{F59321}
\definecolor{popdark}{HTML}{E07A0C}
\definecolor{popcream}{HTML}{FFF6E8}
\definecolor{popink}{HTML}{1C1714}
\definecolor{poppurple}{HTML}{7A5AE0}
\definecolor{popgreen}{HTML}{2E9E6A}
\definecolor{poppink}{HTML}{E0457B}
\definecolor{popblue}{HTML}{2F7DE1}
\definecolor{popgold}{HTML}{C98A12}
\definecolor{popmuted}{HTML}{8A7F76}
\definecolor{popline}{HTML}{EADFD2}
\definecolor{popgray}{HTML}{8A7F76}
\colorlet{popgrey}{popgray}
% Alias tolérants (le modèle nomme parfois les couleurs différemment)
\colorlet{popwhite}{white}
\colorlet{popblack}{popink}
\colorlet{popred}{poppink}
\colorlet{popyellow}{popgold}
% Teintes claires pour fonds de tableaux / zones
\colorlet{popblueL}{popblue!10!white}
\colorlet{poporangeL}{poporange!14!white}
\colorlet{popgreenL}{popgreen!12!white}
\colorlet{poppurpleL}{poppurple!12!white}
\colorlet{poppinkL}{poppink!12!white}
\colorlet{popgoldL}{popgold!14!white}

\pagecolor{popcream}

% ============================================================
% Typographie
% ============================================================
\defaultfontfeatures{Ligatures=TeX}
\setmainfont{PlusJakartaSans-Medium.ttf}[Path=../../../fonts/,BoldFont=PlusJakartaSans-Bold.ttf]
% Maths en sans-serif (Fira Math) avec chiffres et lettres droites de Plus
% Jakarta, pour que les formules se fondent dans le texte.
\usepackage[math-style=ISO,bold-style=ISO]{unicode-math}
\setmathfont{FiraMath-Regular.otf}[Path=../../../fonts/]
\setmathfont{PlusJakartaSans-Medium.ttf}[Path=../../../fonts/,range={up/num,up/{Latin,latin},\mathup}]
\usepackage{icomma} % virgule décimale sans espace en mode math
% Caractères mathématiques Unicode absents des polices de texte (Plus Jakarta,
% Archivo Black) : rendus par leur équivalent mathématique, en texte comme en
% formule — sinon ils s'affichent en carré vide (ex. « Arithmétique dans ℤ »).
\usepackage{newunicodechar}
\newunicodechar{ℝ}{\ensuremath{\mathbb{R}}}
\newunicodechar{ℤ}{\ensuremath{\mathbb{Z}}}
\newunicodechar{ℂ}{\ensuremath{\mathbb{C}}}
\newunicodechar{ℕ}{\ensuremath{\mathbb{N}}}
\newunicodechar{ℚ}{\ensuremath{\mathbb{Q}}}
\newunicodechar{𝔻}{\ensuremath{\mathbb{D}}}
\newunicodechar{α}{\ensuremath{\alpha}}
\newunicodechar{β}{\ensuremath{\beta}}
\newunicodechar{γ}{\ensuremath{\gamma}}
\newunicodechar{δ}{\ensuremath{\delta}}
\newunicodechar{ε}{\ensuremath{\varepsilon}}
\newunicodechar{θ}{\ensuremath{\theta}}
\newunicodechar{λ}{\ensuremath{\lambda}}
\newunicodechar{μ}{\ensuremath{\mu}}
\newunicodechar{σ}{\ensuremath{\sigma}}
\newunicodechar{τ}{\ensuremath{\tau}}
\newunicodechar{φ}{\ensuremath{\varphi}}
\newunicodechar{ψ}{\ensuremath{\psi}}
\newunicodechar{ω}{\ensuremath{\omega}}
\newunicodechar{Σ}{\ensuremath{\Sigma}}
% majuscules produites par \MakeUppercase dans les titres (α-aminés -> Α)
\newunicodechar{Α}{\ensuremath{\alpha}}
\newunicodechar{Β}{\ensuremath{\beta}}
\newunicodechar{Γ}{\ensuremath{\Gamma}}
\newunicodechar{Δ}{\ensuremath{\Delta}}
\newunicodechar{Ε}{\ensuremath{\varepsilon}}
\newunicodechar{Θ}{\ensuremath{\Theta}}
\newunicodechar{Λ}{\ensuremath{\Lambda}}
\newunicodechar{Μ}{\ensuremath{\mu}}
\newunicodechar{Τ}{\ensuremath{\tau}}
\newunicodechar{Φ}{\ensuremath{\Phi}}
\newunicodechar{Ψ}{\ensuremath{\Psi}}
\newunicodechar{Ω}{\ensuremath{\Omega}}
\newunicodechar{↦}{\ensuremath{\mapsto}}
\newunicodechar{∈}{\ensuremath{\in}}
\newunicodechar{∉}{\ensuremath{\notin}}
\newunicodechar{∧}{\ensuremath{\wedge}}
\newunicodechar{⇔}{\ensuremath{\Leftrightarrow}}
\newunicodechar{⟺}{\ensuremath{\Longleftrightarrow}}
\newunicodechar{⇒}{\ensuremath{\Rightarrow}}
\newunicodechar{⟹}{\ensuremath{\Longrightarrow}}
\newunicodechar{∘}{\ensuremath{\circ}}
\newunicodechar{∪}{\ensuremath{\cup}}
\newunicodechar{∩}{\ensuremath{\cap}}
\newunicodechar{≡}{\ensuremath{\equiv}}
\newunicodechar{⊂}{\ensuremath{\subset}}
\newunicodechar{⊥}{\ensuremath{\perp}}
\newunicodechar{∃}{\ensuremath{\exists}}
\newunicodechar{∀}{\ensuremath{\forall}}
\newunicodechar{∛}{\ensuremath{\sqrt[3]{\,}}}
\newunicodechar{∜}{\ensuremath{\sqrt[4]{\,}}}
\newunicodechar{⃗}{\ensuremath{{}^{\to}}}
\newunicodechar{ⁿ}{\ifmmode^{n}\else\textsuperscript{\ensuremath{n}}\fi}
\newunicodechar{ˣ}{\ifmmode^{x}\else\textsuperscript{\ensuremath{x}}\fi}
\newunicodechar{ᵇ}{\ifmmode^{b}\else\textsuperscript{\ensuremath{b}}\fi}
\newunicodechar{ʳ}{\ifmmode^{r}\else\textsuperscript{\ensuremath{r}}\fi}
\newunicodechar{ᵘ}{\ifmmode^{u}\else\textsuperscript{\ensuremath{u}}\fi}
\newunicodechar{ʸ}{\ifmmode^{y}\else\textsuperscript{\ensuremath{y}}\fi}
\newunicodechar{ᵃ}{\ifmmode^{a}\else\textsuperscript{\ensuremath{a}}\fi}
\newunicodechar{ᵉ}{\ifmmode^{e}\else\textsuperscript{\ensuremath{e}}\fi}
\newunicodechar{ᵏ}{\ifmmode^{k}\else\textsuperscript{\ensuremath{k}}\fi}
\newunicodechar{ᵖ}{\ifmmode^{p}\else\textsuperscript{\ensuremath{p}}\fi}
\newunicodechar{ᴺ}{\ifmmode^{N}\else\textsuperscript{\ensuremath{N}}\fi}
\newunicodechar{ᶜ}{\ifmmode^{c}\else\textsuperscript{\ensuremath{c}}\fi}
\newunicodechar{ʰ}{\ifmmode^{h}\else\textsuperscript{\ensuremath{h}}\fi}
\newunicodechar{ᵠ}{\ifmmode^{q}\else\textsuperscript{\ensuremath{q}}\fi}
\newunicodechar{ₙ}{\ifmmode_{n}\else\textsubscript{\ensuremath{n}}\fi}
\newunicodechar{ₐ}{\ifmmode_{a}\else\textsubscript{\ensuremath{a}}\fi}
\newunicodechar{ᵢ}{\ifmmode_{i}\else\textsubscript{\ensuremath{i}}\fi}
\newunicodechar{ᵣ}{\ifmmode_{r}\else\textsubscript{\ensuremath{r}}\fi}
\newunicodechar{ₖ}{\ifmmode_{k}\else\textsubscript{\ensuremath{k}}\fi}
\newunicodechar{ᵦ}{\ifmmode_{\beta}\else\textsubscript{\ensuremath{\beta}}\fi}
\newunicodechar{ₓ}{\ifmmode_{x}\else\textsubscript{\ensuremath{x}}\fi}
\newfontfamily\popdisplay{ArchivoBlack-Regular.ttf}[Path=../../../fonts/]
\newfontfamily\poplogo{SpaceGrotesk-Bold.ttf}[Path=../../../fonts/]
\newfontfamily\popsemi{PlusJakartaSans-SemiBold.ttf}[Path=../../../fonts/]
\newfontfamily\popxbold{PlusJakartaSans-ExtraBold.ttf}[Path=../../../fonts/]
\newfontfamily\popxboldls{PlusJakartaSans-ExtraBold.ttf}[Path=../../../fonts/,LetterSpace=3.0]

\setlist[enumerate]{leftmargin=1.7em,itemsep=5pt,topsep=4pt}
\setlist[itemize]{leftmargin=1.7em,itemsep=4pt,topsep=4pt,label={\color{poporange}\textbullet}}
\setlength{\parindent}{0pt}
\setlength{\parskip}{5pt}
\renewcommand{\arraystretch}{1.25}

% pgfplots : style Pop par défaut
\pgfplotsset{
  every axis/.append style={axis line style={popink,line width=1pt},tick style={popink},
    grid style={popline,line width=0.5pt},label style={font=\small},tick label style={font=\footnotesize}},
  popcurve/.style={poporange,line width=1.8pt,no markers,smooth},
}
% Même style disponible dans un \draw TikZ simple (hors pgfplots)
\tikzset{popcurve/.style={poporange,line width=1.8pt}}
\tikzset{popgrid/.style={popline,very thin}}
\colorlet{popcurve}{poporange} % le nom du style est parfois utilisé comme couleur

% ============================================================
% Pill / squiggle
% ============================================================
\newcommand{\poppill}[3]{%
  \tikz[baseline=-0.55ex]\node[draw=popink,line width=1.2pt,rounded corners=2.6pt,%
    fill=#2,inner xsep=6.5pt,inner ysep=3pt]{%
    \popxboldls\fontsize{7}{7}\selectfont\color{#3}\MakeUppercase{#1}};%
}
\newcommand{\popsquiggle}[1]{%
  \tikz[baseline=-0.4ex]{
    \draw[#1,line width=2.6pt,line cap=round]
      plot [smooth] coordinates {(0,0.09) (0.34,0.01) (0.68,0.21) (1.02,0.01) (1.36,0.10)};
  }%
}
% Mot-clé surligné (marqueur orange sous le texte)
\newcommand{\cle}[1]{{\popxbold\color{popdark}#1}}

% ============================================================
% En-tête / pied
% ============================================================
\pagestyle{fancy}
\fancyhf{}
\renewcommand{\headrulewidth}{0pt}
\renewcommand{\footrulewidth}{0pt}
\lhead{\raisebox{-0.15ex}{\poplogo\fontsize{13}{13}\selectfont\color{popink}OnBuch\color{poporange}+}}
\rhead{\poppill{\DOCMATIERE\ \textbullet\ \DOCNIVEAU\ \textbullet\ Leçon \DOCLECON}{white}{popink}}
\lfoot{\popxbold\fontsize{7.6}{7.6}\selectfont\color{popmuted}\MakeUppercase{Cours signé OnBuch+}\ \textbullet\ {\popsemi\DOCTITRE}}
\rfoot{\tikz[baseline=-0.6ex]\node[rounded corners=8.5pt,fill=popink,minimum height=17pt,minimum width=17pt,inner xsep=5pt,inner sep=0]{\popxbold\fontsize{8}{8}\selectfont\color{popcream}\thepage\,/\,\pageref*{LastPage}};}

% ============================================================
% Titres
% ============================================================
\newcommand{\chaptitre}[2]{%
  \begin{tcolorbox}[enhanced,colback=poporange,colframe=popink,boxrule=2.6pt,arc=11pt,
      drop shadow=popink,left=15pt,right=15pt,top=12pt,bottom=11pt,before skip=3pt,after skip=18pt]
    \poppill{#2}{popink}{popcream}\par\vspace{9pt}
    {\raggedright\hyphenpenalty=10000\exhyphenpenalty=10000\popdisplay\fontsize{22}{24}\selectfont\color{popcream}\MakeUppercase{#1}\par}
  \end{tcolorbox}
}
% Section numérotée : pastille orange avec le numéro + titre Archivo
\newcounter{popsec}
\newcommand{\coursec}[1]{%
  \stepcounter{popsec}\setcounter{popsub}{0}%
  \par\vspace{16pt}\noindent
  \begin{minipage}{\linewidth}
  \tikz[baseline=-0.7ex]\node[draw=popink,line width=1.6pt,fill=poporange,rounded corners=4pt,minimum size=22pt,inner sep=0]{\popdisplay\fontsize{12}{12}\selectfont\color{popcream}\thepopsec};%
  \hspace{8pt}{\popdisplay\fontsize{15}{17}\selectfont\color{popink}\MakeUppercase{#1}}\par
  \vspace{4pt}\noindent{\color{poporange}\rule{36pt}{3.6pt}}
  \end{minipage}\par\nopagebreak\vspace{8pt}
}
\newcounter{popsub}
\newcommand{\courssub}[1]{%
  \stepcounter{popsub}%
  \par\vspace{9pt}\noindent{\popxbold\fontsize{11.5}{13}\selectfont\color{popink}{\color{poporange}\thepopsec.\thepopsub}\hspace{6pt}#1}\par\nopagebreak\vspace{4pt}
}

% ============================================================
% Boîtes pédagogiques — même recette : fond blanc, bordure épaisse colorée,
% ombre dure encre, badge plein. Titre optionnel [#1].
% ============================================================
\tcbset{popbase/.style={breakable,enhanced,colback=white,boxrule=2.3pt,arc=9pt,
  shadow={3pt}{-3pt}{0pt}{popink},left=14pt,right=14pt,top=11pt,bottom=11pt,before skip=11pt,after skip=15pt,
  fontupper=\popsemi\fontsize{10.6}{14.5}\selectfont\color{popink},
  fontlower=\popsemi\fontsize{10.6}{14.5}\selectfont\color{popink}}}
% Coche dessinée (le glyphe ✓ n'existe pas dans les polices du document)
\newcommand{\popcheck}[1]{\tikz[baseline=-0.5ex]\draw[#1,line width=1.6pt,line cap=round,line join=round](-0.13,0.0)--(-0.04,-0.09)--(0.13,0.1);}
\providecommand{\coche}{\popcheck{popgreen}}
\providecommand{\resultat}[1]{\par\smallskip\noindent\fcolorbox{popgreen}{popgreenL}{\parbox{\dimexpr\linewidth-2\fboxsep-2\fboxrule\relax}{\textbf{Résultat.} #1}}\par\smallskip}
\newcommand{\popboxhead}[3]{\poppill{#1}{#2}{white}\ifx\relax#3\relax\else\hspace{6pt}{\popxbold\fontsize{10.6}{12}\selectfont\color{popink}#3}\fi\par\vspace{7pt}}

\newtcolorbox{aretenir}[1][]{popbase,colframe=popgreen,before upper={\popboxhead{À retenir}{popgreen}{#1}}}
\newtcolorbox{exemplebox}[1][]{popbase,colframe=poporange,before upper={\popboxhead{Exemple}{poporange}{#1}}}
\newtcolorbox{definition}[1][]{popbase,colframe=popblue,before upper={\popboxhead{Définition}{popblue}{#1}}}
\newtcolorbox{propriete}[1][]{popbase,colframe=poppurple,before upper={\popboxhead{Propriété}{poppurple}{#1}}}
\newtcolorbox{methode}[1][]{popbase,colframe=popgold,before upper={\popboxhead{Méthode}{popgold}{#1}}}
\newtcolorbox{attention}[1][]{popbase,colframe=poppink,before upper={\popboxhead{Attention}{poppink}{#1}}}
\newtcolorbox{experience}[1][]{popbase,colframe=popblue,colback=popblueL,before upper={\popboxhead{Expérience}{popblue}{#1}}}
% « Le savais-tu ? » : carte violette pleine (contexte, culture, Cameroun)
\newtcolorbox{savaistu}[1][]{popbase,colback=poppurple,colframe=popink,
  fontupper=\popsemi\fontsize{10.4}{14.2}\selectfont\color{white},
  before upper={\poppill{Le savais-tu\,?}{white}{poppurple}\ifx\relax#1\relax\else\hspace{6pt}{\popxbold\color{white}#1}\fi\par\vspace{7pt}}}
% Exercice résolu : énoncé en haut, \tcblower puis la solution sur fond crème
\newcounter{popexo}\newcounter{popexr}
\newtcolorbox{exoresolu}[1][]{popbase,colframe=poporange,colbacklower=popcream,
  segmentation style={popink,line width=1.2pt,dashed},
  before upper={\stepcounter{popexr}\popboxhead{Exercice résolu \thepopexr}{poporange}{#1}},
  before lower={\poppill{Solution}{popink}{popcream}\par\vspace{6pt}}}
% Exercice (non résolu en place) — la correction est dans la partie Corrigés
\newtcolorbox{exercice}[1][]{popbase,colframe=popink,
  before upper={\stepcounter{popexo}\popboxhead{Exercice \thepopexo}{popink}{#1}}}
% Corrigé
\newtcolorbox{corrige}[1][]{popbase,colframe=popgreen,colback=popgreenL,
  before upper={\popboxhead{Corrigé}{popgreen}{#1}}}
% Objectifs / prérequis / bilan
\newtcolorbox{objectifs}{popbase,colframe=popink,colback=poporangeL,
  before upper={\popboxhead{Objectifs de la leçon}{popink}{}}}
\newtcolorbox{prerequis}{popbase,colframe=popink,colback=popblueL,
  before upper={\popboxhead{Prérequis}{popblue}{}}}
\newtcolorbox{bilan}{popbase,colframe=popink,colback=white,boxrule=2.8pt,
  before upper={\popboxhead{Fiche bilan}{poporange}{L'essentiel à connaître pour l'examen}}}
% Figure encadrée avec légende
\newtcolorbox{popfigure}[1][]{enhanced,breakable=false,colback=white,colframe=popline,boxrule=1.4pt,arc=8pt,
  left=10pt,right=10pt,top=10pt,bottom=8pt,before skip=10pt,after skip=12pt,halign=center,
  lower separated=false,halign lower=center,
  fontlower=\popsemi\fontsize{9}{11.5}\selectfont\color{popmuted},
  #1}
\newcommand{\legende}[1]{\par\vspace{6pt}{\popsemi\fontsize{9}{11.5}\selectfont\color{popmuted}#1}}

% ============================================================
% Signature OnBuch+
% ============================================================
\newcommand{\poplogoinline}[1]{{\poplogo\fontsize{#1}{#1}\selectfont\color{popink}OnBuch\color{poporange}+}}
\newcommand{\signatureonbuch}{%
  \begin{tcolorbox}[enhanced,colback=popink,colframe=popink,boxrule=2.2pt,arc=10pt,
      drop shadow=poporange,left=14pt,right=14pt,top=10pt,bottom=10pt,before skip=0pt,after skip=0pt]
    \tikz[baseline=-0.7ex]\node[circle,fill=popgreen,draw=popcream,line width=1.3pt,minimum size=22pt,inner sep=0]{\popcheck{white}};%
    \hspace{9pt}\begin{minipage}[c]{0.62\linewidth}
      {\popxbold\fontsize{12}{13}\selectfont\color{popcream}Signé\ \textbullet\ {\poplogo OnBuch\color{poporange}+}}\par\vspace{2pt}
      {\raggedright\popsemi\fontsize{8}{10.5}\selectfont\color{popline}\MakeUppercase{Équipe pédagogique OnBuch+}\\\MakeUppercase{Conforme au programme officiel MINESEC}\par}
    \end{minipage}\hfill
    {\poplogo\fontsize{22}{22}\selectfont\color{popcream}OnBuch\color{poporange}+}
  \end{tcolorbox}%
}

% ============================================================
% Page de garde
% #1 titre · #2 matière · #3 niveau · #4 module · #5 n° leçon · #6 durée ·
% #7 description / objectifs (texte ou itemize)
% ============================================================
\newcommand{\pagegarde}[7]{%
  \thispagestyle{empty}
  \newgeometry{a4paper,margin=1.7cm,top=1.6cm,bottom=1.4cm}%
  \begin{tikzpicture}[remember picture,overlay]
    \fill[poporange!18] ([xshift=-3.2cm,yshift=-3.6cm]current page.north east) circle (4.6cm);
    \fill[poppurple!14] ([xshift=2.4cm,yshift=7.6cm]current page.south west) circle (3.8cm);
  \end{tikzpicture}%
  {\poplogo\fontsize{30}{30}\selectfont\color{popink}OnBuch\color{poporange}+}\hfill
  \raisebox{0.6ex}{\poppill{Cours signé \textbullet\ Premium}{popink}{popcream}}\par
  \vspace{1pt}\popsquiggle{poppurple}\par
  \vspace{8pt}
  \begin{tcolorbox}[enhanced,colback=poporange,colframe=popink,boxrule=2.8pt,arc=13pt,
      drop shadow=popink,left=18pt,right=18pt,top=12pt,bottom=13pt,after skip=10pt]
    \poppill{#2 \textbullet\ #3}{popink}{popcream}\hspace{5pt}\poppill{#4}{white}{popink}\par\vspace{8pt}
    {\popdisplay\fontsize{14}{14}\selectfont\color{popink}LEÇON #5}\par\vspace{4pt}
    {\raggedright\hyphenpenalty=10000\exhyphenpenalty=10000\popdisplay\fontsize{25}{27}\selectfont\color{popcream}\MakeUppercase{#1}\par}
    \vspace{8pt}
    {\popxbold\fontsize{9.5}{9.5}\selectfont\color{popink}\MakeUppercase{Durée conseillée : #6}}
  \end{tcolorbox}
  \begin{tcolorbox}[enhanced,colback=white,colframe=popink,boxrule=2.2pt,arc=10pt,
      drop shadow=popink,left=16pt,right=16pt,top=10pt,bottom=10pt,after skip=8pt]
    \poppill{Dans cette leçon}{poppurple}{white}\par\vspace{5pt}
    \popsemi\fontsize{9.3}{12.6}\selectfont\color{popink}#7
  \end{tcolorbox}
  \noindent\begin{minipage}[t]{0.487\linewidth}
    \begin{tcolorbox}[enhanced,colback=popgreen,colframe=popink,boxrule=2.2pt,arc=10pt,
        drop shadow=popink,left=11pt,right=11pt,top=10pt,bottom=10pt,height=3.1cm,valign=center]
      \tcbox[colback=white,colframe=popink,boxrule=1.3pt,arc=5pt,boxsep=0pt,nobeforeafter,
        left=3pt,right=3pt,top=3pt,bottom=3pt,valign=center]{\qrcode[height=1.9cm]{https://play.google.com/store/apps/details?id=com.luvvix.onbuch&hl=fr}}%
      \hspace{8pt}\begin{minipage}[c]{0.5\linewidth}
        {\popdisplay\fontsize{11}{12.5}\selectfont\color{white}\MakeUppercase{Télécharge OnBuch}}\par\vspace{4pt}
        {\raggedright\popsemi\fontsize{8.4}{11}\selectfont\color{white}Gratuit \textbullet\ Google Play\\Tuteur IA, annales, concours, résultats\par}
      \end{minipage}
    \end{tcolorbox}
  \end{minipage}\hfill
  \begin{minipage}[t]{0.487\linewidth}
    \begin{tcolorbox}[enhanced,colback=poppurple,colframe=popink,boxrule=2.2pt,arc=10pt,
        drop shadow=popink,left=11pt,right=11pt,top=10pt,bottom=10pt,height=3.1cm,valign=center]
      \tikz[baseline=-0.7ex]\node[circle,fill=white,draw=popink,line width=1.3pt,minimum size=1.6cm,inner sep=0]{\popdisplay\fontsize{24}{24}\selectfont\color{poppurple}+};%
      \hspace{8pt}\begin{minipage}[c]{0.6\linewidth}
        {\popdisplay\fontsize{11}{12.5}\selectfont\color{white}\MakeUppercase{Passe à OnBuch+}}\par\vspace{4pt}
        {\raggedright\popsemi\fontsize{8.4}{11}\selectfont\color{white}Tous les cours signés, Tuteur IA illimité, fiches PDF — onglet OnBuch+ de l'app\par}
      \end{minipage}
    \end{tcolorbox}
  \end{minipage}
  \vspace{8pt}
  \popcontenu
  \vfill
  \signatureonbuch
  \restoregeometry
  \newpage
}

% Rangée « Ce cours contient » (page de garde)
\newcommand{\poptile}[3]{% #1 couleur #2 gros texte #3 légende
  \begin{tcolorbox}[enhanced,colback=#1,colframe=popink,boxrule=2pt,arc=9pt,shadow={2.5pt}{-2.5pt}{0pt}{popink},
      left=6pt,right=6pt,top=9pt,bottom=9pt,halign=center,valign=center,height=1.75cm,width=\linewidth,nobeforeafter]
    {\popdisplay\fontsize{12.5}{13}\selectfont\color{white}\MakeUppercase{#2}}\par\vspace{3pt}
    {\centering\popsemi\fontsize{7.8}{9.5}\selectfont\color{white}#3\par}
  \end{tcolorbox}}
\newcommand{\popcontenu}{%
  \par\noindent{\popxboldls\fontsize{7.5}{7.5}\selectfont\color{popmuted}CE COURS SIGNÉ CONTIENT}\par\vspace{7pt}
  \noindent\begin{minipage}[t]{0.235\linewidth}\poptile{popblue}{Cours}{complet, illustré, pas à pas}\end{minipage}\hfill
  \begin{minipage}[t]{0.235\linewidth}\poptile{poporange}{Méthodes}{et exercices résolus}\end{minipage}\hfill
  \begin{minipage}[t]{0.235\linewidth}\poptile{poppink}{Exercices}{type examen + corrigés}\end{minipage}\hfill
  \begin{minipage}[t]{0.235\linewidth}\poptile{popgreen}{Fiche bilan}{l'essentiel à retenir}\end{minipage}\par}

% Sommaire
\newtcolorbox{sommaire}{enhanced,colback=white,colframe=popink,boxrule=2.2pt,arc=10pt,
  drop shadow=popink,left=18pt,right=18pt,top=13pt,bottom=13pt,before skip=6pt,after skip=20pt,
  before upper={\poppill{Sommaire}{poppurple}{white}\par\vspace{9pt}\popsemi\fontsize{10.6}{14.5}\selectfont\color{popink}}}

% Signature de fin de cours — poussée tout en bas de la dernière page
\newcommand{\findecours}{%
  \par\vfill\nopagebreak
  \begin{center}\popsquiggle{poporange}\end{center}\vspace{4pt}
  \signatureonbuch
}

%% FIGFIX-BEGIN v5 — correctifs globaux des figures (docs/onbuch-generation/tools/figfix)
\usepackage{adjustbox}\usepackage{etoolbox}\tcbuselibrary{raster}
% toute figure TikZ/pgfplots plus large que la ligne est réduite à la largeur disponible
\BeforeBeginEnvironment{tikzpicture}{\begin{adjustbox}{max width=\linewidth}}
\AfterEndEnvironment{tikzpicture}{\end{adjustbox}}
% légendes pgfplots lisibles par-dessus les courbes
\pgfplotsset{every axis/.append style={legend style={fill=white,fill opacity=0.92,text opacity=1,draw=black!15,inner sep=3pt}}}
% étiquettes d'arêtes (syntaxe edge["..."]) avec fond blanc
\tikzset{every edge quotes/.append style={fill=white,inner sep=1.2pt}}
%% FIGFIX-END
\begin{document}

\pagegarde{La radioactivité et la loi de décroissance}{Physique}{Tle C}{Module 4 : Onde, matière et transformations nucléaires}{P19}{5 h}{Cette leçon étudie le noyau atomique, sa cohésion (énergie de liaison, courbe d'Aston) et ses transformations spontanées : la radioactivité α, β⁻, β⁺ et γ. Elle établit et exploite la loi de décroissance radioactive, outil fondamental de la datation et des applications médicales et industrielles.
\vspace{4pt}
\begin{multicols}{2}\small
\begin{itemize}
\item À la fin de cette leçon, je sais définir un nucléide, des isotopes et représenter un noyau par la notation ᴬzX
\item je sais calculer le défaut de masse, l'énergie de liaison et l'énergie de liaison par nucléon d'un noyau
\item je sais exploiter la courbe d'Aston pour prévoir la stabilité d'un noyau
\item je sais définir la radioactivité et citer ses propriétés (spontanée, aléatoire, inévitable)
\item je sais écrire et équilibrer une équation de désintégration α, β⁻, β⁺ ou γ à l'aide des lois de Soddy
\item je sais établir et exploiter la loi de décroissance N = N₀e\^{}(−λt) et calculer la période radioactive
\end{itemize}\end{multicols}}

\begin{sommaire}
\begin{multicols}{2}
\begin{enumerate}
\item Le noyau atomique et sa cohésion
\item La radioactivité : définition et propriétés
\item Types de désintégrations et lois de Soddy
\item Loi de décroissance radioactive
\item Activité d'un échantillon radioactif
\item Bilan énergétique et applications
\item Activité d'intégration
\item Méthodes et exercices résolus
\item Exercices
\item Corrigés des exercices
\item Fiche bilan
\end{enumerate}
\end{multicols}
\end{sommaire}

\begin{objectifs}
\begin{itemize}
\item À la fin de cette leçon, je sais définir un nucléide, des isotopes et représenter un noyau par la notation ᴬzX
\item je sais calculer le défaut de masse, l'énergie de liaison et l'énergie de liaison par nucléon d'un noyau
\item je sais exploiter la courbe d'Aston pour prévoir la stabilité d'un noyau
\item je sais définir la radioactivité et citer ses propriétés (spontanée, aléatoire, inévitable)
\item je sais écrire et équilibrer une équation de désintégration α, β⁻, β⁺ ou γ à l'aide des lois de Soddy
\item je sais établir et exploiter la loi de décroissance N = N₀e\^{}(−λt) et calculer la période radioactive
\item je sais calculer l'activité d'un échantillon radioactif à un instant donné
\item je sais effectuer le bilan énergétique d'une désintégration
\item je sais appliquer la loi de décroissance à la datation au carbone 14
\end{itemize}
\end{objectifs}

\begin{prerequis}
\begin{itemize}
\item Modèle de l'atome (noyau, électrons, numéro atomique) vu en Seconde et Première
\item Relation d'Einstein E = mc² et conversions d'unités d'énergie (J, MeV)
\item Fonction exponentielle et logarithme népérien (mathématiques de Terminale)
\item Équations différentielles du type y' = −λy (mathématiques de Terminale)
\item Conservation de la charge électrique et du nombre de nucléons dans une réaction
\end{itemize}
\end{prerequis}

\newpage

\coursec{Le noyau atomique et sa cohésion}

En 2023, une équipe d'archéologues découvre près de Foumban des fragments d'ossements et de charbon de bois dans un site ancien : comment déterminer leur âge sans aucun document écrit ? À l'hôpital de référence de Douala, on injecte à des patients des substances « radioactives » pour diagnostiquer des cancers : pourquoi ces substances disparaissent-elles d'elles-mêmes au bout de quelques heures ? Et pourquoi le minerai d'uranium reste-t-il dangereux pendant des millénaires alors qu'un médicament radioactif devient inoffensif en une journée ? Toutes ces questions trouvent leur réponse dans un phénomène naturel découvert par Henri Becquerel en 1896 : la \cle{radioactivité}, régie par une loi mathématique remarquable, la loi de décroissance exponentielle.

Mais avant de comprendre pourquoi certains noyaux se transforment spontanément, il faut d'abord comprendre de quoi un noyau est fait et, surtout, \emph{ce qui le tient ensemble}. C'est l'objet de cette première section : la constitution du noyau atomique et le secret de sa cohésion, qui nous mènera à l'une des plus belles courbes de la physique : la courbe d'Aston.

\courssub{Constitution du noyau : protons, neutrons et notation symbolique}

\textbf{De l'atome au noyau.} Tu sais depuis la classe de Seconde que la matière est constituée d'atomes. L'atome possède une structure lacunaire : presque toute sa masse est concentrée dans un noyau central minuscule, autour duquel gravitent les électrons. Les ordres de grandeur sont saisissants :

\begin{itemize}
\item diamètre de l'atome : environ $10^{-10}\ \text{m}$ (un dixième de nanomètre) ;
\item diamètre du noyau : environ $10^{-15}\ \text{m}$ à $10^{-14}\ \text{m}$ (quelques femtomètres).
\end{itemize}

Le noyau est donc environ \num{100000} fois plus petit que l'atome ! Si l'atome avait la taille du stade Ahmadou Ahidjo de Yaoundé, le noyau serait une bille posée au centre du terrain. Et pourtant, cette bille contiendrait plus de \SI{99,97}{\%} de la masse du stade tout entier.

\textbf{Les nucléons.} Le noyau est constitué de deux types de particules, appelées collectivement \cle{nucléons} :

\begin{itemize}
\item les \cle{protons}, de charge électrique $+e$ avec $e = \num{1,602e-19}\ \text{C}$, et de masse $m_p = \num{1,6726e-27}\ \text{kg}$ ;
\item les \cle{neutrons}, électriquement neutres, de masse $m_n = \num{1,6749e-27}\ \text{kg}$, légèrement supérieure à celle du proton.
\end{itemize}

On caractérise un noyau par deux nombres entiers :

\begin{definition}[Numéro atomique et nombre de masse]
Le \cle{numéro atomique} $Z$ est le nombre de protons du noyau ; il définit l'élément chimique. Le \cle{nombre de masse} $A$ est le nombre total de nucléons (protons + neutrons). Le nombre de neutrons est donc $N = A - Z$. Un noyau se note symboliquement :
\[ \ce{^{A}_{Z}X} \]
où $X$ est le symbole chimique de l'élément.
\end{definition}

\begin{exemplebox}[Lecture de la notation symbolique]
Le noyau d'uranium \ce{^{235}_{92}U} contient $Z = 92$ protons et $N = 235 - 92 = 143$ neutrons, soit $A = 235$ nucléons au total. Le noyau de fer \ce{^{56}_{26}Fe} contient 26 protons et 30 neutrons.
\end{exemplebox}

\begin{popfigure}
\begin{center}
\begin{tikzpicture}[scale=1.0]
% Noyau : amas de nucléons
\shade[ball color=popblue] (0.15,0.35) circle (0.34);
\shade[ball color=poporange] (-0.25,0.1) circle (0.34);
\shade[ball color=popblue] (0.3,-0.25) circle (0.34);
\shade[ball color=poporange] (-0.1,-0.45) circle (0.34);
\shade[ball color=popblue] (-0.45,-0.35) circle (0.34);
\shade[ball color=poporange] (0.55,0.05) circle (0.34);
\shade[ball color=popblue] (-0.6,0.35) circle (0.34);
\shade[ball color=poporange] (0.0,0.85) circle (0.34);
\draw[popdark, thick, dashed] (0,0.1) circle (1.15);
% Annotations
\draw[->, popblue, thick] (0.15,0.35) -- (2.6,1.4) node[right, popdark] {\textbf{proton} : charge $+e$, $m_p$};
\draw[->, poporange, thick] (-0.25,0.1) -- (-2.9,-1.3) node[left, popdark] {\textbf{neutron} : sans charge, $m_n$};
\draw[->, popdark] (1.15,0.1) -- (2.6,-0.7) node[right, popdark] {rayon $\approx 10^{-15}\ \text{m}$};
% Notation symbolique
\node[popdark] at (0,-2.3) {$\ce{^{A}_{Z}X}$ \quad : \quad $Z$ protons, $(A-Z)$ neutrons, $A$ nucléons};
\end{tikzpicture}
\end{center}
\legende{Le noyau $\ce{^{A}_{Z}X}$ : un amas compact de protons (bleu) et de neutrons (orange), cent mille fois plus petit que l'atome.}
\end{popfigure}

\courssub{Nucléides et isotopes}

\begin{definition}[Nucléide]
Un \cle{nucléide} est l'ensemble des noyaux (ou des atomes) possédant le même couple $(Z, A)$, c'est-à-dire le même nombre de protons et le même nombre de neutrons. Chaque nucléide est représenté par la notation $\ce{^{A}_{Z}X}$.
\end{definition}

Il existe environ 300 nucléides stables naturels et plusieurs milliers de nucléides instables (radioactifs), naturels ou artificiels.

\begin{definition}[Isotopes]
Des \cle{isotopes} sont des nucléides de \textbf{même numéro atomique $Z$} (même élément chimique) mais de \textbf{nombres de masse $A$ différents} : ils diffèrent donc par leur nombre de neutrons.
\end{definition}

\begin{exemplebox}[Des isotopes célèbres]
\begin{itemize}
\item Le carbone possède notamment deux isotopes : $\ce{^{12}_{6}C}$ (6 protons, 6 neutrons, stable, 98,9 \% du carbone naturel) et $\ce{^{14}_{6}C}$ (6 protons, 8 neutrons, radioactif) — c'est lui qui permettra de dater les ossements de Foumban !
\item L'uranium naturel est un mélange de $\ce{^{238}_{92}U}$ (99,3 \%) et de $\ce{^{235}_{92}U}$ (0,7 \%), tous deux radioactifs mais de comportements très différents dans un réacteur nucléaire.
\item L'hydrogène a trois isotopes : $\ce{^{1}_{1}H}$ (hydrogène), $\ce{^{2}_{1}H}$ (deutérium) et $\ce{^{3}_{1}H}$ (tritium, radioactif).
\end{itemize}
\end{exemplebox}

\begin{attention}[Isotopes = même chimie, mais pas même noyau]
Des isotopes ont les mêmes propriétés \emph{chimiques} (mêmes électrons, mêmes réactions chimiques) car la chimie ne dépend que de $Z$. En revanche, leurs propriétés \emph{nucléaires} (stabilité, radioactivité) peuvent être totalement différentes. Ne dis jamais que deux isotopes « ont le même nombre de nucléons » : c'est précisément le contraire !
\end{attention}

\courssub{L'unité de masse atomique}

Les masses des noyaux, de l'ordre de $10^{-27}\ \text{kg}$, sont minuscules. Pour manipuler des nombres plus parlants, les physiciens ont défini une unité adaptée.

\begin{definition}[Unité de masse atomique]
L'\cle{unité de masse atomique}, de symbole $\text{u}$, est le douzième de la masse d'un atome de carbone 12 :
\[ 1\ \text{u} = \frac{m(\ce{^{12}_{6}C})}{12} = \num{1,66054e-27}\ \text{kg}. \]
\end{definition}

Dans cette unité, les masses des nucléons sont très simples à retenir :
\[ m_p = \num{1,00728}\ \text{u}, \qquad m_n = \num{1,00866}\ \text{u}, \qquad m_e = \num{0,00055}\ \text{u}. \]

À cette unité de masse correspond une unité d'énergie très utilisée en physique nucléaire, le \cle{mégélectronvolt} : $1\ \text{MeV} = \num{1,602e-13}\ \text{J}$. La relation d'Einstein $E = mc^2$ donne l'équivalence fondamentale que tu utiliseras dans tous tes calculs :

\begin{aretenir}[Équivalence masse-énergie]
\[ 1\ \text{u} \cdot c^2 = \num{931,5}\ \text{MeV}. \]
Autrement dit, une masse de 1 u « vaut » 931,5 MeV d'énergie. Cette conversion sera ta clé pour tous les bilans énergétiques nucléaires.
\end{aretenir}

\courssub{Le défaut de masse : l'énigme de la masse manquante}

\textbf{Une observation troublante.} Prenons le noyau d'hélium $\ce{^{4}_{2}He}$, constitué de 2 protons et 2 neutrons. Calculons la masse de ses constituants séparés, au repos :
\[ 2m_p + 2m_n = 2 \times \num{1,00728} + 2 \times \num{1,00866} = \num{4,03188}\ \text{u}. \]
Or la masse mesurée du noyau d'hélium est $m(\ce{^{4}_{2}He}) = \num{4,00150}\ \text{u}$. Le noyau est \emph{moins lourd} que la somme de ses morceaux ! Il manque environ $\num{0,03038}\ \text{u}$, soit près de \SI{0,8}{\%} de la masse. Ce phénomène est général : \textbf{la masse d'un noyau est toujours inférieure à la somme des masses de ses nucléons séparés au repos}.

\begin{definition}[Défaut de masse]
Le \cle{défaut de masse} d'un noyau $\ce{^{A}_{Z}X}$ est la différence, toujours positive, entre la masse des nucléons séparés au repos et la masse du noyau :
\[ \Delta m = Z\, m_p + (A - Z)\, m_n - m(\ce{^{A}_{Z}X}) \quad > 0. \]
\end{definition}

Où est passée cette masse ? Einstein répond en 1905 avec sa célèbre relation $E = mc^2$ : la masse est une forme d'énergie. Lorsque les nucléons s'assemblent pour former le noyau, une partie de leur masse est convertie en énergie, qui est \emph{libérée} vers l'extérieur (sous forme de rayonnement). Le noyau formé est donc dans un état d'énergie plus basse : il est stable, et il faudra lui \emph{rendre} cette énergie pour le casser.

\courssub{Énergie de liaison et cohésion du noyau}

\textbf{Le problème de la cohésion.} Réfléchis un instant : le noyau contient des protons, tous chargés positivement, confinés dans un volume de $10^{-15}\ \text{m}$ de rayon. La répulsion électrostatique entre eux est colossale — la force de Coulomb entre deux protons distants de $10^{-15}\ \text{m}$ vaut environ $F = \dfrac{e^2}{4\pi\varepsilon_0 r^2} \approx 230\ \text{N}$, une force énorme à l'échelle d'une particule ! Pourtant le noyau n'explose pas. Il doit donc exister une interaction attractive, inconnue de la physique classique, qui domine la répulsion électrique à très courte distance.

\begin{propriete}[Interaction forte et énergie de liaison (admise)]
La cohésion du noyau provient de l'\cle{interaction forte}, une force attractive intense entre nucléons (protons et neutrons indifféremment), de très courte portée (environ $10^{-15}\ \text{m}$), qui l'emporte sur la répulsion électrostatique entre protons. On admet que l'\cle{énergie de liaison}, définie par la relation d'Einstein
\[ E_l = \Delta m \cdot c^2, \]
mesure cette cohésion : c'est l'énergie qu'il faut \textbf{fournir} au noyau au repos pour le dissocier en ses nucléons séparés, immobiles et sans interaction.
\end{propriete}

\begin{attention}[Le piège du signe : énergie à fournir ou énergie libérée ?]
C'est l'erreur la plus fréquente au Bac. Retiens bien :
\begin{itemize}
\item Pour \textbf{casser} le noyau en nucléons séparés, il faut \textbf{fournir} $E_l = \Delta m \cdot c^2$ : l'énergie de liaison est une énergie à apporter, elle est comptée positivement.
\item Lors de la \textbf{formation} du noyau à partir des nucléons, cette même énergie est \textbf{libérée} : le bilan du système est alors $-E_l$.
\end{itemize}
Même valeur absolue, signes opposés selon le sens de la transformation. Une énergie de liaison élevée signifie un noyau \emph{très lié}, donc \emph{très stable} — et non « plein d'énergie prête à exploser » !
\end{attention}

\courssub{Énergie de liaison par nucléon : le critère de stabilité}

Comparer les énergies de liaison totales de deux noyaux n'a guère de sens : un noyau d'uranium (238 nucléons) est évidemment plus lié en valeur absolue qu'un noyau d'hélium (4 nucléons). Pour comparer les stabilités, on raisonne \emph{par nucléon}.

\begin{definition}[Énergie de liaison par nucléon]
L'\cle{énergie de liaison par nucléon} d'un noyau est le quotient
\[ \frac{E_l}{A} = \frac{\Delta m \cdot c^2}{A}. \]
Elle s'exprime en MeV/nucléon. \textbf{Plus $E_l/A$ est grande, plus le noyau est stable.}
\end{definition}

\begin{methode}[Calculer $E_l/A$ d'un noyau à partir des masses]
\begin{enumerate}
\item Identifier $Z$ et $A$ dans la notation $\ce{^{A}_{Z}X}$ ; en déduire le nombre de neutrons $A - Z$.
\item Calculer la masse des constituants séparés : $Z\, m_p + (A-Z)\, m_n$, en unités u.
\item Calculer le défaut de masse : $\Delta m = Z m_p + (A-Z) m_n - m(\text{noyau})$.
\item Convertir en énergie : $E_l = \Delta m \times \num{931,5}\ \text{MeV/u}$.
\item Diviser par $A$ pour obtenir $E_l/A$ en MeV/nucléon, puis conclure sur la stabilité (valeur élevée $\Rightarrow$ noyau stable ; le maximum est $\approx 8,8\ \text{MeV/nucléon}$).
\end{enumerate}
\end{methode}

\begin{exoresolu}[Énergie de liaison par nucléon du fer 56]
On donne la masse du noyau de fer : $m(\ce{^{56}_{26}Fe}) = \num{55,9207}\ \text{u}$ ; $m_p = \num{1,00728}\ \text{u}$ ; $m_n = \num{1,00866}\ \text{u}$ ; $1\ \text{u}\cdot c^2 = \num{931,5}\ \text{MeV}$. Calculer l'énergie de liaison par nucléon du fer 56 et commenter.
\tcblower
\textbf{Solution détaillée.}

Le noyau $\ce{^{56}_{26}Fe}$ contient $Z = 26$ protons et $A - Z = 56 - 26 = 30$ neutrons.

\emph{Étape 1 : masse des nucléons séparés.}
\begin{align*}
Z m_p + (A-Z) m_n &= 26 \times \num{1,00728} + 30 \times \num{1,00866} \\
&= \num{26,18928} + \num{30,25980} = \num{56,44908}\ \text{u}.
\end{align*}

\emph{Étape 2 : défaut de masse.}
\[ \Delta m = \num{56,44908} - \num{55,9207} = \num{0,52838}\ \text{u}. \]

\emph{Étape 3 : énergie de liaison.}
\[ E_l = \Delta m \cdot c^2 = \num{0,52838} \times \num{931,5} \approx 492,2\ \text{MeV}. \]

\emph{Étape 4 : énergie de liaison par nucléon.}
\[ \frac{E_l}{A} = \frac{492,2}{56} \approx 8,79\ \text{MeV/nucléon} \approx 8,8\ \text{MeV/nucléon}. \]

\emph{Conclusion.} Cette valeur est parmi les plus élevées de tous les noyaux : le fer 56 est l'un des noyaux les plus stables de l'Univers. Pour arracher un seul nucléon au noyau de fer, il faudrait fournir environ 8,8 MeV, soit $\num{1,4e-12}\ \text{J}$ — une énergie gigantesque à l'échelle atomique (environ dix millions de fois l'énergie d'une liaison chimique).
\end{exoresolu}

\courssub{La courbe d'Aston : la carte de la stabilité nucléaire}

En 1920, le physicien britannique Francis Aston mesura précisément les masses de nombreux noyaux grâce à son spectrographe de masse. En portant l'opposé de l'énergie de liaison par nucléon, $-E_l/A$, en fonction du nombre de masse $A$, on obtient la \cle{courbe d'Aston}. La convention de signe (porter $-E_l/A$) fait que \textbf{les noyaux les plus stables sont au fond de la courbe}.

\begin{popfigure}
\begin{center}
\begin{tikzpicture}
\begin{axis}[
width=0.92\linewidth, height=7.2cm,
xlabel={$A$ (nombre de masse)}, ylabel={$-E_l/A$ (MeV/nucléon)},
xmin=0, xmax=250, ymin=-9.5, ymax=0.5,
grid=both, grid style={popline!40},
axis lines=left, thick,
legend style={draw=none}
]
\addplot[popcurve, domain=2:250, samples=250]
{-8.8 + 8.8*exp(-0.06*x) + 0.012*x - 1.2*exp(-0.5*(x-4))};
% Points remarquables
\addplot[only marks, mark=*, poporange, mark size=2pt] coordinates {(4,-7.07) (12,-7.68) (56,-8.79) (235,-7.6)};
\node[popdark, anchor=south west, font=\small] at (axis cs:60,-8.9) {$\ce{^{56}Fe}$ : $-8{,}8$};
\node[popdark, anchor=south, font=\small] at (axis cs:4,-6.6) {$\ce{^{4}He}$};
\node[popdark, anchor=south, font=\small] at (axis cs:12,-7.2) {$\ce{^{12}C}$};
\node[popdark, anchor=north, font=\small] at (axis cs:235,-7.5) {$\ce{^{235}U}$};
% Zones
\draw[<-, popgreen, very thick] (axis cs:20,-4.5) -- (axis cs:45,-4.5) node[midway, above, popgreen!60!popdark, font=\small\bfseries] {fusion};
\draw[<-, poppink, very thick] (axis cs:200,-5.2) -- (axis cs:160,-5.2) node[midway, above, poppink!60!popdark, font=\small\bfseries] {fission};
\draw[popdark, dashed] (axis cs:56,-9.5) -- (axis cs:56,-8.79);
\end{axis}
\end{tikzpicture}
\end{center}
\legende{Courbe d'Aston : $-E_l/A$ en fonction de $A$. Le minimum (fond de la « cuvette ») se situe autour du fer ($A \approx 56$) : ce sont les noyaux les plus stables.}
\end{popfigure}

\textbf{Lecture qualitative de la courbe.} Trois zones se distinguent :

\begin{itemize}
\item \textbf{Petits $A$ (noyaux légers, $A < 20$)} : la courbe descend fortement. Deux noyaux légers qui s'unissent (\cle{fusion nucléaire}) forment un noyau plus profond dans la cuvette, donc plus stable : la fusion \emph{libère} de l'énergie. C'est la source d'énergie du Soleil et des étoiles.
\item \textbf{Grands $A$ (noyaux lourds, $A > 190$)} : la courbe remonte doucement. Un noyau lourd qui se scinde en deux noyaux moyens (\cle{fission nucléaire}) produit des noyaux plus stables : la fission aussi libère de l'énergie. C'est le principe des réacteurs nucléaires.
\item \textbf{Autour de $A \approx 56$ (fer, nickel)} : le fond de la cuvette, $E_l/A \approx 8,8\ \text{MeV/nucléon}$. Ces noyaux ne peuvent ni fusionner ni fissionner en libérant de l'énergie : ce sont les « cendres » ultimes des transformations nucléaires.
\end{itemize}

\begin{aretenir}[Ce qu'il faut retenir de la courbe d'Aston]
La courbe d'Aston représente $-E_l/A$ en fonction de $A$. Son minimum, autour du fer 56 ($E_l/A \approx 8,8\ \text{MeV/nucléon}$), correspond aux noyaux les plus stables. Les noyaux légers peuvent libérer de l'énergie par \textbf{fusion}, les noyaux lourds par \textbf{fission}. Les noyaux situés loin du fond de la cuvette sont instables : beaucoup sont radioactifs — ce sera l'objet de la suite de la leçon.
\end{aretenir}

\begin{savaistu}[Pourquoi les étoiles finissent en fer]
Dans le cœur des étoiles, la fusion transforme l'hydrogène en hélium, puis l'hélium en carbone, oxygène, silicium\ldots{} Chaque étape libère de l'énergie car le noyau produit est plus profond dans la cuvette d'Aston. Mais la chaîne s'arrête au fer 56 : fusionner du fer \emph{coûterait} de l'énergie au lieu d'en libérer. Le cœur de l'étoile, devenu une boule de fer, ne produit plus d'énergie : il s'effondre, puis explose en supernova, dispersant dans l'Univers les éléments lourds dont tu es fait. Le fer de ton organisme et l'or des bijoux sont littéralement de la poussière d'étoiles !
\end{savaistu}

\begin{exercice}[Pour t'entraîner : le carbone 12]
On donne $m(\ce{^{12}_{6}C}) = \num{12,00000}\ \text{u}$ (par définition), $m_p = \num{1,00728}\ \text{u}$, $m_n = \num{1,00866}\ \text{u}$.
\begin{enumerate}
\item Calculer le défaut de masse du noyau de carbone 12.
\item En déduire son énergie de liaison $E_l$ en MeV.
\item Calculer $E_l/A$ et comparer à celle du fer 56. Lequel des deux noyaux est le plus stable ?
\end{enumerate}
\end{exercice}

\begin{corrige}[Exercice : le carbone 12]
\begin{enumerate}
\item Le noyau $\ce{^{12}_{6}C}$ contient 6 protons et 6 neutrons :
\[ \Delta m = 6 \times \num{1,00728} + 6 \times \num{1,00866} - \num{12,00000} = \num{12,09564} - \num{12,00000} = \num{0,09564}\ \text{u}. \]
\item $E_l = \num{0,09564} \times \num{931,5} \approx 89,1\ \text{MeV}$.
\item $\dfrac{E_l}{A} = \dfrac{89,1}{12} \approx 7,4\ \text{MeV/nucléon}$. C'est moins que les $8,8\ \text{MeV/nucléon}$ du fer 56 : le carbone 12 est stable, mais \emph{moins lié} que le fer. Sur la courbe d'Aston, il se situe au-dessus du fond de la cuvette : cohérent avec le fait que les étoiles gagnent de l'énergie en fusionnant le carbone vers des noyaux plus lourds.
\end{enumerate}
\end{corrige}

Nous savons désormais ce qui tient le noyau ensemble et comment mesurer sa stabilité. Mais que se passe-t-il lorsqu'un noyau est \emph{trop loin} du fond de la cuvette d'Aston — trop lourd, ou avec un mauvais équilibre entre protons et neutrons ? Il se transforme spontanément : c'est la radioactivité, que nous étudions dans la section suivante.

\coursec{La radioactivité : définition et propriétés}

Dans la section précédente, nous avons découvert que le noyau atomique est un assemblage de nucléons maintenus par l'interaction forte, et que sa cohésion se mesure par l'énergie de liaison par nucléon $\dfrac{E_{\ell}}{A}$. La courbe d'Aston nous a appris un fait capital : tous les noyaux ne sont pas également stables. Certains, trop lourds ou mal équilibrés en protons et neutrons, se situent loin du creux de la courbe. Que deviennent ces noyaux instables ? La nature leur offre une issue : ils se transforment d'eux-mêmes en noyaux plus stables, en éjectant des particules et de l'énergie. Ce phénomène naturel, découvert à la toute fin du XIX\textsuperscript{e} siècle, s'appelle la \cle{radioactivité}. C'est lui que nous allons définir et caractériser dans cette section.

\courssub{Une découverte qui a bouleversé la physique}

\textbf{Le hasard heureux de Becquerel (1896)}\par
En janvier 1896, Wilhelm Röntgen découvre les rayons X, émis par des tubes à décharge. Le physicien français Henri Becquerel se demande alors si les corps \emph{fluorescents} — qui émettent de la lumière après avoir été éclairés — émettraient aussi des rayons X. Il choisit un sel d'uranium (le sulfate double d'uranyle et de potassium), le pose sur une plaque photographique enveloppée de papier noir, et l'expose au soleil. La plaque est impressionnée : le sel semble bien émettre un rayonnement.

Mais les jours suivants, le temps est couvert à Paris. Becquerel range le sel et la plaque dans un tiroir, en attendant le soleil. Quelques jours plus tard, par scrupule, il développe quand même la plaque : elle est \emph{fortement} impressionnée, alors que le sel n'a reçu \textbf{aucune lumière}. Conclusion révolutionnaire : le sel d'uranium émet un rayonnement \textbf{de lui-même}, en permanence, sans aucune excitation extérieure. Becquerel vient de découvrir la radioactivité (qu'il appelle alors « rayonnement uranique »).

\textbf{Pierre et Marie Curie : la chasse aux éléments nouveaux}\par
Marie Curie, jeune physicienne polonaise installée à Paris, choisit ce phénomène comme sujet de thèse. Avec son mari Pierre Curie, elle mesure systématiquement le rayonnement de minerais d'uranium (la pechblende) et constate qu'il est bien plus intense que celui de l'uranium pur : le minerai doit contenir des éléments inconnus, beaucoup plus actifs. Après des années de traitement manuel de tonnes de minerai dans un hangar délabré de la rue Lhomond, le couple isole en 1898 deux nouveaux éléments : le \cle{polonium} (hommage à la Pologne natale de Marie) et le \cle{radium}, dont le rayonnement est environ un million de fois plus intense que celui de l'uranium à masse égale. C'est Marie Curie qui forge le mot \emph{radioactivité}.

\textbf{Rutherford : l'identification des rayonnements}\par
Dès 1899, Ernest Rutherford soumet le rayonnement à un champ magnétique et distingue deux composantes : des rayons $\alpha$, déviés faiblement (particules lourdes et chargées positivement), et des rayons $\beta$, fortement déviés (particules légères et chargées négativement). Paul Villard identifie en 1900 une troisième composante, insensible au champ : les rayons $\gamma$, de nature électromagnétique. Nous étudierons ces trois types d'émission en détail dans la section suivante.

\begin{popfigure}
\begin{tikzpicture}[scale=1, transform shape]
\draw[line width=1.2pt, popblue, -{Stealth}] (0,0) -- (14.6,0);
\foreach \x/\annee in {0.6/1896, 3.4/1898, 6.2/1899, 9.0/1900, 11.8/1903, 14.2/1911}{
  \draw[popblue, line width=1pt] (\x,0.12) -- (\x,-0.12);
  \node[below, popdark, font=\small] at (\x,-0.12) {\annee};
}
\node[above, align=center, text width=2.5cm, font=\footnotesize, popdark] at (0.6,0.25) {\textbf{Becquerel}\\ rayonnement de l'uranium};
\node[above, align=center, text width=2.5cm, font=\footnotesize, popdark] at (3.4,0.25) {\textbf{P. et M. Curie}\\ polonium et radium};
\node[above, align=center, text width=2.3cm, font=\footnotesize, popdark] at (6.2,0.25) {\textbf{Rutherford}\\ rayons $\alpha$ et $\beta$};
\node[above, align=center, text width=2.3cm, font=\footnotesize, popdark] at (9.0,0.25) {\textbf{Villard}\\ rayons $\gamma$};
\node[above, align=center, text width=2.5cm, font=\footnotesize, popdark] at (11.8,0.25) {\textbf{Nobel de physique}\\ Becquerel et les Curie};
\node[above, align=center, text width=2.3cm, font=\footnotesize, popdark] at (14.2,0.25) {\textbf{Nobel de chimie}\\ Marie Curie};
\end{tikzpicture}
\legende{Frise chronologique des grandes découvertes sur la radioactivité (1896--1911).}
\end{popfigure}

\begin{savaistu}[Marie Curie, une scientifique hors norme]
Née Maria Skłodowska à Varsovie en 1867, Marie Curie reste à ce jour la \textbf{seule personne} à avoir reçu deux prix Nobel dans deux sciences différentes : physique en 1903 (avec Pierre Curie et Henri Becquerel, pour les recherches sur la radioactivité) et chimie en 1911 (pour l'isolement du radium et du polonium). Elle est aussi la première femme professeure à la Sorbonne. Sa fille Irène Joliot-Curie recevra à son tour le prix Nobel de chimie en 1935. Pendant la Première Guerre mondiale, Marie Curie organise des unités mobiles de radiographie — les « petites Curie » — pour soigner les blessés au front.
\end{savaistu}

\courssub{Définition de la radioactivité}

Revenons au sens physique du phénomène. Le noyau d'uranium contenu dans le sel de Becquerel est instable : son énergie de liaison par nucléon est trop faible. Pour gagner en stabilité, il se transforme en un noyau différent, plus lié, en éjectant une particule et de l'énergie. C'est cette transformation spontanée que l'on appelle radioactivité.

\begin{definition}[Radioactivité et radionucléide]
La \cle{radioactivité} est la transformation \textbf{spontanée} d'un noyau atomique \textbf{instable} en un noyau \textbf{plus stable}, accompagnée de l'émission d'un \textbf{rayonnement} (particules $\alpha$, $\beta^-$ ou $\beta^+$, et le plus souvent rayonnement électromagnétique $\gamma$).

Un noyau qui subit cette transformation est appelé \cle{noyau radioactif} ou \cle{radionucléide}. Le noyau initial est le noyau \textbf{père}, le noyau obtenu est le noyau \textbf{fils}.
\end{definition}

Trois points de cette définition méritent d'être soulignés :

\begin{itemize}
\item La transformation concerne le \textbf{noyau}, pas le nuage électronique : l'élément chimique change (le numéro atomique $Z$ peut varier), contrairement à une réaction chimique où les noyaux restent intacts.
\item Elle s'accompagne d'une \textbf{libération d'énergie} : le noyau fils est plus lié que le noyau père, donc la masse totale diminue et l'énergie correspondante est emportée par les particules émises (nous ferons ce bilan en section 6).
\item Le noyau fils est \emph{souvent} formé dans un état excité ; il revient à son état fondamental en émettant un photon $\gamma$, très énergétique.
\end{itemize}

\begin{exemplebox}[Quelques radionucléides]
\begin{itemize}
\item L'uranium 238 : \ce{^{238}_{92}U}, radioactif $\alpha$, présent dans les minerais.
\item Le carbone 14 : \ce{^{14}_{6}C}, radioactif $\beta^-$, utilisé pour la datation des vestiges archéologiques.
\item Le cobalt 60 : \ce{^{60}_{27}Co}, radioactif $\beta^-$ et émetteur $\gamma$, utilisé en radiothérapie dans les services d'oncologie des grands hôpitaux.
\item L'iode 131 : \ce{^{131}_{53}I}, radioactif $\beta^-$, utilisé en médecine nucléaire pour explorer la thyroïde.
\end{itemize}
\end{exemplebox}

\courssub{Les quatre propriétés fondamentales de la radioactivité}

Toute la physique de la décroissance radioactive repose sur quatre propriétés expérimentales qu'il faut connaître et savoir commenter. Elles justifieront, dans la section 4, la forme mathématique de la loi de décroissance $N = N_0 e^{-\lambda t}$.

\textbf{Propriété 1 : la spontanéité}\par
La désintégration d'un noyau radioactif ne dépend d'\textbf{aucune condition extérieure} : ni la température, ni la pression, ni le champ électrique ou magnétique, ni l'état chimique de l'atome (solide, liquide, gaz, corps pur ou engagé dans une molécule) ne modifient le comportement du noyau. Cela s'explique : la radioactivité met en jeu l'interaction forte et l'interaction faible \emph{à l'intérieur du noyau}, à des échelles d'énergie de l'ordre du MeV. Or l'agitation thermique, même à \SI{3000}{K}, correspond à une énergie moyenne $k_B T$ de l'ordre de \SI{0,26}{eV} par particule : c'est plusieurs millions de fois trop faible pour influencer le noyau. Chauffer, broyer, comprimer ou dissoudre un échantillon radioactif ne change donc \emph{rien} à sa décroissance.

\textbf{Propriété 2 : le caractère aléatoire}\par
Il est \textbf{impossible de prévoir l'instant} où un noyau donné va se désintégrer. Deux noyaux identiques, placés dans des conditions identiques, ne se désintègrent pas au même moment : l'un peut « vivre » une fraction de seconde, l'autre des millénaires. La désintégration d'un noyau individuel est un phénomène quantique, fondamentalement probabiliste. En revanche, pour une population de noyaux \emph{très grande} (un gramme de matière contient de l'ordre de $10^{22}$ noyaux), les lois des grands nombres s'appliquent : le comportement \emph{moyen} de l'échantillon devient parfaitement prévisible et obéit à la loi de décroissance exponentielle. La radioactivité est donc aléatoire au niveau d'un noyau, mais déterministe au niveau d'un échantillon macroscopique.

\textbf{Propriété 3 : l'inévitabilité}\par
Un noyau radioactif finit \textbf{toujours} par se désintégrer : c'est une question de temps, pas de chance définitive. Aucun noyau instable n'échappe à sa transformation. L'échelle de temps caractéristique varie énormément selon les radionucléides : de quelques microsecondes (polonium 214, dont la demi-vie vaut environ \SI{164}{\micro s}) à plusieurs milliards d'années (uranium 238, dont la demi-vie vaut environ \SI{4,5e9}{ans}, comparable à l'âge de la Terre).

\textbf{Propriété 4 : l'indépendance vis-à-vis de l'âge}\par
Un noyau radioactif \textbf{ne vieillit pas} : sa probabilité de se désintégrer dans la prochaine seconde est la même, qu'il vienne d'être formé ou qu'il existe depuis mille ans. On dit que la désintégration est un phénomène \emph{sans mémoire}. C'est cette propriété qui donnera à la loi de décroissance sa forme exponentielle : à chaque instant, la proportion de noyaux qui se désintègrent par unité de temps est constante.

\begin{propriete}[Propriétés de la radioactivité — à retenir]
La radioactivité est un phénomène :
\begin{enumerate}
\item \textbf{spontané} : il ne dépend d'aucune condition extérieure (température, pression, état chimique, champ) ;
\item \textbf{aléatoire} : on ne peut pas prévoir l'instant de désintégration d'un noyau donné ;
\item \textbf{inévitable} : tout noyau radioactif finit par se désintégrer ;
\item \textbf{indépendant de l'âge} : un noyau ne vieillit pas, sa probabilité de désintégration par unité de temps est constante.
\end{enumerate}
Conséquence : pour un échantillon contenant un grand nombre de noyaux, le nombre moyen de désintégrations par unité de temps est proportionnel au nombre $N$ de noyaux présents — ce qui conduira à la loi $N(t) = N_0 e^{-\lambda t}$.
\end{propriete}

\begin{attention}[Erreur fréquente : croire qu'on peut accélérer la décroissance]
Beaucoup d'élèves (et de candidats au Bac !) imaginent qu'en chauffant fortement, en broyant finement ou en dissolvant un échantillon radioactif, on pourrait « épuiser » plus vite sa radioactivité. C'est \textbf{faux}. Ces opérations modifient l'agitation des atomes et les liaisons chimiques, c'est-à-dire le cortège électronique, mais pas le noyau. La demi-vie d'un radionucléide est une constante de la nature : on ne connaît \textbf{aucun procédé} pour la modifier de façon notable. C'est d'ailleurs tout le problème du traitement des déchets nucléaires : il faut simplement attendre, parfois des milliers d'années, qu'ils se désintègrent d'eux-mêmes.
\end{attention}

\courssub{Rayonnements émis et pénétration dans la matière}

Lors d'une désintégration, le noyau émet des \textbf{particules} (les particules $\alpha$ sont des noyaux d'hélium \ce{^{4}_{2}He}, les particules $\beta^-$ sont des électrons, les particules $\beta^+$ sont des positrons) et/ou un \textbf{rayonnement électromagnétique} $\gamma$ (photons très énergétiques). Il faut bien distinguer les deux vocabulaires : une particule possède une masse, un rayonnement $\gamma$ est de la lumière de très haute énergie, sans masse.

Ces émissions n'ont pas du tout la même capacité de pénétration dans la matière :

\begin{itemize}
\item les particules $\alpha$, lourdes et doublement chargées, interagissent fortement avec la matière : elles sont arrêtées par une simple \textbf{feuille de papier} ou par quelques centimètres d'air ;
\item les particules $\beta$, beaucoup plus légères, traversent le papier mais sont arrêtées par quelques millimètres d'\textbf{aluminium} ;
\item les rayonnements $\gamma$, sans charge ni masse, sont très pénétrants : il faut une épaisseur importante de \textbf{plomb} ou de béton pour les atténuer fortement.
\end{itemize}

\begin{popfigure}
\begin{tikzpicture}[scale=1, transform shape]
% Source
\fill[poporange] (0,0) circle (0.28);
\node[below, font=\small, popdark] at (0,-0.3) {source radioactive};
% alpha
\draw[line width=2.2pt, popgreen, -{Stealth}] (0.3,0.9) -- (2.6,0.9);
\node[above, popgreen, font=\small] at (1.4,0.95) {$\alpha$};
% beta
\draw[line width=1.6pt, popblue, -{Stealth}] (0.3,0) -- (5.6,0);
\node[above, popblue, font=\small] at (2.2,0.05) {$\beta$};
% gamma
\draw[line width=1.6pt, poppurple, decorate, decoration={snake, amplitude=0.8mm, segment length=3mm}, -{Stealth}] (0.3,-0.9) -- (9.6,-0.9);
\node[below, poppurple, font=\small] at (2.2,-0.95) {$\gamma$};
% papier
\fill[popgoldL] (2.8,-1.7) rectangle (3.2,1.7);
\node[above, font=\small, popdark] at (3.0,1.75) {papier};
% aluminium
\fill[popblueL] (5.8,-1.7) rectangle (6.6,1.7);
\node[above, font=\small, popdark] at (6.2,1.75) {aluminium};
\node[below, font=\footnotesize, popdark] at (6.2,-1.75) {quelques mm};
% plomb
\fill[popmuted!60] (9.8,-1.7) rectangle (11.4,1.7);
\node[above, font=\small, popdark] at (10.6,1.75) {plomb};
\node[below, font=\footnotesize, popdark] at (10.6,-1.75) {quelques cm};
% gamma attenuated after lead
\draw[line width=0.9pt, poppurple!50, decorate, decoration={snake, amplitude=0.5mm, segment length=3mm}, -{Stealth}] (11.4,-0.9) -- (13.2,-0.9);
\node[right, font=\footnotesize, popmuted] at (13.2,-0.9) {atténué};
\end{tikzpicture}
\legende{Pouvoir de pénétration comparé des rayonnements $\alpha$, $\beta$ et $\gamma$ : les $\alpha$ sont arrêtés par le papier, les $\beta$ par quelques millimètres d'aluminium, les $\gamma$ seulement atténués par une forte épaisseur de plomb.}
\end{popfigure}

\textbf{Effets biologiques (complément utile au Bac)}\par
Ces rayonnements sont dits \emph{ionisants} : en traversant la matière vivante, ils arrachent des électrons aux molécules, ce qui peut endommager l'ADN des cellules. C'est pourquoi on se protège des sources radioactives par la \textbf{distance}, la \textbf{durée d'exposition} réduite et les \textbf{écrans} adaptés (tabliers de plomb en radiologie, murs de béton épais autour des sources de cobalt 60 en radiothérapie). Retour de médaille : à dose contrôlée, ces mêmes rayonnements permettent de détruire des cellules cancéreuses (radiothérapie) ou d'imager des organes (scintigraphie) — nous y reviendrons dans la section 6.

\courssub{Application : raisonner sur le caractère aléatoire}

\begin{exoresolu}[Désintégrations observées en une minute]
Un compteur Geiger-Müller placé près d'un échantillon de césium 137 enregistre, lors de quatre mesures successives d'une minute chacune, les nombres de désintégrations suivants : $n_1 = 487$, $n_2 = 512$, $n_3 = 469$, $n_4 = 501$.
\begin{enumerate}
\item Ces résultats sont-ils compatibles avec les propriétés de la radioactivité ? Justifier.
\item Un élève affirme : « Pour obtenir toujours la même valeur, il suffirait de maintenir l'échantillon à température constante. » Commenter.
\item Calculer le nombre moyen de désintégrations par minute, puis par seconde.
\end{enumerate}
\tcblower
\textbf{Solution détaillée}\par
1. Les quatre mesures sont \emph{voisines mais différentes}. C'est exactement ce qu'on attend d'un phénomène \textbf{aléatoire} : l'instant de désintégration de chaque noyau est imprévisible, donc le nombre de désintégrations pendant une durée fixée fluctue autour d'une valeur moyenne. Ces résultats illustrent le caractère aléatoire de la radioactivité ; ils sont pleinement compatibles avec ses propriétés.

2. L'affirmation de l'élève est \textbf{fausse}. La radioactivité est \textbf{spontanée} : elle ne dépend d'aucune condition extérieure, et en particulier pas de la température. Même à température rigoureusement constante, les fluctuations statistiques subsisteraient, car elles sont d'origine quantique et non thermique. La seule façon de réduire l'écart relatif des fluctuations est d'augmenter la durée de comptage ou le nombre de noyaux observés (loi des grands nombres).

3. Nombre moyen par minute :
\[
\bar{n} = \dfrac{n_1 + n_2 + n_3 + n_4}{4} = \dfrac{487 + 512 + 469 + 501}{4} = \dfrac{1969}{4} \approx 492 \text{ désintégrations par minute.}
\]
Par seconde :
\[
\bar{n}' = \dfrac{492}{60} \approx \SI{8,2}{désintégrations.s^{-1}}.
\]
Ce nombre moyen de désintégrations par seconde n'est autre que l'\emph{activité} de l'échantillon, grandeur que nous définirons rigoureusement en section 5 : ici $A \approx \SI{8,2}{Bq}$ pour la zone comptée.
\end{exoresolu}

\begin{aretenir}[L'essentiel de la section]
\begin{itemize}
\item La radioactivité, découverte par Becquerel en 1896 puis explorée par Pierre et Marie Curie, est la transformation \textbf{spontanée} d'un noyau instable en un noyau plus stable, avec émission de rayonnement ($\alpha$, $\beta$ ou $\gamma$).
\item Elle est \textbf{spontanée} (indépendante de toute condition extérieure), \textbf{aléatoire} (instant de désintégration imprévisible pour un noyau), \textbf{inévitable} (tout noyau radioactif finit par se désintégrer) et \textbf{indépendante de l'âge} du noyau.
\item Aléatoire pour un noyau, elle devient prévisible en moyenne pour un grand nombre de noyaux : c'est le fondement de la loi de décroissance exponentielle que nous établirons en section 4.
\item Les rayonnements $\alpha$ sont peu pénétrants (papier), les $\beta$ moyennement (aluminium), les $\gamma$ très pénétrants (plomb épais) ; tous sont ionisants, d'où des règles de protection strictes.
\end{itemize}
\end{aretenir}

Maintenant que la radioactivité est définie et que ses propriétés sont posées, une question naturelle se pose : comment écrire précisément la transformation d'un noyau père en noyau fils, et quelles grandeurs se conservent lors de cette métamorphose ? C'est l'objet de la section suivante, consacrée aux types de désintégrations et aux lois de conservation de Soddy.

\coursec{Types de désintégrations et lois de Soddy}

Dans la section précédente, nous avons découvert la \cle{radioactivité} : un phénomène \emph{spontané}, \emph{aléatoire} et \emph{inévitable}, par lequel un noyau instable se transforme en un autre noyau en émettant un rayonnement. Mais une question demeure : \textbf{que se passe-t-il exactement à l'intérieur du noyau ?} Quelles particules sont éjectées, et le noyau fils est-il prévisible ? C'est toute la beauté du travail de Frederick Soddy (prix Nobel de chimie 1921) : quelques règles de conservation simples permettent d'écrire \emph{toutes} les équations de désintégration, comme on équilibre une équation chimique. C'est ce que nous allons apprendre à faire dans cette section.

\courssub{Les lois de conservation de Soddy}

\textbf{Observation fondatrice}\par

Lorsqu'un noyau d'uranium se désintègre, on ne trouve jamais n'importe quoi dans les produits : la nature respecte des règles strictes. En compilant des centaines de désintégrations, Soddy a dégagé deux lois de conservation, qui sont la traduction nucléaire de principes fondamentaux de la physique : la conservation de la charge électrique et la conservation du nombre de nucléons.

\begin{propriete}[Lois de conservation de Soddy]
Dans toute réaction nucléaire (désintégration ou réaction provoquée), il y a :
\begin{itemize}
\item \textbf{conservation du nombre de nucléons} $A$ : la somme des nombres de masse des réactifs est égale à la somme des nombres de masse des produits ;
\item \textbf{conservation du nombre de charge} $Z$ : la somme des numéros atomiques des réactifs est égale à la somme des numéros atomiques des produits.
\end{itemize}
Pour une désintégration $\ce{^{A}_{Z}X -> ^{A'}_{Z'}Y + ^{a}_{z}p}$ (où $p$ est la particule émise), cela s'écrit :
\[ A = A' + a \qquad \text{et} \qquad Z = Z' + z. \]
\end{propriete}

\begin{experience}[Démonstration guidée : équilibrer une équation de désintégration]
Cherchons le noyau fils obtenu lorsque le thorium \ce{^{232}_{90}Th} émet une particule $\alpha$, c'est-à-dire un noyau d'hélium \ce{^{4}_{2}He}.
\begin{enumerate}
\item Écrivons l'équation avec les inconnues $A'$ et $Z'$ :
\[ \ce{^{232}_{90}Th -> ^{A'}_{Z'}Y + ^{4}_{2}He} \]
\item Appliquons la conservation de $A$ : $232 = A' + 4$, donc $A' = 232 - 4 = 228$.
\item Appliquons la conservation de $Z$ : $90 = Z' + 2$, donc $Z' = 88$.
\item Identifions l'élément : $Z' = 88$ correspond au \textbf{radium} (Ra) dans le tableau périodique.
\end{enumerate}
Conclusion : $\ce{^{232}_{90}Th -> ^{228}_{88}Ra + ^{4}_{2}He}$. Le noyau fils est parfaitement déterminé : il n'y a aucune ambiguïté.
\end{experience}

\begin{attention}[La notation fait tout le travail]
Le symbole chimique (X, Y, Ra\ldots) est \emph{redondant} avec le numéro atomique $Z$ : écrire \ce{^{228}_{88}Ra} suffit, puisque $Z = 88$ désigne forcément le radium. Au Baccalauréat, on ne te demandera jamais de mémoriser le tableau périodique : l'élément fils sera donné, ou seul le couple $(A', Z')$ sera demandé. En revanche, \textbf{oublier de vérifier les deux conservations} est l'erreur la plus fréquente : vérifie toujours $A$ \emph{et} $Z$ séparément.
\end{attention}

\courssub{La radioactivité $\alpha$}

\textbf{Intuition et définition}\par

Certains noyaux sont très lourds : ils contiennent plus de 200 nucléons. Les forces répulsives entre les nombreux protons deviennent si grandes que le noyau gagne en stabilité en \emph{expulsant un paquet compact} de 2 protons et 2 neutrons — précisément un noyau d'hélium \ce{^{4}_{2}He}, particulièrement stable (on le justifie avec la courbe d'Aston). C'est la \cle{radioactivité alpha}.

\begin{definition}[Radioactivité $\alpha$]
La radioactivité $\alpha$ est l'émission spontanée d'un noyau d'hélium \ce{^{4}_{2}He} (appelé \textbf{particule} $\alpha$) par un noyau lourd instable. L'équation générale s'écrit :
\[ \ce{^{A}_{Z}X -> ^{A-4}_{Z-2}Y + ^{4}_{2}He} \qquad (\text{en général pour } A > 200) \]
Le noyau fils a \textbf{4 nucléons de moins} et \textbf{2 charges de moins} que le noyau père.
\end{definition}

Vérifions les lois de Soddy sur l'équation générale : $A = (A-4) + 4$ (vérifié) et $Z = (Z-2) + 2$ (vérifié). Les particules $\alpha$, chargées positivement ($q = +2e$) et massives, sont peu pénétrantes : une feuille de papier ou quelques centimètres d'air suffisent à les arrêter. Elles sont en revanche dangereuses en cas d'ingestion ou d'inhalation (radon dans certaines habitations).

\begin{exemplebox}[Désintégration $\alpha$ du radium 226 (découverte historique)]
Le radium 226, découvert par Marie et Pierre Curie, est $\alpha$-radioactif :
\[ \ce{^{226}_{88}Ra -> ^{222}_{86}Rn + ^{4}_{2}He} \]
Vérification : $226 = 222 + 4$ et $88 = 86 + 2$. Le noyau fils est le \textbf{radon} 222, un gaz lui-même radioactif. Chaque noyau de radium désintégré libère une énergie de l'ordre de $\SI{4,9}{MeV}$, emportée essentiellement sous forme d'énergie cinétique de la particule $\alpha$.
\end{exemplebox}

\courssub{La radioactivité $\beta^-$}

\textbf{Un électron qui sort du noyau ?!}\par

Voici un paradoxe apparent : certains noyaux émettent des \emph{électrons}. Or il n'y a pas d'électrons dans le noyau ! La résolution de ce paradoxe est l'une des plus belles idées de la physique nucléaire : l'électron est \emph{créé au moment de la désintégration}, lorsqu'un \textbf{neutron se transforme en proton} à l'intérieur du noyau :
\[ \ce{^{1}_{0}n -> ^{1}_{1}p + ^{0}_{-1}e} \]
(En réalité, il est aussi émis une particule quasi insaisissable, l'antineutrino $\bar{\nu}$, non exigible au programme de Terminale mais indispensable à la conservation de l'énergie.)

\begin{definition}[Radioactivité $\beta^-$]
La radioactivité $\beta^-$ est l'émission spontanée d'un électron \ce{^{0}_{-1}e} par un noyau instable \textbf{riche en neutrons}. Un neutron du noyau se transforme en proton. L'équation générale s'écrit :
\[ \ce{^{A}_{Z}X -> ^{A}_{Z+1}Y + ^{0}_{-1}e} \]
Le nombre de masse $A$ est \textbf{inchangé} ; le numéro atomique \textbf{augmente de 1}.
\end{definition}

Vérification de Soddy : $A = A + 0$ (vérifié) et $Z = (Z+1) + (-1)$ (vérifié). Les électrons $\beta^-$ sont plus pénétrants que les particules $\alpha$ : il faut quelques millimètres d'aluminium pour les arrêter.

\begin{exemplebox}[Désintégration $\beta^-$ du carbone 14]
Le carbone 14, célèbre pour la datation des vestiges archéologiques (nous l'exploiterons dans la section 6), est $\beta^-$-radioactif :
\[ \ce{^{14}_{6}C -> ^{14}_{7}N + ^{0}_{-1}e} \]
Vérification : $14 = 14 + 0$ et $6 = 7 - 1$. Le noyau fils est l'\textbf{azote} 14, stable. Un neutron de trop dans le carbone 14 s'est transformé en proton.
\end{exemplebox}

\begin{attention}[L'électron $\beta^-$ ne vient PAS du cortège électronique]
Erreur classique de copie : « l'électron émis est un électron de l'atome qui passe près du noyau ». \textbf{Faux !} L'électron $\beta^-$ est \emph{créé dans le noyau} par la transformation d'un neutron en proton. Conséquence spectaculaire : le noyau fils possède un proton de plus, donc appartient à un \emph{élément chimique différent} (l'azote dans l'exemple du carbone 14), alors que le cortège électronique n'a pas encore changé. Retiens : la radioactivité est un phénomène \textbf{nucléaire}, indépendant du cortège électronique — c'est d'ailleurs pourquoi aucune condition extérieure (température, pression) ne peut la modifier.
\end{attention}

\courssub{La radioactivité $\beta^+$}

\textbf{Le miroir de $\beta^-$}\par

Certains noyaux artificiels (produits en médecine nucléaire, par exemple) souffrent du défaut inverse : ils ont \emph{trop de protons}. Le noyau se stabilise alors en transformant un \textbf{proton en neutron}, avec émission d'un \cle{positon} (ou positron) \ce{^{0}_{+1}e}, l'antiparticule de l'électron : même masse, charge opposée $+e$.
\[ \ce{^{1}_{1}p -> ^{1}_{0}n + ^{0}_{+1}e} \]

\begin{definition}[Radioactivité $\beta^+$]
La radioactivité $\beta^+$ est l'émission spontanée d'un positon \ce{^{0}_{+1}e} par un noyau instable \textbf{riche en protons}. L'équation générale s'écrit :
\[ \ce{^{A}_{Z}X -> ^{A}_{Z-1}Y + ^{0}_{+1}e} \]
Le nombre de masse $A$ est \textbf{inchangé} ; le numéro atomique \textbf{diminue de 1}.
\end{definition}

Vérification : $A = A + 0$ (vérifié) et $Z = (Z-1) + 1$ (vérifié). Le positon émis ne vit pas longtemps : dès qu'il rencontre un électron de la matière environnante, les deux s'\emph{annihilent} en produisant deux photons $\gamma$ d'énergie $\SI{0,511}{MeV}$ chacun. C'est le principe de la \textbf{tomographie par émission de positons (TEP)}, utilisée en imagerie médicale dans les grands hôpitaux pour détecter les cancers.

\begin{exemplebox}[Désintégration $\beta^+$ du fluor 18 (imagerie TEP)]
\[ \ce{^{18}_{9}F -> ^{18}_{8}O + ^{0}_{+1}e} \]
Vérification : $18 = 18 + 0$ et $9 = 8 + 1$. Le fluor 18, injecté au patient lié à du glucose, se concentre dans les cellules cancéreuses très actives ; les photons d'annihilation détectés permettent de localiser la tumeur.
\end{exemplebox}

\courssub{La désexcitation $\gamma$}

\textbf{Un noyau fils « stressé »}\par

Après une désintégration $\alpha$ ou $\beta$, le noyau fils est souvent produit dans un \textbf{état excité}, noté $Y^*$ : il possède un excès d'énergie, exactement comme un atome dont un électron aurait sauté sur une orbite supérieure. Pour revenir à son état fondamental, il émet un \textbf{photon $\gamma$}, rayonnement électromagnétique de très grande énergie (de l'ordre du MeV, contre quelques eV pour la lumière visible).

\begin{definition}[Désexcitation $\gamma$]
La désexcitation $\gamma$ est l'émission d'un photon $\gamma$ par un noyau fils excité $Y^*$ qui revient à son état fondamental :
\[ \ce{^{A}_{Z}Y^* -> ^{A}_{Z}Y + \gamma} \]
Ni $A$ ni $Z$ ne changent : ce n'est \textbf{pas une transmutation}, mais une simple perte d'énergie. L'énergie du photon vaut $E_\gamma = h\nu = E_i - E_f$, différence des énergies des deux niveaux nucléaires.
\end{definition}

\begin{attention}[Le $\gamma$ accompagne, il ne remplace pas]
L'émission $\gamma$ ne modifie ni $A$ ni $Z$ : ne l'utilise \textbf{jamais} pour équilibrer une équation de désintégration ! Elle \emph{accompagne} très souvent une désintégration $\alpha$ ou $\beta$. Exemple complet avec le cobalt 60, utilisé en radiothérapie (notamment dans les centres de lutte contre le cancer de Yaoundé et Douala) :
\[ \ce{^{60}_{27}Co -> ^{60}_{28}Ni^* + ^{0}_{-1}e} \qquad \text{puis} \qquad \ce{^{60}_{28}Ni^* -> ^{60}_{28}Ni + \gamma} \]
Les photons $\gamma$ très pénétrants (il faut du plomb ou du béton épais pour s'en protéger) détruisent les cellules tumorales.
\end{attention}

\begin{popfigure}
\begin{center}
\begin{tikzpicture}[scale=1.0, font=\small]
% Axes
\draw[->, thick, popdark] (0,0) -- (0,4.6) node[above] {$E$};
% Niveau noyau père
\draw[very thick, popblue] (0.6,3.6) -- (3.2,3.6) node[right, popblue] {père $\ce{^{60}_{27}Co}$};
% Niveau excité fils
\draw[very thick, poporange] (4.4,2.6) -- (7.0,2.6) node[right, poporange] {fils excité $\ce{^{60}_{28}Ni^*}$};
% Niveau fondamental fils
\draw[very thick, popgreen!70!black] (4.4,0.5) -- (7.0,0.5) node[right, popgreen!70!black] {fils $\ce{^{60}_{28}Ni}$ (fondamental)};
% Flèche beta-
\draw[->, very thick, poppurple, decorate, decoration={snake, amplitude=0.6mm}] (3.3,3.55) .. controls (3.9,3.4) and (3.9,2.9) .. (4.35,2.65);
\node[poppurple] at (3.2,2.6) {$\beta^-$};
% Flèche gamma
\draw[->, very thick, poppink, decorate, decoration={snake, amplitude=0.8mm}] (5.7,2.55) -- (5.7,0.55);
\node[poppink, right] at (5.75,1.55) {$\gamma$ : $E_\gamma = h\nu$};
% Légende énergies
\node[popmuted, align=left] at (1.9,1.2) {Le noyau fils, créé\\ dans un état excité,\\ redescend en émettant\\ un photon $\gamma$.};
\end{tikzpicture}
\end{center}
\legende{Diagramme des niveaux d'énergie : la désintégration $\beta^-$ du cobalt 60 laisse le noyau de nickel dans un état excité, qui se désexcite en émettant un photon $\gamma$.}
\end{popfigure}

\courssub{Vue d'ensemble : le diagramme $(Z, N)$}

Une image vaut mille équations. Représentons chaque noyau par un point de coordonnées $(Z, N)$, où $N = A - Z$ est le nombre de neutrons. Chaque type de désintégration correspond à un \emph{déplacement caractéristique} dans ce diagramme :
\begin{itemize}
\item $\alpha$ : $Z$ diminue de 2 et $N$ diminue de 2 (diagonale vers le bas-gauche) ;
\item $\beta^-$ : $Z$ augmente de 1 et $N$ diminue de 1 (un neutron devient proton) ;
\item $\beta^+$ : $Z$ diminue de 1 et $N$ augmente de 1 (un proton devient neutron) ;
\item $\gamma$ : aucun déplacement.
\end{itemize}

\begin{popfigure}
\begin{center}
\begin{tikzpicture}[scale=0.95, font=\small]
% Grille et axes
\draw[->, thick, popdark] (0,0) -- (8.2,0) node[below right] {$Z$ (protons)};
\draw[->, thick, popdark] (0,0) -- (0,6.2) node[above left] {$N$ (neutrons)};
\foreach \x in {1,...,7} \draw[popline] (\x,0.05) -- (\x,-0.05) node[below] {\x};
\foreach \y in {1,...,5} \draw[popline] (0.05,\y) -- (-0.05,\y) node[left] {\y};
% Noyau initial
\fill[popblue] (5,4) circle (3pt);
\node[popblue, above right] at (5,4) {père $\ce{^{A}_{Z}X}$};
% Alpha
\fill[poporange] (3,2) circle (3pt);
\draw[->, very thick, poporange] (5,4) -- (3.1,2.1);
\node[poporange, below] at (3.4,2.9) {$\alpha$ : $(-2,-2)$};
% Beta moins
\fill[poppurple] (6,3) circle (3pt);
\draw[->, very thick, poppurple] (5.05,3.95) -- (5.95,3.1);
\node[poppurple, right] at (6.05,3.4) {$\beta^-$ : $(+1,-1)$};
% Beta plus
\fill[popgreen!70!black] (4,5) circle (3pt);
\draw[->, very thick, popgreen!70!black] (4.95,4.05) -- (4.1,4.9);
\node[popgreen!70!black, above left] at (3.7,4.6) {$\beta^+$ : $(-1,+1)$};
% Gamma
\draw[->, very thick, poppink, decorate, decoration={snake, amplitude=0.5mm}] (5.15,4.15) .. controls (5.9,4.6) and (5.5,5.1) .. (5.05,4.25);
\node[poppink] at (6.4,4.9) {$\gamma$ : $(0,0)$};
\end{tikzpicture}
\end{center}
\legende{Déplacements du noyau dans le diagramme $(Z, N)$ selon le type de désintégration. La désexcitation $\gamma$ laisse le noyau sur place.}
\end{popfigure}

\begin{methode}[Identifier un type de radioactivité ou compléter une équation]
\begin{enumerate}
\item \textbf{Regarder la particule émise} (si elle est donnée) : \ce{^{4}_{2}He} $\Rightarrow$ $\alpha$ ; \ce{^{0}_{-1}e} $\Rightarrow$ $\beta^-$ ; \ce{^{0}_{+1}e} $\Rightarrow$ $\beta^+$ ; photon $\gamma$ avec $A$ et $Z$ inchangés $\Rightarrow$ désexcitation.
\item \textbf{Ou comparer père et fils} : $A$ diminue de 4 et $Z$ de 2 $\Rightarrow$ $\alpha$ ; $A$ constant et $Z$ augmente de 1 $\Rightarrow$ $\beta^-$ ; $A$ constant et $Z$ diminue de 1 $\Rightarrow$ $\beta^+$.
\item \textbf{Pour compléter une équation incomplète}, écrire les deux équations de Soddy (conservation de $A$ et de $Z$) et résoudre : deux équations, deux inconnues au maximum.
\item \textbf{Vérifier} systématiquement les deux conservations sur l'équation finale, et identifier l'élément fils grâce à son $Z$.
\end{enumerate}
\end{methode}

\begin{exoresolu}[Identifier et compléter]
\textbf{Énoncé.} (1) Le polonium 218 se désintègre selon : $\ce{^{218}_{84}Po -> ^{214}_{82}Pb + X}$. Déterminer la nature de $X$ et le type de radioactivité. (2) Le phosphore 30 est radioactif $\beta^+$. Écrire l'équation de sa désintégration.
\tcblower
\textbf{Solution.} (1) Écrivons $X = \ce{^{a}_{z}X}$ et appliquons les lois de Soddy :
\begin{align*}
\text{Conservation de } A : \quad & 218 = 214 + a \Rightarrow a = 4 \\
\text{Conservation de } Z : \quad & 84 = 82 + z \Rightarrow z = 2
\end{align*}
Donc $X = \ce{^{4}_{2}He}$ : c'est une particule $\alpha$, et le polonium 218 est \textbf{$\alpha$-radioactif} :
\[ \ce{^{218}_{84}Po -> ^{214}_{82}Pb + ^{4}_{2}He} \]
(2) Pour une émission $\beta^+$, l'équation générale est $\ce{^{A}_{Z}X -> ^{A}_{Z-1}Y + ^{0}_{+1}e}$. Avec $A = 30$ et $Z = 15$ (phosphore) :
\[ \ce{^{30}_{15}P -> ^{30}_{14}Si + ^{0}_{+1}e} \]
Vérification : $30 = 30 + 0$ (vérifié) et $15 = 14 + 1$ (vérifié). Le noyau fils est le silicium 30 ($Z = 14$), stable.
\end{exoresolu}

\courssub{Les familles radioactives}

Une désintégration n'est souvent que la \emph{première étape} d'une longue cascade : le noyau fils est lui-même instable et se désintègre à son tour, et ainsi de suite, jusqu'à l'obtention d'un noyau \textbf{stable}. L'ensemble de ces désintégrations en chaîne constitue une \cle{famille radioactive}.

\begin{definition}[Famille radioactive]
Une famille radioactive est une suite de désintégrations successives ($\alpha$ et $\beta^-$ essentiellement) partant d'un noyau père très lourd et s'achevant sur un noyau stable, généralement un isotope du \textbf{plomb}.
\end{definition}

L'exemple le plus célèbre est la \textbf{famille de l'uranium 238}, qui compte 14 désintégrations successives (8 $\alpha$ et 6 $\beta^-$) et aboutit au plomb 206, stable :
\[ \ce{^{238}_{92}U ->[$\alpha$] ^{234}_{90}Th ->[$\beta^-$] ^{234}_{91}Pa ->[$\beta^-$] ^{234}_{92}U ->[$\alpha$] $\cdots$ ->[$\alpha$] ^{206}_{82}Pb} \; \text{(stable)} \]
Vérifions le bilan global : chaque $\alpha$ retire 4 nucléons, chaque $\beta^-$ n'en retire aucun. Avec 8 désintégrations $\alpha$ : $238 - 8 \times 4 = 238 - 32 = 206$ (vérifié). Pour la charge : $92 - 8 \times 2 + 6 \times 1 = 92 - 16 + 6 = 82$ (vérifié). Les lois de Soddy s'appliquent à \emph{chaque étape}, donc aussi au bilan global.

\begin{savaistu}[Le radon, héritier invisible de l'uranium]
Dans la famille de l'uranium 238 apparaît le radon 222, un gaz radioactif $\alpha$ issu de la désintégration du radium 226. Présent naturellement dans les sols granitiques et volcaniques, il peut s'accumuler dans les habitations mal ventilées et constitue, dans le monde, la deuxième cause de cancer du poumon après le tabac. Les familles radioactives expliquent aussi pourquoi les minerais d'uranium contiennent toujours du radium et du plomb : ce sont les « descendants » de l'uranium, accumulés depuis la formation de la Terre !
\end{savaistu}

\begin{aretenir}[L'essentiel de la section]
\begin{itemize}
\item \textbf{Lois de Soddy} : dans toute désintégration, le nombre de nucléons $A$ et le nombre de charge $Z$ se conservent.
\item \textbf{$\alpha$} : $\ce{^{A}_{Z}X -> ^{A-4}_{Z-2}Y + ^{4}_{2}He}$ — noyaux lourds ($A > 200$).
\item \textbf{$\beta^-$} : $\ce{^{A}_{Z}X -> ^{A}_{Z+1}Y + ^{0}_{-1}e}$ — un neutron devient un proton (noyaux riches en neutrons).
\item \textbf{$\beta^+$} : $\ce{^{A}_{Z}X -> ^{A}_{Z-1}Y + ^{0}_{+1}e}$ — un proton devient un neutron (noyaux riches en protons).
\item \textbf{$\gamma$} : $\ce{^{A}_{Z}Y^* -> ^{A}_{Z}Y + \gamma}$ — désexcitation sans changement de $A$ ni de $Z$, accompagnant souvent $\alpha$ ou $\beta$.
\item Les désintégrations s'enchaînent en \textbf{familles radioactives} jusqu'à un noyau stable (ex. : $\ce{^{238}_{92}U}$ $\rightarrow$ $\ce{^{206}_{82}Pb}$).
\end{itemize}
\end{aretenir}

Maintenant que nous savons \emph{quoi} émet un noyau instable et \emph{en quoi} il se transforme, il reste à répondre à une question temporelle : \emph{quand} ces désintégrations se produisent-elles, et combien de noyaux reste-t-il au bout d'un temps donné ? C'est l'objet de la section suivante : la loi de décroissance radioactive.

\coursec{Loi de décroissance radioactive}

Dans la section précédente, nous avons appris à écrire les équations de désintégration $\alpha$, $\beta^-$, $\beta^+$ et $\gamma$ grâce aux lois de conservation de Soddy. Mais ces équations décrivent \emph{un seul} événement : la transformation d'\emph{un} noyau. Or, dans la pratique — au laboratoire, à l'hôpital de Yaoundé, dans une roche du Mont Cameroun — on manipule des échantillons contenant des milliards de milliards de noyaux. Une question naturelle surgit alors : \textbf{comment évolue, au cours du temps, le nombre de noyaux radioactifs d'un échantillon ?} Peut-on prévoir combien de noyaux resteront dans un an, dans mille ans ? C'est toute l'ambition de la \cle{loi de décroissance radioactive}, l'un des plus beaux exemples de rencontre entre la physique et les mathématiques de Terminale.

\courssub{Approche statistique : l'équation différentielle de la décroissance}

\textbf{Une intuition pour commencer.} Imagine une immense salle contenant un million de pièces de monnaie posées sur pile. À chaque minute, chaque pièce a une petite probabilité de « basculer » sur face, sans qu'on puisse prévoir laquelle. Après une minute, le nombre de pièces ayant basculé sera d'autant plus grand qu'il y avait de pièces au départ, et d'autant plus grand que la durée d'observation est longue. La radioactivité fonctionne exactement ainsi : chaque noyau est un \emph{individu} dont le destin est imprévisible (propriété \cle{aléatoire}), mais la \emph{population} de noyaux obéit à une loi parfaitement régulière.

\begin{experience}[Du hasard individuel à la loi collective]
Considérons un échantillon contenant, à l'instant $t$, un nombre $N(t)$ de noyaux radioactifs d'un même nucléide ($N$ est très grand, typiquement $10^{20}$). Observons l'échantillon pendant une durée très courte $\mathrm{d}t$.
\begin{itemize}
\item La désintégration est \textbf{spontanée} et \textbf{aléatoire} : aucun facteur extérieur (température, pression, champ électrique) ne peut la déclencher ni l'empêcher, et on ne peut pas prévoir quel noyau va se désintégrer ni à quel instant.
\item Chaque noyau possède donc, par unité de temps, une même probabilité de se désintégrer, notée $\lambda$ et appelée \cle{constante radioactive}.
\end{itemize}
Le nombre de noyaux qui se désintègrent pendant $\mathrm{d}t$ est alors proportionnel au nombre $N$ de noyaux présents et à la durée $\mathrm{d}t$ :
\[ \text{nombre de désintégrations} = \lambda\, N\, \mathrm{d}t. \]
Comme chaque désintégration fait \emph{diminuer} $N$, la variation de $N$ s'écrit :
\[ \mathrm{d}N = -\lambda N\, \mathrm{d}t. \]
\end{experience}

\begin{definition}[Constante radioactive]
La \cle{constante radioactive} $\lambda$ d'un nucléide est la probabilité de désintégration d'un noyau par unité de temps. Elle s'exprime en $\mathrm{s^{-1}}$ (ou en $\mathrm{an^{-1}}$, $\mathrm{j^{-1}}$...). Elle est \textbf{caractéristique du nucléide} : elle ne dépend ni de la quantité de matière, ni des conditions extérieures, ni de l'âge de l'échantillon.
\end{definition}

Vérifions l'homogénéité : dans $\mathrm{d}N = -\lambda N\,\mathrm{d}t$, le membre de gauche est un nombre sans dimension, $N$ aussi, donc $[\lambda]\times[\mathrm{d}t] = 1$, c'est-à-dire $[\lambda] = \mathrm{T^{-1}}$. La constante $\lambda$ s'exprime bien en $\mathrm{s^{-1}}$.

En divisant $\mathrm{d}N = -\lambda N\,\mathrm{d}t$ par $\mathrm{d}t$, on obtient l'\textbf{équation différentielle de la décroissance radioactive} :
\[ \dfrac{\mathrm{d}N}{\mathrm{d}t} = -\lambda N. \]

\courssub{Résolution de l'équation différentielle : la loi $N(t) = N_0\,e^{-\lambda t}$}

Tu as rencontré en mathématiques l'équation différentielle $y' = ay$, dont les solutions sont les fonctions exponentielles $y = C\,e^{at}$. La physique te donne ici l'occasion de voir ce théorème à l'œuvre.

\begin{propriete}[Loi de décroissance radioactive — démonstration exigible]
La solution de l'équation différentielle $\dfrac{\mathrm{d}N}{\mathrm{d}t} = -\lambda N$ vérifiant la condition initiale $N(0) = N_0$ est :
\[ \boxed{N(t) = N_0\,e^{-\lambda t}} \]
où $N_0$ est le nombre de noyaux présents à l'instant $t = 0$ choisi comme origine des dates.
\end{propriete}

\begin{experience}[Démonstration guidée pas à pas]
\begin{enumerate}
\item \textbf{Forme générale des solutions.} D'après le cours de mathématiques, les solutions de l'équation différentielle $y' = ay$ (ici $a = -\lambda$) sont les fonctions de la forme $N(t) = C\,e^{-\lambda t}$, où $C$ est une constante réelle.
\item \textbf{Vérification.} Dérivons : $\dfrac{\mathrm{d}N}{\mathrm{d}t} = C\times(-\lambda)e^{-\lambda t} = -\lambda\,C\,e^{-\lambda t} = -\lambda N(t)$. L'équation est bien vérifiée.
\item \textbf{Prise en compte de la condition initiale.} À $t = 0$ : $N(0) = C\,e^{0} = C$. Or $N(0) = N_0$, donc $C = N_0$.
\item \textbf{Conclusion.} L'unique solution est $N(t) = N_0\,e^{-\lambda t}$. \hfill$\blacksquare$
\end{enumerate}
\end{experience}

\begin{aretenir}[La loi de décroissance]
Le nombre de noyaux radioactifs encore présents à la date $t$ dans un échantillon est :
\[ N(t) = N_0\,e^{-\lambda t} \]
avec $N_0$ le nombre initial de noyaux et $\lambda$ la constante radioactive du nucléide ($\mathrm{s^{-1}}$). Le nombre de noyaux décroît \textbf{exponentiellement} : il diminue de plus en plus lentement, sans jamais s'annuler rigoureusement.
\end{aretenir}

\courssub{Courbe de décroissance et interprétation graphique}

La fonction $t \mapsto N_0 e^{-\lambda t}$ est une exponentielle décroissante : sa courbe part de $N_0$ à $t = 0$, décroît rapidement au début (car il y a beaucoup de noyaux, donc beaucoup de désintégrations par seconde), puis de plus en plus lentement. La tangente à l'origine a pour équation $N = N_0(1 - \lambda t)$ : elle coupe l'axe des temps en $t = \dfrac{1}{\lambda} = \tau$, appelée \cle{temps caractéristique} (qui est aussi la \emph{vie moyenne} d'un noyau). À la date $t = \tau$, il reste $N = N_0 e^{-1} \approx 0{,}37\,N_0$.

\begin{popfigure}
\begin{center}
\begin{tikzpicture}
\begin{axis}[
  width=0.95\linewidth, height=7cm,
  axis lines=left,
  xlabel={$t$ (en unit\'es de $t_{1/2}$)}, ylabel={$N(t)/N_0$},
  xmin=0, xmax=3.4, ymin=0, ymax=1.08,
  xtick={0,1,2,3}, xticklabels={$0$,$t_{1/2}$,$2t_{1/2}$,$3t_{1/2}$},
  ytick={0.125,0.25,0.5,1}, yticklabels={$N_0/8$,$N_0/4$,$N_0/2$,$N_0$},
  grid=major, grid style={popline!40},
  tick label style={font=\small}, label style={font=\small}
]
\addplot[popcurve,domain=0:3.4,samples=200]{exp(-0.6931*x)};
\addplot[dashed,poporange,thick] coordinates {(1,0)(1,0.5)(0,0.5)};
\addplot[dashed,popgreen,thick] coordinates {(2,0)(2,0.25)(0,0.25)};
\addplot[dashed,poppurple,thick] coordinates {(3,0)(3,0.125)(0,0.125)};
\addplot[only marks,mark=*,mark size=2pt,popdark] coordinates {(1,0.5)(2,0.25)(3,0.125)};
\node[poporange,font=\small,anchor=west] at (axis cs:1.05,0.56) {$N_0/2$};
\node[popgreen,font=\small,anchor=west] at (axis cs:2.05,0.31) {$N_0/4$};
\node[poppurple,font=\small,anchor=west] at (axis cs:3.05,0.185) {$N_0/8$};
\end{axis}
\end{tikzpicture}
\end{center}
\legende{Courbe de décroissance radioactive $N(t)/N_0 = e^{-\lambda t}$. À chaque période $t_{1/2}$ supplémentaire, le nombre de noyaux restants est divisé par deux.}
\end{popfigure}

\courssub{Période radioactive (demi-vie)}

\begin{definition}[Période radioactive]
La \cle{période radioactive} (ou \cle{demi-vie}), notée $t_{1/2}$, d'un nucléide radioactif est la durée au bout de laquelle la \textbf{moitié} des noyaux initialement présents s'est désintégrée :
\[ N(t_{1/2}) = \dfrac{N_0}{2}. \]
\end{definition}

\begin{propriete}[Relation entre $t_{1/2}$ et $\lambda$ — démonstration guidée]
La période radioactive et la constante radioactive sont liées par :
\[ \boxed{t_{1/2} = \dfrac{\ln 2}{\lambda} \approx \dfrac{0{,}693}{\lambda}} \]
\end{propriete}

\begin{experience}[Démonstration guidée pas à pas]
\begin{enumerate}
\item \textbf{Traduction de la définition.} À $t = t_{1/2}$, on a $N(t_{1/2}) = \dfrac{N_0}{2}$, soit $N_0\,e^{-\lambda t_{1/2}} = \dfrac{N_0}{2}$.
\item \textbf{Simplification.} En divisant par $N_0$ : $e^{-\lambda t_{1/2}} = \dfrac{1}{2}$.
\item \textbf{Passage au logarithme.} Prenons le logarithme népérien des deux membres : $-\lambda t_{1/2} = \ln\!\left(\dfrac{1}{2}\right) = -\ln 2$.
\item \textbf{Conclusion.} En divisant par $-\lambda$ : $t_{1/2} = \dfrac{\ln 2}{\lambda} \approx \dfrac{0{,}693}{\lambda}$. \hfill$\blacksquare$
\end{enumerate}
\end{experience}

\begin{attention}[Analyse dimensionnelle et unités]
$\lambda$ s'exprime en $\mathrm{s^{-1}}$, donc $t_{1/2} = \ln 2 / \lambda$ s'exprime en secondes : tout est cohérent. Mais prends garde : si $\lambda$ est en $\mathrm{an^{-1}}$, alors $t_{1/2}$ obtenu est en \textbf{années}. Dans un calcul, $\lambda$ et $t_{1/2}$ doivent toujours être exprimés avec des unités de temps \emph{compatibles}.
\end{attention}

\begin{propriete}[Décroissance par moitiés successives]
Après $n$ périodes radioactives ($t = n\,t_{1/2}$), le nombre de noyaux restants est :
\[ N = \dfrac{N_0}{2^n}. \]
En effet : $N(n\,t_{1/2}) = N_0 e^{-\lambda n t_{1/2}} = N_0 e^{-n\ln 2} = N_0\,(e^{\ln 2})^{-n} = \dfrac{N_0}{2^n}$.
\end{propriete}

Ainsi, après 1 période il reste $N_0/2$, après 2 périodes $N_0/4$, après 3 périodes $N_0/8$, après 10 périodes $N_0/1024 \approx 0{,}001\,N_0$.

\begin{attention}[Erreur fréquente : la décroissance n'est jamais totale]
Beaucoup d'élèves écrivent : « après une période, la moitié des noyaux s'est désintégrée, donc après deux périodes, tout a disparu ». \textbf{C'est faux !} À chaque période, c'est la moitié des noyaux \emph{encore présents} qui se désintègre : après $2\,t_{1/2}$ il reste $N_0/4$, après $3\,t_{1/2}$ il reste $N_0/8$, etc. La fonction exponentielle ne s'annule jamais : en théorie, il reste toujours des noyaux radioactifs. En pratique, on considère qu'un échantillon est devenu inoffensif après une dizaine de périodes (il reste moins de $0{,}1\,\%$ des noyaux initiaux).
\end{attention}

\courssub{Ordres de grandeur : des horloges pour toutes les échelles de temps}

Les périodes radioactives s'étendent sur des gammes prodigieuses, d'une fraction de seconde à des milliards d'années. C'est cette diversité qui fait la richesse des applications : on choisit le nucléide adapté à la durée du phénomène étudié.

\begin{center}
\begin{tabularx}{\linewidth}{|l|c|c|X|}
\hline
\rowcolor{popblueL} \textbf{Nucléide} & \textbf{Type} & $\boldsymbol{t_{1/2}}$ & \textbf{Utilisation typique} \\
\hline
\ce{^{131}_{53}I} & $\beta^-$ & 8,0 jours & Traitement des cancers de la thyroïde (médecine nucléaire) \\
\hline
\ce{^{14}_{6}C} & $\beta^-$ & 5730 ans & Datation des objets archéologiques et des restes organiques \\
\hline
\ce{^{238}_{92}U} & $\alpha$ & $4{,}5\times 10^{9}$ ans & Datation des roches et de la Terre (géologie) \\
\hline
\end{tabularx}
\end{center}

\begin{popfigure}
\begin{center}
\begin{tikzpicture}[scale=1]
\draw[->,thick,popdark] (0,0) -- (11.5,0);
\foreach \x/\lab in {0/{$1$ j},2/{$10^2$ j},4/{$10^4$ j $\approx 27$ an},6/{$10^6$ j $\approx 2700$ an},8/{$10^8$ j},10/{$10^{10}$ j $\approx 2{,}7\times 10^7$ an}}{
  \draw[popdark] (\x,0.08) -- (\x,-0.08);
  \node[below,font=\scriptsize,align=center] at (\x,-0.1) {\lab};
}
\node[above,font=\small\itshape] at (5.75,1.35) {P\'eriodes radioactives (\'echelle logarithmique, en jours)};
\fill[poporange] (0.9,0) circle (2.2pt); \node[above,font=\scriptsize,poporange] at (0.9,0.08) {\ce{^{131}I} : 8 j};
\fill[popgreen] (6.32,0) circle (2.2pt); \node[above,font=\scriptsize,popgreen] at (6.32,0.08) {\ce{^{14}C} : 5730 an};
\fill[poppurple] (11.3,0) circle (2.2pt); \node[above,font=\scriptsize,poppurple,align=center] at (10.6,0.08) {\ce{^{238}U} : $4{,}5\times10^9$ an};
\end{tikzpicture}
\end{center}
\legende{Échelle logarithmique des périodes : chaque graduation correspond à une multiplication par 100. L'écart entre l'iode 131 et l'uranium 238 est de plus de 11 ordres de grandeur.}
\end{popfigure}

\begin{savaistu}[Une horloge universelle qui ne dépend que du noyau]
La demi-vie d'un nucléide est rigoureusement constante : ni la chaleur d'une forge, ni le froid d'un glacier, ni les champs électriques les plus intenses ne peuvent la modifier. C'est une véritable « horloge » gravée dans le cœur de la matière. Les géologues utilisent l'uranium 238 pour dater les roches (la Terre a environ $4{,}5$ milliards d'années, soit exactement une période de l'\ce{^{238}U} !), et les archéologues utilisent le carbone 14 pour dater le bois, les ossements ou les tissus des sites préhistoriques. Au Cameroun, la datation au carbone 14 a notamment permis d'étudier les sites de l'âge du fer du Sahel camerounais.
\end{savaistu}

\courssub{Méthodes et calculs autour de la loi de décroissance}

\begin{methode}[Déterminer graphiquement $t_{1/2}$ puis $\lambda$]
À partir d'une courbe expérimentale $N(t)$ (ou $A(t)$) :
\begin{enumerate}
\item Repérer l'ordonnée initiale $N_0$ (à $t = 0$).
\item Calculer $N_0/2$ et placer cette valeur sur l'axe des ordonnées.
\item Tracer l'horizontale d'ordonnée $N_0/2$ jusqu'à la courbe, puis abaisser la verticale : l'abscisse du point d'intersection est $t_{1/2}$.
\item \textbf{Vérification} : répéter l'opération de $N_0/2$ à $N_0/4$ ; on doit retrouver la même durée $t_{1/2}$. C'est la signature d'une décroissance exponentielle.
\item En déduire la constante radioactive : $\lambda = \dfrac{\ln 2}{t_{1/2}}$, en veillant à convertir $t_{1/2}$ en secondes pour obtenir $\lambda$ en $\mathrm{s^{-1}}$.
\end{enumerate}
\end{methode}

\begin{methode}[Calculer une date ou une durée à partir de la loi de décroissance]
Pour trouver la date $t$ à laquelle il reste $N$ noyaux (ou pour dater un échantillon) :
\begin{enumerate}
\item Écrire la loi $N = N_0\,e^{-\lambda t}$ et isoler l'exponentielle : $e^{-\lambda t} = \dfrac{N}{N_0}$.
\item Prendre le logarithme népérien : $-\lambda t = \ln\!\left(\dfrac{N}{N_0}\right)$.
\item Conclure : $t = \dfrac{1}{\lambda}\ln\!\left(\dfrac{N_0}{N}\right)$, avec $\lambda$ et $t$ dans des unités de temps compatibles.
\item Vérifier la cohérence : si $N < N_0/2$, on doit trouver $t > t_{1/2}$.
\end{enumerate}
\end{methode}

\begin{exemplebox}[Constante radioactive du carbone 14]
Le carbone 14 a une période $t_{1/2} = 5730$ ans. Calculons $\lambda$ en $\mathrm{an^{-1}}$ puis en $\mathrm{s^{-1}}$.
\begin{align*}
\lambda &= \dfrac{\ln 2}{t_{1/2}} = \dfrac{0{,}693}{5730} = 1{,}21\times 10^{-4}\ \mathrm{an^{-1}}.\\
\lambda &= \dfrac{1{,}21\times 10^{-4}}{365{,}25\times 24\times 3600} = \dfrac{1{,}21\times 10^{-4}}{3{,}156\times 10^{7}} = 3{,}83\times 10^{-12}\ \mathrm{s^{-1}}.
\end{align*}
Interprétation : chaque seconde, un noyau de carbone 14 n'a qu'environ $4$ chances sur mille milliards de se désintégrer... mais un gramme de carbone contient tellement de noyaux que la radioactivité reste mesurable !
\end{exemplebox}

\begin{exoresolu}[Décroissance de l'iode 131 en médecine]
Un hôpital reçoit un échantillon d'iode 131 ($t_{1/2} = 8{,}0$ jours) contenant $N_0 = 4{,}0\times 10^{15}$ noyaux.
\begin{enumerate}
\item Calculer la constante radioactive $\lambda$ en $\mathrm{s^{-1}}$.
\item Combien de noyaux reste-t-il après 24 jours ? Après 80 jours ?
\item Au bout de combien de temps reste-t-il moins de $1\,\%$ des noyaux initiaux ?
\end{enumerate}
\tcblower
\textbf{1.} Convertissons la période en secondes : $t_{1/2} = 8{,}0\times 24\times 3600 = 6{,}91\times 10^{5}\ \mathrm{s}$.
\[ \lambda = \dfrac{\ln 2}{t_{1/2}} = \dfrac{0{,}693}{6{,}91\times 10^{5}} = 1{,}0\times 10^{-6}\ \mathrm{s^{-1}}. \]
\textbf{2.} 24 jours correspondent à $n = \dfrac{24}{8{,}0} = 3$ périodes. Donc :
\[ N = \dfrac{N_0}{2^3} = \dfrac{4{,}0\times 10^{15}}{8} = 5{,}0\times 10^{14}\ \text{noyaux}. \]
80 jours correspondent à $n = 10$ périodes : $N = \dfrac{N_0}{2^{10}} = \dfrac{4{,}0\times 10^{15}}{1024} \approx 3{,}9\times 10^{12}$ noyaux.
\textbf{3.} On cherche $t$ tel que $N \leq 0{,}01\,N_0$, soit $e^{-\lambda t} \leq 0{,}01$ :
\[ -\lambda t \leq \ln(0{,}01) \quad\Longrightarrow\quad t \geq \dfrac{\ln 100}{\lambda} = \dfrac{4{,}61}{1{,}0\times 10^{-6}} = 4{,}61\times 10^{6}\ \mathrm{s} \approx 53\ \text{jours}. \]
(On pouvait aussi remarquer que $2^7 = 128 > 100$, donc il faut 7 périodes, soit $7\times 8 = 56$ jours — le léger écart vient de ce que $2^7 = 128$ correspond à $0{,}78\,\%$ et non exactement $1\,\%$.)
\end{exoresolu}

\begin{exercice}[Datation et décroissance]
Un échantillon de technétium 99m, utilisé en imagerie médicale, a une période $t_{1/2} = 6{,}0$ h. Un flacon contient $N_0 = 1{,}2\times 10^{13}$ noyaux à 8 h du matin.
\begin{enumerate}
\item Calculer $\lambda$ en $\mathrm{s^{-1}}$.
\item Combien de noyaux restera-t-il à 20 h le même jour ?
\item À quelle heure ne restera-t-il plus que $1{,}5\times 10^{11}$ noyaux ?
\end{enumerate}
\end{exercice}

\begin{corrige}[Exercice : Datation et décroissance]
\textbf{1.} $t_{1/2} = 6{,}0\times 3600 = 2{,}16\times 10^{4}\ \mathrm{s}$, donc $\lambda = \dfrac{\ln 2}{t_{1/2}} = \dfrac{0{,}693}{2{,}16\times 10^{4}} = 3{,}2\times 10^{-5}\ \mathrm{s^{-1}}$.

\textbf{2.} De 8 h à 20 h, il s'écoule $12$ h $= 2\,t_{1/2}$. Donc $N = \dfrac{N_0}{2^2} = \dfrac{1{,}2\times 10^{13}}{4} = 3{,}0\times 10^{12}$ noyaux.

\textbf{3.} On résout $N_0 e^{-\lambda t} = N$ :
\[ t = \dfrac{1}{\lambda}\ln\!\left(\dfrac{N_0}{N}\right) = \dfrac{1}{3{,}2\times 10^{-5}}\ln\!\left(\dfrac{1{,}2\times 10^{13}}{1{,}5\times 10^{11}}\right) = \dfrac{\ln 80}{3{,}2\times 10^{-5}} = \dfrac{4{,}38}{3{,}2\times 10^{-5}} \approx 1{,}37\times 10^{5}\ \mathrm{s} \approx 38\ \mathrm{h}. \]
Il restera $1{,}5\times 10^{11}$ noyaux environ 38 h après 8 h, soit vers 22 h le lendemain.
\end{corrige}

\begin{aretenir}[L'essentiel de la section]
\begin{itemize}
\item La décroissance radioactive obéit à l'équation différentielle $\dfrac{\mathrm{d}N}{\mathrm{d}t} = -\lambda N$, dont la solution est $N(t) = N_0\,e^{-\lambda t}$.
\item $\lambda$ (en $\mathrm{s^{-1}}$) est la probabilité de désintégration par unité de temps ; elle est caractéristique du nucléide.
\item La période $t_{1/2} = \dfrac{\ln 2}{\lambda}$ est la durée de disparition de la moitié des noyaux ; après $n$ périodes, $N = \dfrac{N_0}{2^n}$.
\item Pour trouver une durée : $t = \dfrac{1}{\lambda}\ln\!\left(\dfrac{N_0}{N}\right)$.
\item La décroissance n'est jamais totale : c'est une exponentielle, qui tend vers zéro sans l'atteindre.
\end{itemize}
\end{aretenir}

Maintenant que nous savons compter les noyaux restants, une question pratique demeure : en laboratoire ou à l'hôpital, on ne compte pas les noyaux un par un — on mesure le nombre de désintégrations par seconde à l'aide d'un compteur. Cette grandeur, appelée \emph{activité}, fera l'objet de la section suivante.

\coursec{Activité d'un échantillon radioactif}

Dans la section précédente, nous avons établi la loi de décroissance radioactive : le nombre de noyaux encore présents dans un échantillon décroît exponentiellement selon $N(t) = N_0 e^{-\lambda t}$, avec une période $t_{1/2} = \dfrac{\ln 2}{\lambda}$. Cette loi décrit le \emph{stock} de noyaux radioactifs. Mais en pratique, personne ne compte les noyaux un à un : ce que l'on mesure dans un laboratoire, dans un hôpital de Yaoundé ou de Douala, ou autour d'une source radioactive, c'est le \emph{rythme} des désintégrations, c'est-à-dire le nombre de noyaux qui se désintègrent chaque seconde. Cette grandeur, fondamentale en radioprotection et en médecine nucléaire, s'appelle l'\cle{activité}.

\courssub{Définition de l'activité}

\textbf{Intuition.} Imagine deux échantillons : l'un contient $10^{20}$ noyaux radioactifs, l'autre seulement $10^{10}$. Lequel est le plus « dangereux » ? Intuitivement, celui qui émet le plus de rayonnements par seconde, c'est-à-dire celui où il se produit le plus de désintégrations par unité de temps. Ce rythme de désintégration est précisément ce que mesure l'activité.

\begin{definition}[Activité d'un échantillon radioactif]
L'\cle{activité} $A$ d'un échantillon radioactif est le nombre moyen de désintégrations qui s'y produisent par unité de temps :
\[ A = -\dfrac{dN}{dt} = \lambda N \]
où $N$ est le nombre de noyaux radioactifs présents à l'instant $t$ et $\lambda$ la constante radioactive du nucléide (en $\mathrm{s^{-1}}$).

L'unité SI de l'activité est le \cle{becquerel} (symbole Bq) :
\[ 1~\mathrm{Bq} = 1~\text{désintégration par seconde} \]
\end{definition}

Quelques remarques essentielles pour bien comprendre cette définition :

\begin{itemize}
\item Le signe moins dans $-\dfrac{dN}{dt}$ traduit le fait que $N$ \emph{diminue} : $\dfrac{dN}{dt} < 0$, mais l'activité est une grandeur positive.
\item L'égalité $A = \lambda N$ s'obtient directement en dérivant la loi de décroissance : $\dfrac{dN}{dt} = -\lambda N_0 e^{-\lambda t} = -\lambda N$, d'où $A = -\dfrac{dN}{dt} = \lambda N$.
\item L'activité est une moyenne statistique : la radioactivité étant un phénomène aléatoire, le nombre de désintégrations comptées pendant une seconde fluctue légèrement autour de $A$.
\end{itemize}

\begin{savaistu}[Le curie, une unité historique]
Avant le becquerel, on utilisait le \cle{curie} (Ci), défini comme l'activité d'un gramme de radium 226 (l'élément découvert par Marie et Pierre Curie) :
\[ 1~\mathrm{Ci} = 3{,}7 \times 10^{10}~\mathrm{Bq} \]
Le curie est une unité énorme : on utilise plutôt le millicurie ($1~\mathrm{mCi} = 37~\mathrm{MBq}$) en médecine nucléaire. Dans les sujets de Bac, on peut encore rencontrer cette unité : sache convertir dans les deux sens. Par exemple, une source de $2~\mathrm{mCi}$ a une activité de $7{,}4 \times 10^{7}~\mathrm{Bq}$.
\end{savaistu}

\courssub{Décroissance de l'activité au cours du temps}

Puisque $A = \lambda N$ et que $N(t) = N_0 e^{-\lambda t}$, on obtient immédiatement :

\begin{propriete}[Loi de décroissance de l'activité]
L'activité d'un échantillon radioactif décroît exponentiellement avec la même constante de temps que le nombre de noyaux :
\[ A(t) = A_0 e^{-\lambda t} \qquad \text{avec} \qquad A_0 = \lambda N_0 \]
En particulier, l'activité est divisée par deux à chaque période radioactive : $A(t + t_{1/2}) = \dfrac{A(t)}{2}$.
\end{propriete}

C'est un point capital : $N(t)$ et $A(t)$ sont proportionnelles (le coefficient de proportionnalité est $\lambda$), donc les deux courbes ont exactement la même allure exponentielle. Après une période, il reste la moitié des noyaux \emph{et} l'activité a été divisée par deux. Après $n$ périodes, $A = \dfrac{A_0}{2^n}$.

\begin{popfigure}
\begin{center}
\begin{tikzpicture}
\begin{axis}[
  width=0.9\linewidth, height=6.5cm,
  xlabel={$t$ (en périodes $t_{1/2}$)}, ylabel={},
  xmin=0, xmax=6, ymin=0, ymax=1.15,
  xtick={0,1,2,3,4,5,6},
  ytick={0,0.25,0.5,0.75,1},
  yticklabels={0, 0{,}25, 0{,}50, 0{,}75, 1},
  grid=both, grid style={popline!40},
  legend style={at={(0.97,0.95)}, anchor=north east, font=\small},
  axis lines=left,
]
\addplot[popcurve, popblue, very thick, domain=0:6, samples=200] {exp(-0.6931*x)};
\addlegendentry{$N(t)/N_0$}
\addplot[popcurve, poporange, very thick, dashed, domain=0:6, samples=200] {exp(-0.6931*x)};
\addlegendentry{$A(t)/A_0$}
\addplot[popmuted, thin, domain=0:6, samples=2] {0.5};
\draw[popmuted, thin] (axis cs:1,0) -- (axis cs:1,0.5);
\draw[popmuted, thin] (axis cs:2,0) -- (axis cs:2,0.25);
\addplot[popmuted, thin, domain=0:6, samples=2] {0.25};
\node[font=\small, popdark] at (axis cs:3.6,0.62) {$A = \lambda N$ : m\^eme d\'ecroissance};
\end{axis}
\end{tikzpicture}
\end{center}
\legende{Les courbes $N(t)/N_0$ et $A(t)/A_0$ sont confondues : comme $A = \lambda N$, l'activité suit la même loi exponentielle que le nombre de noyaux. À chaque période $t_{1/2}$, les deux grandeurs sont divisées par deux.}
\end{popfigure}

\begin{attention}[Ne pas confondre activité et nombre de noyaux]
Trois erreurs reviennent sans cesse dans les copies :
\begin{enumerate}
\item \textbf{Confondre $A$ et $N$.} Le nombre de noyaux $N$ est sans unité (c'est un comptage), l'activité $A$ s'exprime en becquerels. Dire « l'activité est de $10^{15}$ noyaux » n'a aucun sens.
\item \textbf{Croire que $\lambda$ diminue avec le temps.} C'est faux : la constante radioactive $\lambda$ est une caractéristique du nucléide, elle reste \emph{constante}. C'est $N$ qui décroît, et donc $A = \lambda N$ décroît parce que $N$ décroît.
\item \textbf{Oublier que $A$ se divise par deux à chaque période.} Si $A_0 = 8 \times 10^{8}~\mathrm{Bq}$ et $t_{1/2} = 6~\mathrm{h}$, alors après $18~\mathrm{h} = 3\,t_{1/2}$, on a $A = A_0/2^3 = 10^{8}~\mathrm{Bq}$, et non $A_0/3$.
\end{enumerate}
\end{attention}

\courssub{Relation entre activité et masse de l'échantillon}

Pour calculer l'activité d'un échantillon dont on connaît la masse $m$, il faut d'abord déterminer le nombre de noyaux qu'elle contient. Un échantillon de nucléide de masse molaire $M$ et de masse $m$ contient une quantité de matière $n = \dfrac{m}{M}$, donc :
\[ N = n \times N_A = \dfrac{m}{M}\,N_A \]
où $N_A = 6{,}022 \times 10^{23}~\mathrm{mol^{-1}}$ est le nombre d'Avogadro. L'activité s'écrit alors :
\[ A = \lambda N = \lambda \, \dfrac{m\,N_A}{M} \qquad \text{avec} \qquad \lambda = \dfrac{\ln 2}{t_{1/2}} \]

Vérifions l'homogénéité : $\lambda$ est en $\mathrm{s^{-1}}$, $m/M$ en mol, $N_A$ en $\mathrm{mol^{-1}}$, donc $A$ est bien en $\mathrm{s^{-1}}$, c'est-à-dire en becquerels.

\begin{methode}[Calculer l'activité d'un échantillon de masse connue]
\begin{enumerate}
\item \textbf{Convertir la période en secondes} : $t_{1/2}$ doit être exprimée en secondes pour obtenir $\lambda$ en $\mathrm{s^{-1}}$.
\item \textbf{Calculer la constante radioactive} : $\lambda = \dfrac{\ln 2}{t_{1/2}}$.
\item \textbf{Calculer le nombre de noyaux} : $N = \dfrac{m}{M}\,N_A$, avec $m$ en grammes et $M$ en $\mathrm{g\,mol^{-1}}$ (numériquement, $M \approx A$ en $\mathrm{g\,mol^{-1}}$, où $A$ est le nombre de masse).
\item \textbf{Calculer l'activité} : $A = \lambda N$, en becquerels.
\item \textbf{Conclure avec un nombre raisonnable de chiffres significatifs} et l'unité (Bq, MBq, GBq...).
\end{enumerate}
\end{methode}

\begin{exoresolu}[Activité d'un gramme d'iode 131]
L'iode 131 ($\ce{^{131}_{53}I}$), utilisé en médecine nucléaire pour traiter les pathologies de la thyroïde, a une période radioactive $t_{1/2} = 8{,}0$ jours. Calculer l'activité d'un échantillon de masse $m = 1{,}0~\mathrm{g}$ d'iode 131. On donne $N_A = 6{,}02 \times 10^{23}~\mathrm{mol^{-1}}$ et $M(\ce{^{131}I}) = 131~\mathrm{g\,mol^{-1}}$.

\tcblower

\textbf{Étape 1 : conversion de la période en secondes.}
\[ t_{1/2} = 8{,}0~\mathrm{j} = 8{,}0 \times 24 \times 3600~\mathrm{s} = 6{,}91 \times 10^{5}~\mathrm{s} \]

\textbf{Étape 2 : constante radioactive.}
\[ \lambda = \dfrac{\ln 2}{t_{1/2}} = \dfrac{0{,}693}{6{,}91 \times 10^{5}~\mathrm{s}} = 1{,}00 \times 10^{-6}~\mathrm{s^{-1}} \]

\textbf{Étape 3 : nombre de noyaux dans 1,0 g.}
\[ N = \dfrac{m}{M}\,N_A = \dfrac{1{,}0~\mathrm{g}}{131~\mathrm{g\,mol^{-1}}} \times 6{,}02 \times 10^{23}~\mathrm{mol^{-1}} = 4{,}60 \times 10^{21}~\text{noyaux} \]

\textbf{Étape 4 : activité.}
\[ A = \lambda N = 1{,}00 \times 10^{-6} \times 4{,}60 \times 10^{21} = 4{,}6 \times 10^{15}~\mathrm{Bq} \]

\textbf{Conclusion.} Un seul gramme d'iode 131 présente une activité d'environ $4{,}6 \times 10^{15}~\mathrm{Bq}$, soit environ 4,6 millions de milliards de désintégrations par seconde (environ $1{,}2 \times 10^{5}$ curies !). Cet exemple montre à quel point le becquerel est une unité « petite » à l'échelle des échantillons macroscopiques, et pourquoi on manipule en médecine des masses infimes de radionucléides.
\end{exoresolu}

\courssub{Mesure de l'activité : le compteur Geiger-Müller}

Comment mesure-t-on concrètement une activité ? L'appareil le plus répandu est le \cle{compteur Geiger-Müller}. Son principe (à connaître comme complément culturel) est le suivant : un tube métallique rempli d'un gaz rare (argon ou hélium) sous faible pression contient un fil central porté à une haute tension (quelques centaines de volts) par rapport à la paroi. Lorsqu'un rayonnement ionisant pénètre dans le tube par une fenêtre mince (en mica), il ionise des atomes du gaz ; les électrons arrachés sont accélérés vers le fil et ionisent à leur tour d'autres atomes : c'est l'« avalanche » électronique. Chaque avalanche produit une impulsion de courant, amplifiée puis comptée — c'est le « tic-tic » caractéristique du compteur.

\begin{popfigure}
\begin{center}
\begin{tikzpicture}[scale=1.0, font=\small]
% Tube cylindrique
\draw[very thick, popblue] (0,-1.1) -- (7,-1.1) arc(-90:90:0.35 and 1.1) -- (0,1.1) arc(90:270:0.35 and 1.1);
% Fenêtre d'entrée
\draw[very thick, poporange, fill=poporangeL] (-0.35,-1.1) arc(270:90:0.35 and 1.1) -- cycle;
\node[poporange, below, align=center] at (-0.4,-1.5) {fen\^etre\\ (mica)};
% Fil central
\draw[very thick, poppink] (0,0) -- (7,0);
\node[poppink, above] at (3.5,0.15) {fil anodique ($+U$)};
% Gaz
\node[popmuted] at (3.5,-0.75) {gaz (argon) basse pression};
% Rayonnement entrant
\draw[-{Stealth[length=3mm]}, very thick, popgreen] (-2.3,0.5) -- (-0.5,0.1);
\node[popgreen, left] at (-2.3,0.5) {rayonnement};
% Avalanche
\foreach \p in {(2.2,0.45),(2.6,-0.5),(3.1,0.6),(3.6,-0.35)}{
  \node[popgold] at \p {\scriptsize $\bullet$};
}
\draw[-{Stealth[length=2mm]}, popgold] (2.6,0.35) -- (2.7,0.05);
\node[popgold, align=center] at (5.6,0.75) {avalanche\\ d'ionisation};
% Circuit extérieur
\draw[very thick, popdark] (7,0) -- (8.2,0) -- (8.2,-2.3) -- (1.5,-2.3);
\draw[very thick, popdark] (1.5,-2.3) -- (0,-2.3) -- (0,-1.1);
\draw[popdark, very thick] (3.4,-2.45) -- (3.4,-2.15);
\draw[popdark, very thick] (3.7,-2.5) -- (3.7,-2.1);
\node[popdark, below] at (3.55,-2.55) {$U$};
% Compteur
\draw[thick, poppurple, fill=poppurpleL] (7.4,-2.75) rectangle (9.0,-1.85);
\node[poppurple] at (8.2,-2.3) {compteur};
\draw[very thick, popdark] (8.2,-1.1) -- (8.2,-1.85);
\node[popmuted, right, align=left] at (9.1,-2.3) {impulsions\\ compt\'ees};
\end{tikzpicture}
\end{center}
\legende{Schéma de principe d'un compteur Geiger-Müller : le rayonnement ionise le gaz contenu dans le tube ; l'avalanche d'électrons vers le fil anodique produit une impulsion électrique comptée par l'électronique.}
\end{popfigure}

\begin{attention}[Le compteur ne mesure pas directement l'activité totale]
Le compteur Geiger-Müller ne détecte qu'une fraction des désintégrations : seuls les rayonnements qui pénètrent dans le tube sont comptés, et la source n'émet pas seulement vers le détecteur (elle émet dans tout l'espace). Le comptage (en coups par seconde) est donc \emph{proportionnel} à l'activité, mais l'exploitation quantitative complète nécessite un étalonnage. Au Bac, on retient : le compteur Geiger-Müller permet de \emph{détecter} un rayonnement et de \emph{suivre l'évolution} d'une activité au cours du temps.
\end{attention}

\courssub{Applications : radioprotection et médecine nucléaire}

\textbf{Radioprotection.} Puisque l'activité décroît exponentiellement, le danger présenté par une source radioactive diminue spontanément avec le temps. C'est le principe de la gestion des déchets radioactifs à vie courte : on les stocke en attendant que leur activité devienne négligeable (règle pratique : après une dizaine de périodes, l'activité est divisée par environ $2^{10} \approx 10^{3}$). À l'inverse, les déchets à vie longue (comme le plutonium 239, $t_{1/2} = 24\,000$ ans) posent un problème de stockage sur des durées considérables.

\textbf{Choix des traceurs en médecine.} En imagerie médicale (scintigraphie), on injecte au patient un \cle{traceur} radioactif qui se fixe sur l'organe à étudier. Pour limiter l'irradiation du patient, on choisit un nucléide à \emph{courte période} : le technétium 99m ($t_{1/2} = 6{,}0~\mathrm{h}$), le plus utilisé au monde, permet de réaliser l'examen avec une activité élevée (donc de bonnes images), puis son activité devient rapidement négligeable — divisée par 16 en 24 heures. À l'inverse, l'iode 131 ($t_{1/2} = 8$ jours), émetteur $\beta^-$ plus pénétrant, est réservé à la \emph{radiothérapie} de la thyroïde : on veut alors irradier durablement les cellules malades pour les détruire.

\begin{savaistu}[Ton corps est radioactif !]
Toi aussi, tu es naturellement radioactif. Le corps humain contient du potassium 40 (un isotope radioactif présent dans tout le potassium, notamment des bananes !) et du carbone 14, entretenu par l'alimentation et la respiration. L'activité totale d'un adulte d'environ 70 kg est de l'ordre de $8\,000~\mathrm{Bq}$ : ton corps subit environ huit mille désintégrations par seconde, sans aucune conséquence, car cette activité est faible et fait partie de notre environnement naturel depuis toujours. Cela relativise la notion de « zéro radioactivité » : ce qui compte en radioprotection, ce sont les \emph{doses} reçues et leur comparaison avec le bruit de fond naturel (environ $2$ à $3~\mathrm{mSv}$ par an au Cameroun comme ailleurs).
\end{savaistu}

\begin{aretenir}[L'essentiel sur l'activité]
\begin{itemize}
\item L'activité est le nombre moyen de désintégrations par seconde : $A = -\dfrac{dN}{dt} = \lambda N$, en becquerels ($1~\mathrm{Bq} = 1$ désintégration/s).
\item Ancienne unité : $1~\mathrm{Ci} = 3{,}7 \times 10^{10}~\mathrm{Bq}$.
\item L'activité décroît comme $N$ : $A(t) = A_0 e^{-\lambda t}$ avec $A_0 = \lambda N_0$ ; elle est divisée par deux à chaque période.
\item Pour une masse $m$ de nucléide : $A = \lambda\,\dfrac{m\,N_A}{M}$ avec $\lambda = \dfrac{\ln 2}{t_{1/2}}$ (période en secondes).
\item Le compteur Geiger-Müller détecte les rayonnements par ionisation d'un gaz et permet de suivre l'évolution de l'activité.
\item En médecine, on choisit des traceurs à courte période pour limiter l'irradiation du patient.
\end{itemize}
\end{aretenir}

\begin{exercice}[Activité du technétium 99m]
Un hôpital prépare, pour une scintigraphie, une dose de technétium 99m ($t_{1/2} = 6{,}0~\mathrm{h}$ ; $M = 99~\mathrm{g\,mol^{-1}}$) d'activité $A_0 = 400~\mathrm{MBq}$ à 6 h du matin.
\begin{enumerate}
\item Calculer la masse de technétium 99m contenue dans cette dose.
\item L'examen n'a lieu qu'à 15 h. Quelle est alors l'activité de la dose ? Commenter.
\end{enumerate}
\end{exercice}

\begin{corrige}[Exercice : Activité du technétium 99m]
\begin{enumerate}
\item Constante radioactive : $\lambda = \dfrac{\ln 2}{t_{1/2}} = \dfrac{0{,}693}{6{,}0 \times 3600~\mathrm{s}} = 3{,}21 \times 10^{-5}~\mathrm{s^{-1}}$.\\
Nombre de noyaux : $N_0 = \dfrac{A_0}{\lambda} = \dfrac{4{,}00 \times 10^{8}}{3{,}21 \times 10^{-5}} = 1{,}25 \times 10^{13}$ noyaux.\\
Masse : $m = \dfrac{N_0\,M}{N_A} = \dfrac{1{,}25 \times 10^{13} \times 99}{6{,}02 \times 10^{23}} = 2{,}1 \times 10^{-9}~\mathrm{g}$, soit environ $2~\mathrm{ng}$. Une masse infime suffit !
\item Entre 6 h et 15 h, il s'écoule $\Delta t = 9{,}0~\mathrm{h} = 1{,}5\,t_{1/2}$. Donc :
\[ A = A_0 \times 2^{-1{,}5} = 400 \times 0{,}354 \approx 141~\mathrm{MBq} \]
L'activité a chuté de près de deux tiers en 9 heures : en médecine nucléaire, les doses doivent être préparées au plus près de l'heure de l'examen, et les hôpitaux se font livrer le technétium quotidiennement.
\end{enumerate}
\end{corrige}

\coursec{Bilan énergétique et applications}

Dans la section précédente, nous avons appris à quantifier la radioactivité d'un échantillon grâce à son activité $A = \lambda N$, exprimée en becquerels. Mais une question fondamentale reste en suspens : d'où vient l'énergie dégagée lors d'une désintégration ? Pourquoi certaines désintégrations sont-elles possibles et d'autres non ? Et surtout, comment l'humanité a-t-elle su transformer ce phénomène naturel en outils extraordinaires : dater un ossement vieux de dix mille ans, soigner un cancer à l'hôpital de Yaoundé, ou déterminer l'âge des roches du socle camerounais ? C'est tout l'objet de cette dernière section.

\courssub{Bilan énergétique d'une désintégration}

\courssub{L'idée intuitive : la masse est un réservoir d'énergie}

Considère une désintégration $\alpha$, par exemple celle du radium 226 :
\[
\ce{^{226}_{88}Ra -> ^{222}_{86}Rn + ^{4}_{2}He}
\]
Si l'on pèse avec une précision extrême les noyaux avant et après la réaction, on fait une découverte troublante : la somme des masses des produits est \emph{légèrement inférieure} à la masse du noyau père. De la masse a « disparu ». Or, d'après la célèbre relation d'Einstein $E = mc^2$, la masse est une forme d'énergie. La masse manquante n'a pas disparu : elle a été \emph{convertie en énergie}, emportée par les produits de la désintégration.

C'est exactement l'inverse du raisonnement tenu pour l'énergie de liaison : assembler un noyau à partir de nucléons séparés libère de l'énergie (défaut de masse) ; ici, un noyau instable se transforme en des noyaux \emph{plus stables}, c'est-à-dire plus liés, et la différence de masse correspondante est libérée.

\begin{propriete}[Énergie libérée par une désintégration]
Lors d'une désintégration $\ce{^{A}_{Z}X -> ^{A'}_{Z'}Y + p}$ (où $p$ désigne la particule émise : $\alpha$, $\beta^{-}$, $\beta^{+}$...), l'énergie libérée vaut :
\[
E = \left[ m(\ce{X}) - m(\ce{Y}) - m(p) \right] \cdot c^2 = \Delta m \cdot c^2
\]
où $\Delta m = m_{\text{avant}} - m_{\text{après}}$ est la \cle{perte de masse} du système, exprimée en kilogrammes (ou en unités de masse atomique u, avec $1\ \text{u} \cdot c^2 = \num{931,5}\ \text{MeV}$), et $c = \num{3,00e8}\ \text{m.s}^{-1}$.
\end{propriete}

Vérifions l'homogénéité de cette relation par analyse dimensionnelle : $[\Delta m \cdot c^2] = \text{kg} \cdot (\text{m.s}^{-1})^2 = \text{kg.m}^2\text{.s}^{-2} = \text{J}$. C'est bien une énergie.

\begin{attention}[Condition de spontanéité]
Une désintégration n'est \cle{spontanée} que si elle libère de l'énergie, c'est-à-dire si :
\[
E > 0 \quad \Longleftrightarrow \quad \Delta m > 0 \quad \Longleftrightarrow \quad m_{\text{avant}} > m_{\text{après}}
\]
Si la masse totale des produits est supérieure à celle du noyau père, la réaction ne peut pas se produire spontanément. Erreur fréquente au Bac : inverser la différence ($m_{\text{après}} - m_{\text{avant}}$) et conclure à tort qu'une réaction impossible est spontanée, ou trouver une énergie négative sans s'alarmer. Une énergie libérée négative signifie simplement que la désintégration ne peut pas avoir lieu seule.
\end{attention}

\begin{attention}[Masses nucléaires ou masses atomiques ?]
Les tables fournissent souvent des \emph{masses atomiques} (noyau + électrons). Pour la désintégration $\alpha$, on peut utiliser directement les masses atomiques : le noyau père possède $Z$ électrons, et les produits (atome fils à $Z-2$ électrons + atome d'hélium à 2 électrons) en possèdent aussi $Z$ au total ; les masses d'électrons se compensent. Pour la désintégration $\beta^{-}$, la compensation est également exacte avec les masses atomiques. En revanche, pour la désintégration $\beta^{+}$, l'utilisation des masses atomiques impose la condition $m(\ce{X}) > m(\ce{Y}) + 2\,m_e$ : retiens simplement que la désintégration $\beta^{+}$ exige une perte de masse plus importante. Au Bac, l'énoncé précise toujours les masses à utiliser : lis-le attentivement.
\end{attention}

\courssub{Sous quelle forme cette énergie est-elle emportée ?}

L'énergie $E = \Delta m \cdot c^2$ ne reste pas « dans » les noyaux : elle se répartit entre :

\begin{itemize}
\item l'\cle{énergie cinétique} de la particule émise (noyau $\alpha$, électron $\beta^{-}$ ou positon $\beta^{+}$), qui emporte généralement la plus grande part car elle est beaucoup plus légère que le noyau fils ;
\item l'énergie cinétique de \emph{recul} du noyau fils (conservation de la quantité de mouvement : le noyau père étant au repos, les produits partent en sens opposés) ;
\item un éventuel \cle{rayonnement $\gamma$}, émis lorsque le noyau fils est produit dans un état excité et redescend vers son état fondamental : $E_\gamma = h\nu$.
\end{itemize}

On écrit donc le bilan :
\[
\Delta m \cdot c^2 = E_c(\text{particule}) + E_c(\text{recul du noyau fils}) + E_\gamma
\]

\begin{popfigure}
\begin{tikzpicture}[scale=1.0]
% Avant
\node[draw, circle, fill=popblue!40, minimum size=1.1cm] (pere) at (0,2.6) {\ce{^{A}_{Z}X}};
\node[below=0.15cm of pere, text=popdark] {\small Noyau père au repos};
\node[below=0.55cm of pere, text=popmuted] {\small Énergie : $m(\ce{X})\,c^2$};
% Flèche
\draw[->, very thick, poporange] (1.1,2.6) -- (2.6,2.6) node[midway, above, text=poporange] {\small désintégration};
% Après
\node[draw, circle, fill=popgreen!40, minimum size=0.85cm] (fils) at (4.6,2.6) {\small \ce{Y}};
\node[draw, circle, fill=poppink!50, minimum size=0.45cm] (part) at (6.6,2.6) {};
\node[right=0.05cm of part, text=popdark] {\small particule $\alpha$ ou $\beta$};
\draw[->, thick, poppink] (4.95,2.75) -- (5.9,2.75) node[midway, above, sloped] {\footnotesize $\vec{v}$};
\draw[->, thick, poppurple, decorate, decoration={snake, amplitude=0.8mm}] (4.6,2.1) -- (5.8,1.4) node[right, text=poppurple] {\small photon $\gamma$};
% Niveaux d'énergie
\draw[very thick, popblue] (-0.6,0.2) -- (1.6,0.2) node[right, text=popdark] {\small $m(\ce{X})\,c^2$};
\draw[very thick, popgreen] (3.4,0.2) -- (5.6,0.2) node[right, text=popdark] {\small $[m(\ce{Y}) + m(p)]\,c^2$};
\draw[<->, thick, poporange] (2.9,0.2) -- (2.9,-0.9);
\draw[very thick, poporange, dashed] (3.4,-0.9) -- (5.6,-0.9);
\node[text=poporange] at (2.9,-1.25) {\small $\Delta m \cdot c^2$};
\node[text=popmuted, align=center] at (7.6,-0.5) {\small convertie en $E_c$\\ \small des produits $+\ E_\gamma$};
\end{tikzpicture}
\legende{Diagramme énergétique d'une désintégration : la perte de masse $\Delta m$ correspond à une énergie $\Delta m \cdot c^2$ emportée sous forme d'énergie cinétique des produits et de rayonnement $\gamma$.}
\end{popfigure}

\begin{exoresolu}[Énergie libérée par la désintégration $\alpha$ du radium 226]
Énoncé. Le radium 226 se désintègre selon : $\ce{^{226}_{88}Ra -> ^{222}_{86}Rn + ^{4}_{2}He}$. On donne les masses atomiques : $m(\ce{Ra}) = \num{226,0254}\ \text{u}$, $m(\ce{Rn}) = \num{222,0176}\ \text{u}$, $m(\ce{He}) = \num{4,0026}\ \text{u}$, et $1\ \text{u} \cdot c^2 = \num{931,5}\ \text{MeV}$. Calculer l'énergie libérée en MeV, puis en joules. Conclure sur la spontanéité.
\tcblower
Solution. Calculons d'abord la perte de masse (les masses atomiques conviennent ici, les électrons se compensant) :
\begin{align*}
\Delta m &= m(\ce{Ra}) - m(\ce{Rn}) - m(\ce{He}) \\
&= 226,0254 - 222,0176 - 4,0026 \\
&= \num{0,0052}\ \text{u}
\end{align*}
L'énergie libérée vaut alors :
\[
E = \Delta m \cdot c^2 = 0,0052 \times 931,5 \approx \num{4,84}\ \text{MeV}
\]
Conversion en joules, avec $1\ \text{eV} = 1,60 \times 10^{-19}\ \text{J}$ :
\[
E = 4,84 \times 10^6 \times 1,60 \times 10^{-19} \approx \num{7,7e-13}\ \text{J}
\]
Comme $E > 0$ (diminution de masse), la désintégration est bien spontanée. Cette énergie, minuscule à l'échelle d'un noyau, devient considérable à l'échelle macroscopique : un gramme de radium contient environ $N = \dfrac{1}{226} \times 6,02 \times 10^{23} \approx 2,7 \times 10^{21}$ noyaux !
\end{exoresolu}

\courssub{Application 1 : la datation au carbone 14}

\courssub{Principe : une horloge naturelle dans la matière vivante}

Le carbone 14 est un isotope radioactif $\beta^{-}$ du carbone, de période $t_{1/2} = \num{5730}$ ans. Son génie, pour l'archéologue, tient à trois faits :

\begin{enumerate}
\item \textbf{Il est continuellement formé dans la haute atmosphère} : les neutrons du rayonnement cosmique transmutent l'azote 14 selon $\ce{^{14}_{7}N + ^{1}_{0}n -> ^{14}_{6}C + ^{1}_{1}H}$.
\item \textbf{Il s'incorpore aux êtres vivants} : le \ce{^{14}C} s'oxyde en \ce{CO2}, absorbé par les plantes (photosynthèse), puis par les animaux qui les mangent. Tant que l'organisme est vivant, il échange du carbone avec l'atmosphère : la proportion de \ce{^{14}C} dans son corps reste \emph{constante}, égale à celle de l'atmosphère. Son activité spécifique est $A_0 \approx \num{13,6}$ désintégrations par minute et par gramme de carbone.
\item \textbf{À la mort, les échanges cessent} : le \ce{^{14}C} n'est plus renouvelé et décroît selon la loi $A(t) = A_0\, e^{-\lambda t}$. L'activité mesurée aujourd'hui renseigne donc sur la date de la mort.
\end{enumerate}

\begin{popfigure}
\begin{tikzpicture}[scale=0.95]
% Étape 1 : formation
\node[draw, rounded corners, fill=popblueL, text width=3.1cm, align=center] (form) at (0,0) {\small \textbf{1. Formation}\\ \footnotesize Rayonnement cosmique : $\ce{^{14}N} + n \to \ce{^{14}C}$};
% Étape 2 : absorption
\node[draw, rounded corners, fill=popgreenL, text width=3.1cm, align=center] (abs) at (5.2,0) {\small \textbf{2. Vie}\\ \footnotesize \ce{CO2} absorbé (photosynthèse) : $A = A_0$ constante};
% Étape 3 : mort
\node[draw, rounded corners, fill=poporangeL, text width=3.1cm, align=center] (mort) at (10.4,0) {\small \textbf{3. Mort}\\ \footnotesize Plus d'échanges : $A = A_0 e^{-\lambda t}$};
\draw[->, very thick, popblue] (form) -- (abs);
\draw[->, very thick, popgreen] (abs) -- (mort);
% Courbe de décroissance
\begin{scope}[shift={(5.2,-3.6)}, scale=0.9]
\draw[->, thick] (0,0) -- (6,0) node[right] {\footnotesize $t$};
\draw[->, thick] (0,0) -- (0,2.4) node[above] {\footnotesize $A$};
\draw[domain=0:5.6, samples=100, very thick, poporange] plot (\x, {2*exp(-0.5*\x)});
\node[left] at (0,2) {\footnotesize $A_0$};
\draw[dashed, popmuted] (1.386,1) -- (1.386,0) node[below] {\footnotesize $t_{1/2}$};
\draw[dashed, popmuted] (0,1) -- (1.386,1);
\node[left] at (0,1) {\footnotesize $A_0/2$};
\node[align=center, text=popmuted] at (3,-0.7) {\footnotesize après la mort};
\end{scope}
\end{tikzpicture}
\legende{Le cycle du carbone 14 : formation dans la haute atmosphère, équilibre chez les êtres vivants, puis décroissance exponentielle après la mort.}
\end{popfigure}

\begin{methode}[Dater un échantillon au carbone 14]
\begin{enumerate}
\item \textbf{Mesurer} l'activité $A$ de l'échantillon (en \ce{^{14}C}) à l'aide d'un compteur.
\item \textbf{Connaître} l'activité $A_0$ d'un échantillon vivant équivalent (même masse de carbone), ou l'activité spécifique de référence.
\item \textbf{Exploiter} la loi de décroissance $A = A_0 e^{-\lambda t}$, d'où :
\[
t = \frac{1}{\lambda}\, \ln\!\left(\frac{A_0}{A}\right) = \frac{t_{1/2}}{\ln 2}\, \ln\!\left(\frac{A_0}{A}\right)
\]
car $\lambda = \dfrac{\ln 2}{t_{1/2}}$.
\item \textbf{Vérifier} la cohérence : l'âge trouvé doit rester inférieur à une dizaine de périodes pour que la méthode soit fiable.
\end{enumerate}
\end{methode}

\begin{exemplebox}[Âge d'un charbon de bois préhistorique]
Un charbon de bois trouvé dans un site archéologique présente une activité en \ce{^{14}C} égale au \emph{quart} de celle d'un bois vivant de même masse : $A = A_0/4$. Quel est son âge ?

Comme $A/A_0 = 1/4 = (1/2)^2$, l'échantillon a vécu exactement \textbf{deux périodes} :
\[
t = 2 \times t_{1/2} = 2 \times 5730 = \num{11460}\ \text{ans}
\]
Vérification par la formule : $t = \dfrac{5730}{\ln 2}\,\ln 4 = \dfrac{5730}{0,693} \times 1,386 \approx \num{11460}$ ans. Le feu qui a produit ce charbon a été allumé il y a environ onze mille cinq cents ans, au début de l'Holocène.
\end{exemplebox}

\begin{attention}[Limites de la datation au carbone 14]
La datation au \ce{^{14}C} n'est fiable que jusqu'à environ \cle{10 périodes}, soit $\approx \num{57000}$ ans. Au-delà, l'activité résiduelle ($A_0/2^{10} \approx A_0/1000$) est trop faible pour être distinguée du bruit de fond. Erreur classique au Bac : appliquer la formule à un échantillon vieux de plusieurs millions d'années ! Pour les âges géologiques, on utilise d'autres couples radioactifs à longue période : uranium-plomb ($t_{1/2}(\ce{^{238}U}) = 4,5$ milliards d'années) ou potassium-argon ($t_{1/2}(\ce{^{40}K}) = 1,25$ milliard d'années). Autre piège : confondre $A_0/A$ et $A/A_0$ dans le logarithme — un âge négatif doit te mettre la puce à l'oreille.
\end{attention}

\begin{savaistu}[Willard Libby, prix Nobel de chimie 1960]
C'est le chimiste américain Willard Libby qui mit au point la datation au carbone 14 à la fin des années 1940, travail couronné par le prix Nobel de chimie en 1960. Sa méthode a révolutionné l'archéologie : pour la première fois, on pouvait dater objectivement un ossement, un tissu ou un charbon de bois. En Afrique, elle a permis de dater les peintures rupestres du Sahara et les sites de l'âge du fer, contribuant à écrire l'histoire ancienne du continent, y compris celle des premières métallurgies de la région camerounaise.
\end{savaistu}

\courssub{Application 2 : la radioactivité en médecine}

La radioactivité, loin d'être seulement un danger, est devenue un outil médical irremplaçable, y compris dans les grands hôpitaux camerounais (Centre Pasteur, hôpitaux universitaires de Yaoundé et Douala).

\textbf{La scintigraphie (diagnostic).} On injecte au patient un \cle{traceur} radioactif, le plus souvent le technétium 99m ($t_{1/2} \approx 6\ \text{h}$), émetteur $\gamma$. Ce noyau se fixe préférentiellement sur certains organes (os, thyroïde, reins). Une caméra à scintillations détecte les photons $\gamma$ émis et reconstitue une image de l'organe : on \emph{voit} son fonctionnement sans ouvrir le patient. La période courte est un atout : le patient n'est exposé que quelques heures.

\textbf{La radiothérapie (traitement).} On utilise au contraire le pouvoir destructeur des rayonnements pour tuer des cellules cancéreuses. La source classique est le cobalt 60 ($t_{1/2} \approx 5,3$ ans), émetteur $\beta^{-}$ et $\gamma$ très énergétiques ($E_\gamma \approx 1,2\ \text{MeV}$). Les faisceaux, soigneusement collimatés et orientés, convergent sur la tumeur : chaque tissu sain traversé reçoit une dose faible, mais la tumeur, au point de convergence, reçoit la dose totale.

\textbf{Les traceurs biologiques.} En agronomie comme en physiologie, on suit le devenir d'une substance (engrais, hormone, médicament) en y incorporant un isotope radioactif détectable en quantités infimes. C'est ainsi qu'on étudie l'absorption des phosphates par les cultures de maïs ou de cacao.

\begin{attention}[Radiothérapie : ne pas tout confondre]
Ne confonds pas scintigraphie et radiothérapie : dans le premier cas, on cherche une dose \emph{minimale} pour une image (diagnostic) ; dans le second, une dose \emph{maximale ciblée} pour détruire des cellules (thérapie). Et rappelle-toi : la radioprotection repose sur trois principes — \emph{distance} (la dose diminue avec le carré de la distance), \emph{durée} d'exposition réduite, et \emph{écrans} (plomb pour les $\gamma$ et $\beta$, une simple feuille de papier arrête les $\alpha$).
\end{attention}

\courssub{Application 3 : autres applications utiles au Bac}

\textbf{Datation géologique uranium-plomb.} L'uranium 238 se transforme, par une cascade de désintégrations, en plomb 206 stable. En mesurant le rapport $\ce{^{206}Pb}/\ce{^{238}U}$ dans un cristal de zircon, on détermine l'âge de la roche. C'est ainsi qu'on a daté le socle granitique précambrien d'Afrique centrale, sur lequel repose le Cameroun, à plus de 2 milliards d'années.

\textbf{Stérilisation et conservation.} L'irradiation $\gamma$ (cobalt 60) stérilise le matériel médical à usage unique (seringues, compresses) et prolonge la conservation de certains aliments en détruisant bactéries et moisissures — sans rendre les produits radioactifs, car l'irradiation ne crée pas de noyaux instables dans la matière aux énergies utilisées.

\textbf{Détecteurs de fumée.} Beaucoup de détecteurs domestiques contiennent une microsource d'américium 241, émetteur $\alpha$. Les particules $\alpha$ ionisent l'air entre deux électrodes, créant un courant permanent ; quand de la fumée pénètre dans la chambre, elle absorbe les ions, le courant chute et l'alarme se déclenche. Une application discrète mais quotidienne de la physique nucléaire.

\begin{exercice}[Datation d'un échantillon de bois]
Un échantillon de bois prélevé sur une pirogue ancienne a une activité en \ce{^{14}C} de $A = \num{6,8}$ désintégrations par minute et par gramme de carbone. Un bois vivant de référence donne $A_0 = \num{13,6}$ dés.min$^{-1}$.g$^{-1}$. On donne $t_{1/2}(\ce{^{14}C}) = \num{5730}$ ans. Déterminer l'âge de la pirogue.
\end{exercice}

\begin{corrige}[Exercice : datation d'un échantillon de bois]
On applique la loi de décroissance de l'activité :
\[
t = \frac{t_{1/2}}{\ln 2}\, \ln\!\left(\frac{A_0}{A}\right) = \frac{5730}{\ln 2}\, \ln\!\left(\frac{13,6}{6,8}\right) = \frac{5730}{\ln 2}\, \ln 2 = 5730\ \text{ans}
\]
Le rapport $A_0/A = 2$ correspond exactement à une période : la pirogue a environ $\num{5730}$ ans.
\end{corrige}

\begin{aretenir}[L'essentiel de la section]
\begin{itemize}
\item L'énergie libérée par une désintégration vaut $E = \Delta m \cdot c^2$, où $\Delta m = m_{\text{avant}} - m_{\text{après}}$ ; elle est emportée en énergie cinétique des produits et en rayonnement $\gamma$.
\item La désintégration n'est spontanée que si $\Delta m > 0$ (diminution de masse).
\item Datation au \ce{^{14}C} : $t = \dfrac{t_{1/2}}{\ln 2}\,\ln\!\left(\dfrac{A_0}{A}\right)$, fiable jusqu'à environ $\num{57000}$ ans ; au-delà, couples U-Pb ou K-Ar.
\item Médecine : scintigraphie au technétium 99m (diagnostic), radiothérapie au cobalt 60 (traitement), traceurs biologiques.
\end{itemize}
\end{aretenir}

\newpage

\coursec{Activité d'intégration}

\courssub{Objectifs de l'activité}

À l'issue de cette activité d'intégration, tu seras capable de :
\begin{itemize}
\item écrire et équilibrer l'équation d'une désintégration $\beta^-$ en appliquant les lois de conservation de Soddy ;
\item calculer une constante radioactive $\lambda$ à partir de la période $t_{1/2}$ et réciproquement ;
\item exploiter la loi de décroissance $N(t) = N_0\,e^{-\lambda t}$ sous sa forme « activité » $A(t) = A_0\,e^{-\lambda t}$ pour dater un échantillon archéologique ;
\item croiser une détermination par le calcul avec une lecture graphique sur une courbe de décroissance ;
\item calculer une activité résiduelle à long terme et discuter de manière argumentée les limites de la méthode de datation au carbone 14 ;
\item rédiger une conclusion scientifique claire, destinée à des partenaires non physiciens (ici, des archéologues).
\end{itemize}

\courssub{Situation}

\textbf{Quel âge ont les vestiges du site de Foumban ?}

Foumban, chef-lieu du département du Noun dans la région de l'Ouest, est l'ancienne capitale du royaume bamoun. Lors de fouilles récentes menées près du palais royal, une équipe d'archéologues de l'Université de Dschang a mis au jour, dans une couche de foyer ancien, des fragments de \cle{charbon de bois} associés à des poteries. Pour dater l'occupation du site, un échantillon de ce charbon est envoyé à un laboratoire de datation.

Tu fais partie du groupe d'élèves de Terminale C invité, dans le cadre d'un projet interdisciplinaire physique--histoire, à exploiter le \emph{rapport de laboratoire} ci-dessous et à rédiger la conclusion qui sera transmise aux archéologues.

\begin{center}
\begin{tabularx}{\linewidth}{|l|X|}
\hline
\rowcolor{popblueL} \multicolumn{2}{|c|}{\textbf{Rapport de laboratoire -- Datation par le carbone 14}} \\
\hline
Échantillon & Charbon de bois, site de Foumban (fouille F-23, couche 4) \\
\hline
Méthode & Mesure de l'activité du carbone 14 par comptage $\beta^-$ \\
\hline
Activité massique mesurée & $a = \SI{0,115}{Bq/g}$ de carbone \\
\hline
Activité massique de référence (bois vivant actuel) & $a_0 = \SI{0,230}{Bq/g}$ de carbone \\
\hline
Période du carbone 14 & $t_{1/2} = \SI{5730}{ans}$ \\
\hline
Noyau fils & Azote 14 : \ce{^{14}_{7}N} \\
\hline
\end{tabularx}
\end{center}

\textbf{Rappel de principe.} Le carbone 14 est produit en permanence dans la haute atmosphère par interaction des rayons cosmiques avec l'azote. Tant qu'un végétal est vivant, les échanges avec l'atmosphère (photosynthèse) maintiennent constante la proportion de \ce{^{14}C} dans ses tissus : son activité massique vaut $a_0$. À la mort de l'arbre (ici, lors de la fabrication du feu qui a produit le charbon), ces échanges cessent : le \ce{^{14}C} n'est plus renouvelé et son activité décroît selon la loi de décroissance radioactive. Mesurer $a < a_0$ permet donc de remonter à la date de la mort du végétal.

\courssub{Consignes}

Travail en groupes de trois ou quatre. Chaque groupe rédige un compte rendu soigné ; la question 6 fera l'objet d'une restitution orale.

\begin{enumerate}
\item \textbf{Équation de la désintégration.} Le carbone 14 est un émetteur $\beta^-$. Écrire l'équation de sa désintégration en précisant les lois de conservation utilisées, et identifier la particule émise.
\item \textbf{Constante radioactive.} Calculer la constante radioactive $\lambda$ du carbone 14 en an$^{-1}$, puis en s$^{-1}$. Vérifier que la valeur de $t_{1/2}$ annoncée par le laboratoire est cohérente.
\item \textbf{Datation par le calcul.} En admettant que l'activité massique suit la loi $a(t) = a_0\,e^{-\lambda t}$, déterminer l'âge $t$ du vestige. Exprimer le résultat en années avec trois chiffres significatifs.
\item \textbf{Datation graphique.} Tracer la courbe $a/a_0 = e^{-\lambda t}$ pour $t$ variant de $0$ à $\SI{20000}{ans}$ (sur papier ou avec un tableur). Déterminer graphiquement l'âge correspondant à $a/a_0 = 0,500$ et comparer avec le résultat de la question 3.
\item \textbf{Activité résiduelle et limite de la méthode.} Calculer l'activité massique de l'échantillon dans $\SI{10000}{ans}$. Sachant que l'appareil du laboratoire ne peut pas mesurer une activité inférieure à $\SI{0,001}{Bq/g}$, estimer l'âge maximal d'un échantillon datable par cette méthode. Conclure sur la limite pratique de la datation au carbone 14.
\item \textbf{Conclusion pour les archéologues.} Rédiger un court texte (10 lignes maximum) indiquant l'âge estimé du foyer, la période historique correspondante, et les précautions d'interprétation.
\item \textbf{Prolongement.} Des vestiges bien plus anciens (plusieurs centaines de milliers d'années) ne peuvent pas être datés au carbone 14. Expliquer pourquoi, et citer une méthode adaptée utilisant un noyau de période beaucoup plus longue.
\end{enumerate}

\courssub{Ressources à mobiliser}

\begin{aretenir}[Boîte à outils de l'activité]
\begin{itemize}
\item Lois de Soddy : conservation du nombre de nucléons $A$ et de la charge $Z$ dans toute équation nucléaire.
\item Désintégration $\beta^-$ : \ce{^{A}_{Z}X -> ^{A}_{Z+1}Y + ^{0}_{-1}e} (un neutron se transforme en proton + électron + antineutrino).
\item Loi de décroissance : $N(t) = N_0\,e^{-\lambda t}$, et de même pour l'activité : $A(t) = A_0\,e^{-\lambda t}$ avec $A = \lambda N$.
\item Période radioactive : $t_{1/2} = \dfrac{\ln 2}{\lambda}$, soit $\lambda = \dfrac{\ln 2}{t_{1/2}}$.
\item Inversion de la loi : $t = \dfrac{1}{\lambda}\ln\!\left(\dfrac{A_0}{A}\right) = \dfrac{t_{1/2}}{\ln 2}\,\ln\!\left(\dfrac{A_0}{A}\right)$.
\item Unité d'activité : le becquerel, $\SI{1}{Bq} = \SI{1}{désintégration/s}$ ; $\SI{1}{an} \approx \SI{3,156e7}{s}$.
\end{itemize}
\end{aretenir}

\begin{attention}[Pièges fréquents]
\begin{itemize}
\item Ne pas confondre $A$ (nombre de masse) et $A(t)$ (activité) : dans cette activité, on note $a$ l'activité massique pour éviter toute ambiguïté.
\item L'activité massique $a$ est proportionnelle au nombre de noyaux $N$ : le rapport $a/a_0$ est égal au rapport $N/N_0$. Inutile de convertir en nombre de noyaux !
\item Dans $\ln(a_0/a)$, le rapport est \emph{sans unité} ; veiller à ce que $a$ et $a_0$ soient exprimés dans la même unité.
\item Le résultat d'une datation est un \emph{âge} (durée écoulée depuis la mort du végétal), pas une date : il faut ensuite la situer par rapport à aujourd'hui.
\end{itemize}
\end{attention}

\courssub{Démarche et corrigé commenté}

\begin{exoresolu}[Datation au carbone 14 des vestiges de Foumban]
Énoncé : un charbon de bois de Foumban présente une activité massique en \ce{^{14}C} de $a = \SI{0,115}{Bq/g}$, contre $a_0 = \SI{0,230}{Bq/g}$ pour un bois vivant. Données : $t_{1/2} = \SI{5730}{ans}$. Déterminer l'âge du vestige, vérifier graphiquement, calculer l'activité dans $\SI{10000}{ans}$ et discuter la limite de la méthode.
\tcblower
\textbf{1. Équation de désintégration $\beta^-$ du carbone 14.}

Le carbone 14 \ce{^{14}_{6}C} se transforme en azote 14. Conservation du nombre de nucléons ($14 = 14 + 0$) et de la charge ($6 = 7 - 1$) :
\[ \ce{^{14}_{6}C -> ^{14}_{7}N + ^{0}_{-1}e} \]
La particule émise est un \cle{électron} (accompagné d'un antineutrino électronique $\bar{\nu}_e$, hors programme). Au sein du noyau, un neutron s'est transformé en proton : $Z$ augmente d'une unité, $A$ reste inchangé.

\textbf{2. Constante radioactive.}

\[ \lambda = \frac{\ln 2}{t_{1/2}} = \frac{0,693}{5730} = \SI{1,21e-4}{an^{-1}} \quad (\text{3 chiffres significatifs}). \]
En secondes : $1\,\text{an} = \SI{3,156e7}{s}$, donc
\[ \lambda = \frac{1,21\times 10^{-4}}{3,156\times 10^{7}} = \SI{3,83e-12}{s^{-1}}. \]
Vérification inverse : $t_{1/2} = \dfrac{\ln 2}{\lambda} = \dfrac{0,693}{1,21\times 10^{-4}} \approx \SI{5730}{ans}$ : la valeur du laboratoire est cohérente.

\textbf{3. Âge du vestige par le calcul.}

L'activité massique suit $a(t) = a_0\,e^{-\lambda t}$, d'où :
\begin{align*}
t &= \frac{1}{\lambda}\ln\!\left(\frac{a_0}{a}\right)
= \frac{t_{1/2}}{\ln 2}\,\ln\!\left(\frac{0,230}{0,115}\right) \\
&= \frac{5730}{0,693}\times \ln 2
= \frac{5730}{0,693}\times 0,693 = \SI{5730}{ans}.
\end{align*}
Le rapport $a/a_0 = 0,500$ vaut exactement $1/2$ : l'échantillon a vécu \emph{exactement une période} de décroissance. L'âge du vestige, exprimé avec trois chiffres significatifs, est $t = \SI{5,73e3}{ans}$ (soit vers 3700 avant notre ère, époque néolithique).

\textbf{4. Vérification graphique.}

On trace $a/a_0 = e^{-\lambda t}$. La courbe passe par les points remarquables $(0\,;\,1)$, $(5730\,;\,0,5)$, $(11460\,;\,0,25)$, $(17190\,;\,0,125)$.

\begin{popfigure}
\begin{center}
\begin{tikzpicture}
\begin{axis}[
width=0.85\linewidth, height=6cm,
xlabel={$t$ (ans)}, ylabel={$a/a_0$},
xmin=0, xmax=20000, ymin=0, ymax=1.05,
xtick={0,5730,11460,17190},
xticklabels={0, 5730, 11460, 17190},
ytick={0,0.125,0.25,0.5,1},
grid=major, grid style={popline!50},
tick label style={font=\small},
]
\addplot[popcurve,domain=0:20000,samples=200]{exp(-0.693*x/5730)};
\addplot[poporange,dashed] coordinates {(5730,0) (5730,0.5) (0,0.5)};
\addplot[mark=*,poporange] coordinates {(5730,0.5)};
\end{axis}
\end{tikzpicture}
\end{center}
\legende{Courbe de décroissance du carbone 14 : lecture graphique de l'âge pour $a/a_0 = 0,500$.}
\end{popfigure}

La lecture graphique à $a/a_0 = 0,500$ redonne $t \approx \SI{5730}{ans}$, en accord avec le calcul. Les deux méthodes se confortent mutuellement : c'est une bonne pratique scientifique.

\textbf{5. Activité résiduelle dans $\SI{10000}{ans}$ et limite de la méthode.}

\[ a(\SI{10000}{ans}) = 0,230\times e^{-1,21\times 10^{-4}\times 10^{4}} = 0,230\times e^{-1,21} = 0,230\times 0,298 \approx \SI{0,0686}{Bq/g}. \]
Limite de détection : on cherche $t_{\max}$ tel que $a = \SI{0,001}{Bq/g}$ :
\[ t_{\max} = \frac{5730}{0,693}\times \ln\!\left(\frac{0,230}{0,001}\right) = 8268\times \ln 230 = 8268\times 5,44 \approx \SI{4,50e4}{ans}. \]
Au-delà d'environ 8 périodes ($\approx \SI{45000}{ans}$), l'activité devient trop faible pour être mesurée avec cet appareil : la datation au carbone 14 atteint sa \cle{limite pratique}.

\textbf{6. Conclusion pour les archéologues (exemple de rédaction).}

\emph{« L'analyse du charbon de bois de la couche 4 indique une activité en carbone 14 égale à la moitié de celle d'un bois vivant. Le foyer a donc fonctionné il y a environ 5 730 ans, au Néolithique ancien. Cette datation suppose que l'échantillon n'a pas été contaminé par du carbone récent (racines, humidité) ; une confirmation sur un second échantillon indépendant est recommandée. »}

\textbf{7. Prolongement.} Pour des vestiges de plusieurs centaines de milliers d'années, le \ce{^{14}C} a quasiment totalement disparu (activité inférieure au bruit de fond). On utilise alors des noyaux à longue période, par exemple la méthode \cle{potassium-argon} : le \ce{^{40}_{19}K} ($t_{1/2} = \SI{1,25e9}{ans}$) se désintègre par capture électronique en \ce{^{40}_{18}Ar} (gaz rare piégé dans la roche volcanique) :
\[ \ce{^{40}_{19}K + ^{0}_{-1}e -> ^{40}_{18}Ar + \nu_e} \]
Le rapport Ar/K mesuré donne l'âge de la solidification de la roche.
\end{exoresolu}

\courssub{Bilan de l'activité}

Cette activité t'a permis de mobiliser l'ensemble des savoir-faire du chapitre dans une situation authentique :
\begin{itemize}
\item les \textbf{lois de Soddy} pour écrire proprement une désintégration $\beta^-$ ;
\item le lien $\lambda = \ln 2 / t_{1/2}$ entre constante radioactive et période, avec un soin particulier apporté aux unités ;
\item l'\textbf{exploitation de la loi de décroissance} sous sa forme inversée $t = \dfrac{t_{1/2}}{\ln 2}\ln\!\left(\dfrac{a_0}{a}\right)$, cœur de toute datation radioactive ;
\item la \textbf{complémentarité calcul / lecture graphique} : une courbe de décroissance exponentielle se lit directement grâce aux points remarquables $t_{1/2}$, $2t_{1/2}$, $3t_{1/2}$ où l'activité est divisée par 2, 4, 8 ;
\item l'esprit critique : toute méthode de mesure a une \textbf{limite} (ici environ $\SI{45000}{ans}$, soit 8 périodes), et le choix du « chronomètre radioactif » dépend de l'ordre de grandeur de l'âge recherché — carbone 14 pour l'archéologie récente, potassium-argon ou uranium-plomb pour la géologie.
\end{itemize}

\begin{savaistu}[Et au Cameroun ?]
Les sites archéologiques camerounais (Foumban, les gravures rupestres du Nord, les fosses à ferraille de la région de Yaoundé) sont datés grâce à ces méthodes, souvent en collaboration avec des laboratoires étrangers faute d'équipement de comptage $\beta$ de très bas bruit sur place. À l'inverse, la radioactivité est bien présente dans les hôpitaux camerounais : le Centre d'oncologie de Yaoundé et l'Hôpital général de Douala utilisent des sources radioactives (cobalt 60, iode 131) pour la radiothérapie et la scintigraphie — les mêmes lois de décroissance y gouvernent le renouvellement des sources et la radioprotection du personnel.
\end{savaistu}

\newpage

\coursec{Méthodes et exercices résolus}

\coursec{Le noyau atomique et sa cohésion}

\begin{definition}[Noyau atomique, nucléides et isotopes]
Un \cle{noyau atomique} est constitué de $Z$ protons (charge $+e$) et $N$ neutrons (charge $0$), appelés collectivement \cle{nucléons}. Le nombre de masse $A = Z + N$ est le nombre total de nucléons. Un \cle{nucléide} est une espèce nucléaire caractérisée par $(Z, A)$, notée $\ce{^{A}_{Z}X}$ où $X$ est le symbole chimique. Des \cle{isotopes} sont des nucléides d'un même élément (même $Z$) mais de nombres de masse $A$ différents (ex. : $\ce{^{12}_{6}C}$, $\ce{^{13}_{6}C}$, $\ce{^{14}_{6}C}$).
\end{definition}

\begin{propriete}[Défaut de masse, énergie de liaison et courbe d'Aston]
La masse d'un noyau $m(\ce{^{A}_{Z}X})$ est \emph{toujours inférieure} à la somme des masses de ses $Z$ protons et $N$ neutrons libres. La différence est le \cle{défaut de masse} :
\[ \Delta m = \left[ Z\,m_p + (A-Z)\,m_n \right] - m(\ce{^{A}_{Z}X}) \quad (>0 \text{ pour un noyau stable}). \]
L'\cle{énergie de liaison} $E_\ell$ est l'énergie nécessaire pour dissocier complètement le noyau en ses nucléons libres (équivalent énergétique du défaut de masse) :
\[ E_\ell = \Delta m \cdot c^2. \]
L'\cle{énergie de liaison par nucléon} $\dfrac{E_\ell}{A}$ caractérise la stabilité nucléaire. La \cle{courbe d'Aston} (représentation de $E_\ell/A$ en fonction de $A$) présente un maximum $\approx 8{,}8\ \mathrm{MeV/nucléon}$ pour $A \approx 56$ (fer 56). Les noyaux légers ($A < 56$) gagnent en stabilité par \cle{fusion} ; les noyaux lourds ($A > 56$) par \cle{fission}.
\end{propriete}

\courssub{Méthode 1 : Identifier un nucléide, des isotopes et calculer l'énergie de liaison}

\begin{methode}[Noyau, isotopes, défaut de masse et énergie de liaison]
\begin{enumerate}
\item \textbf{Lire la notation} $\ce{^{A}_{Z}X}$ : $Z$ = nombre de protons (numéro atomique), $A$ = nombre de nucléons (masse), donc $N = A - Z$ neutrons. Des \cle{isotopes} ont le même $Z$ mais des $A$ différents.
\item \textbf{Calculer le défaut de masse} : $\Delta m = \left[Z\,m_p + (A-Z)\,m_n\right] - m(\ce{^{A}_{Z}X})$. Le défaut de masse est toujours \emph{positif} pour un noyau stable.
\item \textbf{En déduire l'énergie de liaison} : $E_\ell = \Delta m \cdot c^2$. Vérifier l'homogénéité : $[\Delta m \cdot c^2] = \mathrm{kg \cdot m^2 \cdot s^{-2}} = \mathrm{J}$. Pour obtenir des MeV : $1\ \mathrm{u} \cdot c^2 = 931{,}5\ \mathrm{MeV}$.
\item \textbf{Énergie de liaison par nucléon} : $\dfrac{E_\ell}{A}$, en MeV/nucléon. Plus elle est grande, plus le noyau est stable (maximum $\approx 8{,}8$ MeV/nucléon vers le fer 56 sur la courbe d'Aston).
\item \textbf{Rédiger} : donner chaque résultat avec son unité et 3 chiffres significatifs, et conclure par une phrase comparant à la courbe d'Aston.
\end{enumerate}
\end{methode}

\begin{exoresolu}[Cohésion du noyau d'uranium 235]
On donne : $m_p = 1{,}00728\ \mathrm{u}$ ; $m_n = 1{,}00866\ \mathrm{u}$ ; $m(\ce{^{235}_{92}U}) = 235{,}0439\ \mathrm{u}$ ; $1\ \mathrm{u} \cdot c^2 = 931{,}5\ \mathrm{MeV}$.
\begin{enumerate}
\item Donner la composition du noyau $\ce{^{235}_{92}U}$.
\item Calculer son défaut de masse puis son énergie de liaison en MeV.
\item Calculer son énergie de liaison par nucléon. Ce noyau est-il très stable ? Justifier à l'aide de la courbe d'Aston.
\end{enumerate}
\tcblower
\textbf{1. Composition.} Le noyau $\ce{^{235}_{92}U}$ contient $Z = 92$ protons et $N = A - Z = 235 - 92 = 143$ neutrons, soit 235 nucléons au total.

\textbf{2. Défaut de masse et énergie de liaison.} La masse des nucléons séparés vaut :
\begin{align*}
m_{\text{séparés}} &= Z\,m_p + (A-Z)\,m_n = 92 \times 1{,}00728 + 143 \times 1{,}00866 \\
&= 92{,}66976 + 144{,}23838 = 236{,}90814\ \mathrm{u}.
\end{align*}
D'où le défaut de masse :
\[ \Delta m = 236{,}90814 - 235{,}0439 = 1{,}8642\ \mathrm{u}. \]
Le défaut de masse est positif : le noyau est moins massif que ses constituants séparés, la différence s'est transformée en énergie de cohésion. L'énergie de liaison est :
\[ E_\ell = \Delta m \cdot c^2 = 1{,}8642 \times 931{,}5 \approx 1736{,}5\ \mathrm{MeV} \approx 1{,}74 \times 10^{3}\ \mathrm{MeV}. \]

\textbf{3. Énergie de liaison par nucléon.}
\[ \frac{E_\ell}{A} = \frac{1736{,}5}{235} \approx 7{,}39\ \mathrm{MeV/nucléon}. \]
Or le maximum de la courbe d'Aston est d'environ $8{,}8$ MeV/nucléon (région du fer). Comme $7{,}39 < 8{,}8$, le noyau d'uranium 235 est situé \emph{au-delà} du maximum, dans la zone des noyaux lourds : il n'est pas parmi les plus stables et peut gagner en stabilité par \cle{fission} (cassure en deux noyaux plus légers, plus liés). C'est précisément cette propriété qui est exploitée dans les réacteurs nucléaires.

\textbf{Conclusion :} $\Delta m \approx 1{,}864\ \mathrm{u}$, $E_\ell \approx 1{,}74 \times 10^{3}$ MeV et $E_\ell/A \approx 7{,}39$ MeV/nucléon : l'uranium 235 est un noyau lourd, modérément lié, donc fissile.
\end{exoresolu}

\coursec{La radioactivité : définition et propriétés}

\begin{definition}[Radioactivité]
La \cle{radioactivité} est un phénomène nucléaire au cours duquel des noyaux instables (dits \cle{radioactifs}) se transforment spontanément en d'autres noyaux plus stables, en émettant des particules ($\alpha$, $\beta$) et/ou des photons ($\gamma$). Cette transformation est accompagnée d'une libération d'énergie.
\end{definition}

\begin{propriete}[Propriétés fondamentales de la radioactivité]
\begin{enumerate}
\item \textbf{Spontanée :} Elle se produit sans cause extérieure apparente (ni température, ni pression, ni état chimique, ni champ électrique/magnétique ne l'influencent).
\item \textbf{Aléatoire :} Il est impossible de prédire \emph{quand} un noyau donné va se désintégrer. La constante radioactive $\lambda$ représente la \emph{probabilité par unité de temps} qu'un noyau se désintègre.
\item \textbf{Inévitable :} On ne peut ni l'accélérer, ni la ralentir, ni l'arrêter par des moyens physiques ou chimiques usuels.
\end{enumerate}
\end{propriete}

\begin{savaistu}[Histoire et contexte camerounais]
La radioactivité a été découverte en 1896 par Henri Becquerel (sels d'uranium), puis étudiée par Marie et Pierre Curie (polonium, radium). Au Cameroun, les applications médicales sont majeures : le Centre de Radiothérapie de l'Hôpital Général de Yaoundé utilise des sources de cobalt 60 ($\gamma$) pour traiter les cancers ; l'iode 131 ($\beta^-$) traite l'hyperthyroïdie ; le technétium 99m ($\gamma$, $T=6$ h) sert de traceur en scintigraphie. L'Agence Internationale de l'Énergie Atomique (AIEA) accompagne le Cameroun dans la sûreté radiologique et la gestion des déchets.
\end{savaistu}

\coursec{Types de désintégrations et lois de Soddy}

\begin{definition}[Émissions radioactives]
Quatre types d'émissions principaux :
\begin{itemize}
\item \cle{Émission $\alpha$} : noyau d'hélium $\ce{^{4}_{2}He}$ (2 protons, 2 neutrons), charge $+2e$, masse $\approx 4\ \mathrm{u}$, fort pouvoir ionisant, faible portée (quelques cm dans l'air, stoppé par une feuille de papier).
\item \cle{Émission $\beta^-$} : électron $\ce{^{0}_{-1}e}$ (un neutron $\to$ proton + $e^-$ + antineutrino $\bar{\nu}_e$), charge $-e$, masse $m_e$, portée $\sim$ m dans l'air, stoppé par quelques mm d'aluminium.
\item \cle{Émission $\beta^+$} : positon $\ce{^{0}_{+1}e}$ (un proton $\to$ neutron + $e^+$ + neutrino $\nu_e$), charge $+e$, masse $m_e$, s'annihile avec un électron en produisant deux photons $\gamma$ de 511 keV.
\item \cle{Émission $\gamma$} : photon de haute énergie (désexcitation nucléaire), sans charge ni masse au repos, très pénétrant (plomb, béton épais).
\end{itemize}
\end{definition}

\begin{propriete}[Lois de conservation de Soddy]
Dans toute réaction nucléaire, deux grandeurs sont conservées :
\begin{enumerate}
\item Le \textbf{nombre de masse} $A$ (nombre total de nucléons) : $\sum A_{\text{avant}} = \sum A_{\text{après}}$.
\item La \textbf{charge électrique} (nombre atomique $Z$) : $\sum Z_{\text{avant}} = \sum Z_{\text{après}}$.
\end{enumerate}
Ces lois permettent d'identifier le noyau fils.
\end{propriete}

\courssub{Méthode 2 : Écrire et équilibrer une équation de désintégration}

\begin{methode}[Équations nucléaires et lois de Soddy]
\begin{enumerate}
\item \textbf{Identifier le type d'émission} et la particule émise : $\alpha = \ce{^{4}_{2}He}$ ; $\beta^- = \ce{^{0}_{-1}e}$ ; $\beta^+ = \ce{^{0}_{+1}e}$ ; $\gamma$ = photon sans charge ni masse (désexcitation, $Z$ et $A$ inchangés).
\item \textbf{Appliquer les deux lois de conservation de Soddy} : conservation du nombre de nucléons ($\sum A$ identique des deux côtés) et conservation de la charge ($\sum Z$ identique des deux côtés).
\item \textbf{Déterminer le noyau fils} $\ce{^{A'}_{Z'}Y}$ :
\begin{itemize}
\item $\alpha$ : $A' = A - 4$, $Z' = Z - 2$ ;
\item $\beta^-$ : $A' = A$, $Z' = Z + 1$ (un neutron se transforme en proton) ;
\item $\beta^+$ : $A' = A$, $Z' = Z - 1$ (un proton se transforme en neutron).
\end{itemize}
\item \textbf{Identifier l'élément fils} grâce à $Z'$ dans le tableau périodique, puis \textbf{vérifier} en additionnant les $A$ et les $Z$ de part et d'autre de la flèche.
\item Pour le $\beta^+$, ne pas oublier que l'émission s'accompagne d'un positon $\ce{^{0}_{+1}e}$ ; la désintégration $\gamma$ s'écrit $\ce{^{A}_{Z}X^{*} -> ^{A}_{Z}X + \gamma}$ (noyau fils souvent émis dans un état excité).
\end{enumerate}
\end{methode}

\begin{exoresolu}[Famille radioactive du radium et radiothérapie au cobalt 60]
\begin{enumerate}
\item Le radium 226, $\ce{^{226}_{88}Ra}$, est un émetteur $\alpha$. Écrire l'équation de sa désintégration et identifier le noyau fils.
\item Le cobalt 60, $\ce{^{60}_{27}Co}$, utilisé en radiothérapie dans certains hôpitaux, est un émetteur $\beta^-$. Écrire l'équation de sa désintégration.
\item Le noyau fils obtenu en 2. est produit dans un état excité. Écrire l'équation de sa désexcitation.
\end{enumerate}
On donne les numéros atomiques : Rn ($Z=86$), Ra ($Z=88$), Ni ($Z=28$), Co ($Z=27$).
\tcblower
\textbf{1. Désintégration $\alpha$ du radium 226.} La particule émise est $\ce{^{4}_{2}He}$. Les lois de Soddy imposent :
\[ A' = 226 - 4 = 222 \quad ; \quad Z' = 88 - 2 = 86. \]
L'élément de numéro atomique 86 est le radon Rn. D'où l'équation :
\[ \ce{^{226}_{88}Ra -> ^{222}_{86}Rn + ^{4}_{2}He} \]
Vérification : nucléons $226 = 222 + 4$ \checkmark ; charges $88 = 86 + 2$ \checkmark.

\textbf{2. Désintégration $\beta^-$ du cobalt 60.} La particule émise est un électron $\ce{^{0}_{-1}e}$ :
\[ A' = 60 - 0 = 60 \quad ; \quad Z' = 27 - (-1) = 28. \]
L'élément de numéro atomique 28 est le nickel Ni :
\[ \ce{^{60}_{27}Co -> ^{60}_{28}Ni^{*} + ^{0}_{-1}e} \]
Vérification : $60 = 60 + 0$ \checkmark ; $27 = 28 - 1$ \checkmark. (Un neutron du noyau s'est transformé en proton.)

\textbf{3. Désexcitation $\gamma$.} Le noyau de nickel excité redescend vers son état fondamental en émettant un photon $\gamma$, sans changer ni $A$ ni $Z$ :
\[ \ce{^{60}_{28}Ni^{*} -> ^{60}_{28}Ni + \gamma} \]

\textbf{Conclusion :} le radium 226 donne du radon 222 par émission $\alpha$ ; le cobalt 60 donne du nickel 60 par émission $\beta^-$, suivi d'une désexcitation $\gamma$ — ce sont ces photons $\gamma$ très énergétiques qui sont utilisés pour irradier les tumeurs en radiothérapie.
\end{exoresolu}

\coursec{Loi de décroissance radioactive}

\begin{propriete}[Loi de décroissance radioactive, constante $\lambda$ et période $T$]
Pour un échantillon contenant $N_0$ noyaux radioactifs identiques à l'instant $t=0$, le nombre $N(t)$ de noyaux non encore désintégrés à l'instant $t$ suit la loi exponentielle :
\[ N(t) = N_0 e^{-\lambda t} \]
où $\lambda$ est la \cle{constante radioactive} (probabilité de désintégration par unité de temps), en $\mathrm{s^{-1}}$ (ou $\mathrm{h^{-1}}$, $\mathrm{an^{-1}}$ selon l'unité de temps).
La \cle{période radioactive} (ou \cle{demi-vie}) $T$ est le temps au bout duquel la moitié des noyaux initiaux s'est désintégrée : $N(T) = N_0/2$. Le lien fondamental est :
\[ T = \frac{\ln 2}{\lambda} \quad \Longleftrightarrow \quad \lambda = \frac{\ln 2}{T}. \]
\textbf{Attention à la cohérence des unités} : si $T$ est en ans, $\lambda$ est en $\mathrm{an^{-1}}$ et $t$ doit être en ans.
Après $n$ périodes ($t = nT$) : $N = \dfrac{N_0}{2^n}$.
Inverser la loi pour trouver une durée :
\[ t = \frac{1}{\lambda}\ln\!\left(\frac{N_0}{N}\right) = \frac{T}{\ln 2}\ln\!\left(\frac{N_0}{N}\right). \]
\end{propriete}

\courssub{Méthode 3 : Exploiter la loi de décroissance $N(t) = N_0 e^{-\lambda t}$}

\begin{methode}[Décroissance radioactive, constante $\lambda$ et demi-vie $T$]
\begin{enumerate}
\item \textbf{Écrire la loi} $N(t) = N_0 e^{-\lambda t}$ avec $N_0$ nombre de noyaux à $t=0$ et $\lambda$ la constante radioactive (en $\mathrm{s^{-1}}$, $\mathrm{h^{-1}}$ ou $\mathrm{an^{-1}}$ selon l'unité de temps choisie).
\item \textbf{Lien fondamental} avec la demi-vie (période radioactive) $T$ :
\[ T = \frac{\ln 2}{\lambda} \quad \Longleftrightarrow \quad \lambda = \frac{\ln 2}{T}. \]
\textbf{Attention à la cohérence des unités} : si $T$ est en ans, $\lambda$ est en $\mathrm{an^{-1}}$ et $t$ doit être en ans.
\item \textbf{Après $n$ périodes} ($t = nT$) : $N = \dfrac{N_0}{2^n}$. Réflexe : diviser par 2 autant de fois qu'il y a de demi-vies écoulées.
\item \textbf{Inverser la loi} pour trouver une durée : de $\dfrac{N}{N_0} = e^{-\lambda t}$ on tire
\[ t = \frac{1}{\lambda}\ln\!\left(\frac{N_0}{N}\right) = \frac{T}{\ln 2}\ln\!\left(\frac{N_0}{N}\right). \]
\item \textbf{Vérifier la cohérence} : si $t > T$, on doit trouver $N < N_0/2$ ; la décroissance est exponentielle, jamais linéaire.
\end{enumerate}
\end{methode}

\begin{exoresolu}[Déchets d'un centre de médecine nucléaire]
Un laboratoire d'analyses médicales de Yaoundé utilise de l'iode 131, de demi-vie $T = 8{,}0$ jours. Un échantillon contient à $t = 0$ un nombre $N_0 = 4{,}0 \times 10^{16}$ noyaux d'iode 131.
\begin{enumerate}
\item Calculer la constante radioactive $\lambda$ en $\mathrm{jour^{-1}}$ puis en $\mathrm{s^{-1}}$.
\item Calculer le nombre de noyaux restants au bout de 32 jours.
\item Au bout de combien de temps ne reste-t-il que $1{,}0\ \%$ des noyaux initiaux ?
\end{enumerate}
\tcblower
\textbf{1. Constante radioactive.}
\[ \lambda = \frac{\ln 2}{T} = \frac{0{,}693}{8{,}0} \approx 8{,}66 \times 10^{-2}\ \mathrm{jour^{-1}}. \]
Conversion en secondes : $1$ jour $= 86400\ \mathrm{s}$, donc
\[ \lambda = \frac{8{,}66 \times 10^{-2}}{86400} \approx 1{,}00 \times 10^{-6}\ \mathrm{s^{-1}}. \]
Analyse dimensionnelle : $\lambda$ est bien l'inverse d'un temps.

\textbf{2. Noyaux restants après 32 jours.} Remarquons que $32 = 4 \times 8{,}0 = 4T$ : il s'est écoulé exactement 4 demi-vies. Donc :
\[ N = \frac{N_0}{2^4} = \frac{4{,}0 \times 10^{16}}{16} = 2{,}5 \times 10^{15}\ \text{noyaux}. \]
Vérification par la loi exponentielle : $N = 4{,}0 \times 10^{16} \times e^{-0{,}0866 \times 32} = 4{,}0 \times 10^{16} \times e^{-2{,}772} \approx 4{,}0 \times 10^{16} \times 0{,}0625 = 2{,}5 \times 10^{15}$ \checkmark.

\textbf{3. Durée pour qu'il reste $1{,}0\ \%$.} On veut $N = 0{,}010\,N_0$, soit $\dfrac{N_0}{N} = 100$ :
\[ t = \frac{T}{\ln 2}\ln\!\left(\frac{N_0}{N}\right) = \frac{8{,}0}{0{,}693}\times \ln 100 = \frac{8{,}0}{0{,}693}\times 4{,}605 \approx 53{,}2\ \text{jours}. \]
Vérification de cohérence : $53{,}2/8{,}0 \approx 6{,}6$ demi-vies, et $1/2^{6{,}6} \approx 0{,}0103 \approx 1\ \%$ \checkmark.

\textbf{Conclusion :} $\lambda \approx 8{,}7 \times 10^{-2}\ \mathrm{jour^{-1}} \approx 1{,}0 \times 10^{-6}\ \mathrm{s^{-1}}$ ; après 32 jours il reste $2{,}5 \times 10^{15}$ noyaux ; l'activité de la source tombe à $1\ \%$ de sa valeur initiale au bout d'environ 53 jours, ce qui fixe la durée de stockage de sécurité des déchets.
\end{exoresolu}

\coursec{Activité d'un échantillon radioactif}

\begin{definition}[Activité, Becquerel]
L'\cle{activité} $A$ d'un échantillon radioactif est le nombre moyen de désintégrations qui se produisent par seconde dans cet échantillon :
\[ A = \lambda N \qquad \text{(unité : le \cle{becquerel}, } 1~\mathrm{Bq} = 1~\mathrm{s^{-1}}\text{)} \]
où $N$ est le nombre de noyaux radioactifs présents à l'instant considéré. L'ancienne unité, le curie (Ci), vaut $1\ \mathrm{Ci} = 3{,}7 \times 10^{10}\ \mathrm{Bq}$.
\end{definition}

\begin{propriete}[Loi de décroissance de l'activité et lien avec la masse]
Puisque $N(t) = N_0 e^{-\lambda t}$, l'activité suit la même loi exponentielle :
\[ A(t) = A_0 e^{-\lambda t} \quad \text{avec} \quad A_0 = \lambda N_0. \]
Le nombre de noyaux dans une masse $m$ d'échantillon de masse molaire $M$ est $N = \dfrac{m}{M}\,\mathcal{N}_A$ ($\mathcal{N}_A = 6{,}02 \times 10^{23}\ \mathrm{mol^{-1}}$). D'où :
\[ A = \lambda \frac{m}{M}\mathcal{N}_A. \]
\end{propriete}

\courssub{Méthode 4 : Calculer et exploiter l'activité d'un échantillon}

\begin{methode}[Activité $A = \lambda N$]
\begin{enumerate}
\item \textbf{Définition} : l'activité $A$ est le nombre moyen de désintégrations par seconde : $A = \lambda N$, unité le \cle{becquerel} ($1\ \mathrm{Bq} = 1$ désintégration/s).
\item \textbf{Décroissance de l'activité} : comme $N = N_0 e^{-\lambda t}$, on a aussi $A(t) = A_0 e^{-\lambda t}$ avec $A_0 = \lambda N_0$. L'activité suit la même loi exponentielle que $N$.
\item \textbf{Lien avec la masse} : le nombre de noyaux dans une masse $m$ d'échantillon est $N = \dfrac{m}{M}\,\mathcal{N}_A$ ($M$ masse molaire, $\mathcal{N}_A = 6{,}02 \times 10^{23}\ \mathrm{mol^{-1}}$). D'où $A = \lambda \dfrac{m}{M}\mathcal{N}_A$.
\item \textbf{Convertir $\lambda$ en $\mathrm{s^{-1}}$} avant tout calcul d'activité en becquerels : c'est l'erreur la plus fréquente.
\item \textbf{Exploiter une mesure} : connaissant $A$ et $\lambda$, on remonte à $N = A/\lambda$ puis à la masse $m = \dfrac{N\,M}{\mathcal{N}_A}$.
\end{enumerate}
\end{methode}

\begin{exoresolu}[Activité d'une source de césium 137]
Le césium 137 est un émetteur $\beta^-$ de demi-vie $T = 30$ ans, de masse molaire $M = 137\ \mathrm{g \cdot mol^{-1}}$. On dispose d'un échantillon de masse $m = 1{,}0\ \mathrm{mg}$.
\begin{enumerate}
\item Calculer le nombre de noyaux $N_0$ présents dans l'échantillon.
\item Calculer l'activité initiale $A_0$ de cet échantillon en becquerels.
\item Quelle sera son activité au bout de 90 ans ? Commenter.
\end{enumerate}
On donne : $\mathcal{N}_A = 6{,}02 \times 10^{23}\ \mathrm{mol^{-1}}$ ; $\ln 2 = 0{,}693$ ; 1 an $= 3{,}156 \times 10^{7}\ \mathrm{s}$.
\tcblower
\textbf{1. Nombre de noyaux.}
\[ N_0 = \frac{m}{M}\,\mathcal{N}_A = \frac{1{,}0 \times 10^{-3}}{137}\times 6{,}02 \times 10^{23} \approx 4{,}39 \times 10^{18}\ \text{noyaux}. \]

\textbf{2. Activité initiale.} Il faut d'abord exprimer $\lambda$ en $\mathrm{s^{-1}}$ :
\[ \lambda = \frac{\ln 2}{T} = \frac{0{,}693}{30 \times 3{,}156 \times 10^{7}} = \frac{0{,}693}{9{,}468 \times 10^{8}} \approx 7{,}32 \times 10^{-10}\ \mathrm{s^{-1}}. \]
Alors :
\[ A_0 = \lambda N_0 = 7{,}32 \times 10^{-10} \times 4{,}39 \times 10^{18} \approx 3{,}2 \times 10^{9}\ \mathrm{Bq}. \]
Soit environ 3,2 milliards de désintégrations par seconde pour un seul milligramme : l'unité becquerel correspond à des nombres considérables.

\textbf{3. Activité au bout de 90 ans.} $90 = 3 \times 30 = 3T$ : trois demi-vies se sont écoulées, donc
\[ A = \frac{A_0}{2^3} = \frac{3{,}2 \times 10^{9}}{8} = 4{,}0 \times 10^{8}\ \mathrm{Bq}. \]

\textbf{Conclusion :} $N_0 \approx 4{,}4 \times 10^{18}$ noyaux, $A_0 \approx 3{,}2 \times 10^{9}$ Bq ; après 90 ans, l'activité vaut encore $4{,}0 \times 10^{8}$ Bq, soit un huitième de l'activité initiale : les déchets à base de césium 137 restent dangereux pendant plusieurs générations, d'où l'importance du stockage de longue durée.
\end{exoresolu}

\coursec{Bilan énergétique et applications}

\courssub{Méthode 5 : Établir le bilan énergétique d'une désintégration}

\begin{methode}[Énergie libérée par une réaction nucléaire]
\begin{enumerate}
\item \textbf{Écrire l'équation} de la désintégration en appliquant les lois de Soddy (méthode 2).
\item \textbf{Calculer la variation de masse} : $\Delta m = m_{\text{produits}} - m_{\text{père}} = \left[m(\text{fils}) + m(\text{particule})\right] - m(\ce{^{A}_{Z}X})$. Pour une désintégration \emph{spontanée}, on doit trouver $\Delta m < 0$.
\item \textbf{Énergie libérée} : $E = |\Delta m| \cdot c^2$. En pratique : $E\ (\mathrm{MeV}) = |\Delta m|\ (\mathrm{u}) \times 931{,}5\ \mathrm{MeV/u}$. Conversion : $1\ \mathrm{MeV} = 1{,}60 \times 10^{-13}\ \mathrm{J}$.
\item \textbf{Interpréter} : cette énergie est emportée sous forme d'énergie cinétique des produits (particule $\alpha$ ou $\beta$, noyau fils) et, s'il y a lieu, d'un photon $\gamma$ : $E = E_c(\text{fils}) + E_c(\text{particule}) + E_\gamma$.
\item \textbf{Condition de spontanéité} : la désintégration n'est possible que si $E > 0$, c'est-à-dire si le noyau père est plus massif que l'ensemble des produits.
\end{enumerate}
\end{methode}

\begin{exoresolu}[Bilan énergétique de la désintégration $\alpha$ du polonium 210]
Le polonium 210 se désintègre selon : $\ce{^{210}_{84}Po -> ^{206}_{82}Pb + ^{4}_{2}He}$.
On donne : $m(\ce{^{210}_{84}Po}) = 209{,}9829\ \mathrm{u}$ ; $m(\ce{^{206}_{82}Pb}) = 205{,}9745\ \mathrm{u}$ ; $m(\ce{^{4}_{2}He}) = 4{,}0026\ \mathrm{u}$ ; $1\ \mathrm{u} \cdot c^2 = 931{,}5\ \mathrm{MeV}$ ; $1\ \mathrm{MeV} = 1{,}60 \times 10^{-13}\ \mathrm{J}$.
\begin{enumerate}
\item Montrer que cette désintégration est énergétiquement possible.
\item Calculer en MeV puis en joules l'énergie libérée par la désintégration d'un noyau.
\item Sous quelles formes cette énergie est-elle emportée ?
\end{enumerate}
\tcblower
\textbf{1. Condition de spontanéité.} Calculons la variation de masse :
\begin{align*}
\Delta m &= \left[m(\ce{^{206}_{82}Pb}) + m(\ce{^{4}_{2}He})\right] - m(\ce{^{210}_{84}Po}) \\
&= (205{,}9745 + 4{,}0026) - 209{,}9829 \\
&= 209{,}9771 - 209{,}9829 = -5{,}8 \times 10^{-3}\ \mathrm{u}.
\end{align*}
Comme $\Delta m < 0$, la masse des produits est inférieure à celle du noyau père : la désintégration libère de l'énergie, elle est donc énergétiquement possible (spontanée).

\textbf{2. Énergie libérée.}
\[ E = |\Delta m|\cdot c^2 = 5{,}8 \times 10^{-3} \times 931{,}5 \approx 5{,}4\ \mathrm{MeV}. \]
Conversion en joules :
\[ E = 5{,}4 \times 1{,}60 \times 10^{-13} \approx 8{,}6 \times 10^{-13}\ \mathrm{J}. \]
Cette énergie paraît minuscule, mais rapportée à une mole de noyaux elle devient considérable : $E_{\text{mol}} = E \times \mathcal{N}_A \approx 8{,}6 \times 10^{-13} \times 6{,}02 \times 10^{23} \approx 5{,}2 \times 10^{11}\ \mathrm{J/mol}$, des millions de fois plus qu'une réaction chimique.

\textbf{3. Formes d'énergie.} L'énergie libérée se répartit en énergie cinétique de la particule $\alpha$ (l'essentiel, car elle est beaucoup plus légère que le noyau fils), en énergie cinétique de recul du noyau de plomb, et éventuellement en photon $\gamma$ si le noyau fils est produit dans un état excité :
\[ E = E_c(\alpha) + E_c(\ce{Pb}) + E_\gamma. \]

\textbf{Conclusion :} $\Delta m = -5{,}8 \times 10^{-3}$ u $< 0$ : la désintégration est spontanée et libère $E \approx 5{,}4$ MeV $\approx 8{,}6 \times 10^{-13}$ J par noyau, emportée principalement sous forme d'énergie cinétique de la particule $\alpha$.
\end{exoresolu}

\courssub{Méthode 6 : Réaliser une datation au carbone 14}

\begin{methode}[Datation d'un échantillon]
\begin{enumerate}
\item \textbf{Principe} : tant que l'organisme est vivant, il échange du carbone avec l'atmosphère et son taux de $\ce{^{14}_{6}C}$ est constant ($A_0$ connu). À la mort, les échanges cessent et le carbone 14 décroît avec $T = 5730$ ans.
\item \textbf{Mesurer} l'activité $A(t)$ de l'échantillon (ou le rapport $N/N_0$) et la comparer à l'activité $A_0$ d'un échantillon vivant de même nature et même masse.
\item \textbf{Appliquer la loi inversée} :
\[ t = \frac{T}{\ln 2}\,\ln\!\left(\frac{A_0}{A(t)}\right). \]
\item \textbf{Vérifier l'ordre de grandeur} : si $A = A_0/2$, on doit trouver $t = T$ ; si $A = A_0/4$, $t = 2T$, etc.
\item \textbf{Limites de la méthode} : au-delà d'environ $8\,T \approx 45\,000$ ans, l'activité résiduelle est trop faible pour être mesurée correctement ; on utilise alors d'autres couples radioactifs (potassium-argon, uranium-plomb).
\end{enumerate}
\end{methode}

\begin{exoresolu}[Datation d'un objet en bois]
Des archéologues découvrent un fragment de bois dans un site ancien. L'activité en carbone 14 de ce fragment est mesurée à $A = 0{,}115\ \mathrm{Bq}$ par gramme de carbone, alors qu'un bois vivant de même essence présente une activité de $A_0 = 0{,}230\ \mathrm{Bq}$ par gramme. La demi-vie du carbone 14 est $T = 5730$ ans.
\begin{enumerate}
\item Expliquer pourquoi l'activité du fragment est inférieure à celle du bois vivant.
\item Déterminer l'âge du fragment.
\end{enumerate}
\tcblower
\textbf{1. Explication.} Dans le bois vivant, les échanges de $\mathrm{CO_2}$ avec l'atmosphère maintiennent constante la proportion de carbone 14. Une fois l'arbre abattu, ces échanges cessent : le carbone 14 n'est plus renouvelé et son nombre de noyaux décroît selon $N = N_0 e^{-\lambda t}$, donc l'activité $A = \lambda N$ décroît de la même façon. C'est pourquoi $A < A_0$.

\textbf{2. Âge du fragment.} On remarque que $\dfrac{A}{A_0} = \dfrac{0{,}115}{0{,}230} = 0{,}500 = \dfrac{1}{2}$ : l'activité a été divisée par 2, il s'est donc écoulé exactement \emph{une} demi-vie. Vérification par la formule :
\[ t = \frac{T}{\ln 2}\ln\!\left(\frac{A_0}{A}\right) = \frac{5730}{0{,}693}\times \ln 2 = 5730\ \text{ans}. \]

\textbf{Conclusion :} le fragment de bois est âgé d'environ $5{,}7 \times 10^{3}$ ans, soit environ 5700 ans ; l'objet date donc d'environ 3700 ans avant notre ère. La méthode est fiable ici car l'âge est de l'ordre de la demi-vie, donc l'activité résiduelle reste aisément mesurable.
\end{exoresolu}

\begin{savaistu}[Applications médicales, industrielles et limites]
\begin{itemize}
\item \textbf{Médecine nucléaire (Cameroun) :} Radiothérapie externe (Co-60, accélérateurs linéaires) à l'Hôpital Général de Yaoundé et à Douala ; Curiethérapie (sources scellées Ir-192, I-125) ; Scintigraphie (Tc-99m, $T=6$ h) pour imagerie osseuse, cardiaque, rénale ; TEP (Tomographie par Émission de Positons) avec F-18 ($T=110$ min) pour détection précoce des cancers.
\item \textbf{Traceurs industriels :} Détection de fuites sur les pipelines (SNPC, SONARA) par injection de radioisotopes ($\beta$ ou $\gamma$) ; jauges de niveau/épaisseur dans les usines (Cimencam, Alucam) ; radiographie industrielle (Ir-192, Co-60) pour contrôle de soudures.
\item \textbf{Énergie :} Fission de l'U-235 (centrales nucléaires, non encore au Cameroun où l'hydroélectricité domine avec les barrages d'\textbf{Edéa}, \textbf{Song Loulou}, \textbf{Lom Pangar}, \textbf{Nachtigal}) ; Fusion (projet ITER, énergie du Soleil).
\item \textbf{Protection radiologique :} Principes ALARA (As Low As Reasonably Achievable) : justification, optimisation, limitation des doses. Unités : Gray (Gy = J/kg) pour la dose absorbée, Sievert (Sv) pour la dose équivalente.
\end{itemize}
\end{savaistu}

\courssub{Les pièges à éviter}

\begin{attention}[Les 5 erreurs qui coûtent des points au Bac]
\begin{enumerate}
\item \textbf{Confondre $A$ et $Z$ dans une équation nucléaire.} En $\beta^-$, le nombre de masse $A$ ne change \emph{pas} et $Z$ \emph{augmente} de 1 ; en $\beta^+$, $Z$ \emph{diminue} de 1. Toujours vérifier explicitement la conservation de $A$ et de $Z$ en écrivant les deux bilans sous l'équation : un correcteur les cherche.
\item \textbf{Oublier de convertir $\lambda$ en $\mathrm{s^{-1}}$ pour calculer une activité en becquerels.} Si $T$ est donné en jours ou en ans, $\lambda = \ln 2 / T$ est en $\mathrm{jour^{-1}}$ ou $\mathrm{an^{-1}}$ : il faut convertir avant d'écrire $A = \lambda N$ en Bq. Une activité calculée sans conversion est fausse d'un facteur 86400 ou plus.
\item \textbf{Croire que la décroissance est linéaire.} Après 2 demi-vies, il ne reste pas « la moitié de la moitié... donc zéro » mais $N_0/4$ ; après $n$ périodes, $N = N_0/2^n$. De même, écrire $t = \dfrac{1}{\lambda}\ln\!\left(\dfrac{N_0}{N}\right)$ et non $t = \dfrac{N_0 - N}{\lambda}$.
\item \textbf{Se tromper de signe dans le bilan énergétique.} L'énergie libérée est $E = |\Delta m|c^2$ avec $\Delta m = m_{\text{produits}} - m_{\text{père}} < 0$ pour une désintégration spontanée. Ne jamais écrire $E = \Delta m c^2$ avec $\Delta m$ négatif en le présentant comme une énergie libérée positive : préciser le signe et prendre la valeur absolue.
\item \textbf{Négliger les unités et les chiffres significatifs dans la conclusion.} Une énergie sans unité (J ou MeV ?), une demi-vie en « 8 » sans « jours », ou un résultat donné avec 8 chiffres significatifs sont pénalisés. Terminer chaque question par une phrase de conclusion avec valeur, unité et 2 ou 3 chiffres significatifs.
\end{enumerate}
\end{attention}

\newpage

\coursec{Exercices}

\courssub{Je vérifie mes connaissances}

\begin{exercice}[QCM : premiers réflexes]
Pour chaque question, une seule réponse est exacte. Entoure-la.
\begin{enumerate}
\item La radioactivité est un phénomène :
\begin{enumerate}
\item[a)] provoqué par une réaction chimique du noyau avec l'air ;
\item[b)] spontané, aléatoire et inévitable ;
\item[c)] qui dépend de la température de l'échantillon.
\end{enumerate}
\item Dans une désintégration $\alpha$, le noyau fils a :
\begin{enumerate}
\item[a)] $A$ diminué de 4 et $Z$ diminué de 2 ;
\item[b)] $A$ diminué de 2 et $Z$ diminué de 4 ;
\item[c)] $A$ inchangé et $Z$ augmenté de 1.
\end{enumerate}
\item La constante radioactive $\lambda$ s'exprime en :
\begin{enumerate}
\item[a)] $\mathrm{s}$ ;
\item[b)] $\mathrm{s^{-1}}$ ;
\item[c)] $\mathrm{Bq \cdot s}$.
\end{enumerate}
\item Deux isotopes d'un même élément possèdent :
\begin{enumerate}
\item[a)] le même nombre de neutrons ;
\item[b)] le même nombre de protons mais des nombres de neutrons différents ;
\item[c)] le même nombre de masse.
\end{enumerate}
\item Au bout de deux périodes radioactives, le nombre de noyaux restants est :
\begin{enumerate}
\item[a)] $N_0/2$ ;
\item[b)] $N_0/4$ ;
\item[c)] nul.
\end{enumerate}
\end{enumerate}
\end{exercice}

\begin{exercice}[Vrai ou faux ? Justifie]
Réponds par vrai ou faux en justifiant chaque réponse par une phrase ou un calcul court.
\begin{enumerate}
\item L'activité d'un échantillon radioactif est constante au cours du temps.
\item La période radioactive $t_{1/2}$ et la constante $\lambda$ sont liées par $t_{1/2} = \dfrac{\ln 2}{\lambda}$.
\item L'émission $\gamma$ modifie le nombre de masse du noyau.
\item Sur la courbe d'Aston, les noyaux les plus stables sont ceux dont l'énergie de liaison par nucléon est la plus grande en valeur absolue.
\item On peut prévoir l'instant exact de désintégration d'un noyau donné.
\end{enumerate}
\end{exercice}

\begin{exercice}[Définitions et grandeurs]
\begin{enumerate}
\item Définis : nucléide, défaut de masse, énergie de liaison, période radioactive, activité.
\item Rappelle l'unité SI de l'activité et sa définition.
\item Écris la relation entre l'énergie de liaison $E_\ell$, le défaut de masse $\Delta m$ et la célérité $c$ de la lumière, puis vérifie l'homogénéité de cette relation par analyse dimensionnelle.
\end{enumerate}
\end{exercice}

\begin{exercice}[Calculs directs]
On donne : $\ln 2 \approx 0,693$ ; $1\ \mathrm{u} = 1,66\times 10^{-27}\ \mathrm{kg}$ ; $c = 3,00\times 10^{8}\ \mathrm{m\cdot s^{-1}}$ ; $1\ \mathrm{eV} = 1,60\times 10^{-19}\ \mathrm{J}$ ; 1 an $= 3,156\times 10^{7}\ \mathrm{s}$.
\begin{enumerate}
\item La période radioactive du cobalt 60 est $t_{1/2} = 5,3$ ans. Calcule sa constante radioactive $\lambda$ en $\mathrm{s^{-1}}$.
\item Un échantillon contient $N = 2,0\times 10^{15}$ noyaux d'un radionucléide de constante $\lambda = 4,0\times 10^{-9}\ \mathrm{s^{-1}}$. Calcule son activité.
\item Le défaut de masse du noyau d'hélium 4 vaut $\Delta m = 0,0304\ \mathrm{u}$. Calcule son énergie de liaison en joules, puis en MeV, puis son énergie de liaison par nucléon.
\end{enumerate}
\end{exercice}

\courssub{Je m'entraîne}

\begin{exercice}[Équations de désintégration et lois de Soddy]
Complète chaque équation en appliquant les lois de conservation de Soddy (conservation du nombre de masse $A$ et du nombre de charge $Z$), et identifie le type de radioactivité ($\alpha$, $\beta^-$, $\beta^+$ ou $\gamma$).
\begin{enumerate}
\item $\ce{^{210}_{84}Po -> ^{206}_{82}Pb} + \ldots$
\item $\ce{^{14}_{6}C -> $\cdots$ + ^{0}_{-1}e}$
\item $\ce{^{22}_{11}Na -> ^{22}_{10}Ne} + \ldots$
\item $\ce{^{137}_{55}Cs -> $\cdots$ + ^{0}_{-1}e}$ suivi de $\ce{^{137m}_{56}Ba -> ^{137}_{56}Ba} + \ldots$
\item $\ce{^{226}_{88}Ra -> $\cdots$ + ^{4}_{2}He}$
\end{enumerate}
\end{exercice}

\begin{exercice}[Énergie libérée par la désintégration $\alpha$ de l'uranium 238]
L'uranium 238 se désintègre selon : $\ce{^{238}_{92}U -> ^{234}_{90}Th + ^{4}_{2}He}$.
On donne les masses des noyaux : $m(\ce{^{238}U}) = 238,0508\ \mathrm{u}$ ; $m(\ce{^{234}Th}) = 234,0436\ \mathrm{u}$ ; $m(\ce{^{4}He}) = 4,0026\ \mathrm{u}$ ; $1\ \mathrm{u} = 931,5\ \mathrm{MeV}/c^{2}$ ; $1\ \mathrm{eV} = 1,60\times 10^{-19}\ \mathrm{J}$.
\begin{enumerate}
\item Calcule la perte de masse $\Delta m$ accompagnant la désintégration d'un noyau, en unité de masse atomique.
\item En déduis l'énergie libérée $E$ en MeV, puis en joules.
\item Cette énergie est-elle absorbée ou libérée ? Justifie le signe du bilan.
\end{enumerate}
\end{exercice}

\begin{exercice}[Rejet d'un déchet hospitalier à l'iode 131]
Dans un hôpital de Yaoundé, on utilise l'iode 131 ($t_{1/2} = 8,0$ jours) pour le traitement de la thyroïde. Un flacon usagé présente une activité $A_0 = 4,0\times 10^{8}\ \mathrm{Bq}$. La réglementation autorise son rejet avec les déchets ordinaires lorsque son activité devient inférieure à $A_{\mathrm{seuil}} = 1,0\times 10^{7}\ \mathrm{Bq}$.
\begin{enumerate}
\item Calcule la constante radioactive $\lambda$ de l'iode 131 en $\mathrm{jour^{-1}}$.
\item Écris la loi d'évolution de l'activité $A(t)$ du flacon.
\item Détermine la durée minimale de stockage avant rejet. Exprime le résultat en jours, puis en nombre de périodes.
\end{enumerate}
\end{exercice}

\begin{exercice}[Courbe d'Aston et stabilité des noyaux]
On donne les énergies de liaison par nucléon : $E_\ell/A(\ce{^{2}H}) = 1,1\ \mathrm{MeV}$ ; $E_\ell/A(\ce{^{56}Fe}) = 8,8\ \mathrm{MeV}$ ; $E_\ell/A(\ce{^{235}U}) = 7,6\ \mathrm{MeV}$.
\begin{enumerate}
\item Rappelle ce que représente l'énergie de liaison par nucléon et ce qu'indique sa valeur sur la stabilité d'un noyau.
\item Classe les trois noyaux par ordre de stabilité croissante en justifiant.
\item Explique, à l'aide de la courbe d'Aston, pourquoi la fission de $\ce{^{235}U}$ et la fusion de deux noyaux de $\ce{^{2}H}$ libèrent toutes deux de l'énergie.
\item Situe approximativement ces trois noyaux sur une courbe d'Aston que tu esquisseras (allure, axes, valeurs remarquables).
\end{enumerate}
\end{exercice}

\courssub{Je me prépare au Bac}

\begin{exercice}[Type Bac : le radium 226 de Marie Curie]
Le radium 226, découvert par Pierre et Marie Curie, est un émetteur $\alpha$ de période radioactive $t_{1/2} = 1600$ ans. Il se désintègre en radon 222 ($\ce{^{222}_{86}Rn}$), gaz radioactif qui s'accumule parfois dans les habitations.
On donne : $M(\ce{^{226}Ra}) = 226\ \mathrm{g\cdot mol^{-1}}$ ; $\mathcal{N}_A = 6,02\times 10^{23}\ \mathrm{mol^{-1}}$ ; $\ln 2 \approx 0,693$ ; 1 an $= 3,156\times 10^{7}\ \mathrm{s}$.
\begin{enumerate}
\item Définis la radioactivité $\alpha$. Écris l'équation de désintégration du radium 226 en précisant les lois de conservation utilisées.
\item Calcule la constante radioactive $\lambda$ du radium 226 en $\mathrm{s^{-1}}$.
\item Calcule le nombre $N_0$ de noyaux contenus dans un échantillon de masse $m = 1,00\ \mathrm{g}$ de radium 226 pur.
\item En déduis l'activité initiale $A_0$ de cet échantillon. Compare à la définition historique du curie ($1\ \mathrm{Ci} = 3,7\times 10^{10}\ \mathrm{Bq}$).
\item Calcule l'activité de cet échantillon au bout de $t = 4800$ ans. Exprime d'abord le rapport $t/t_{1/2}$.
\item Au bout de combien de temps l'activité de l'échantillon sera-t-elle divisée par 1000 ?
\end{enumerate}
\end{exercice}

\begin{exercice}[Type Bac : datation au carbone 14 d'un ossement]
Le carbone 14 ($t_{1/2} = 5730$ ans) est produit dans la haute atmosphère et s'incorpore aux organismes vivants, dont l'activité en carbone 14 reste constante tant qu'ils sont en vie : $A_{\mathrm{vivant}} = 13,6$ désintégrations par minute et par gramme de carbone. À la mort de l'organisme, cette activité décroît.
Des archéologues découvrent lors de fouilles un ossement dont l'activité mesurée est $12,5\ \%$ de celle d'un organisme vivant actuel.
On donne : $\ln 2 \approx 0,693$.
\begin{enumerate}
\item Écris l'équation de désintégration $\beta^-$ du carbone 14 ($\ce{^{14}_{6}C}$) en azote 14 ($\ce{^{14}_{7}N}$).
\item Établis l'expression de l'âge $t$ de l'échantillon en fonction de $t_{1/2}$, $A_0$ et $A(t)$.
\item Montre que $12,5\ \%$ correspond à un nombre entier de périodes, puis déduis l'âge de l'ossement.
\item Un second ossement du même site présente une activité de $3,4$ désintégrations par minute et par gramme de carbone. Calcule son âge.
\item La méthode du carbone 14 est réputée fiable jusqu'à environ 10 périodes. Calcule l'activité résiduelle correspondante (en $\%$ de $A_0$) et explique pourquoi la datation devient incertaine au-delà.
\end{enumerate}
\end{exercice}

\begin{exercice}[Type Bac : détermination expérimentale d'une période radioactive]
Au laboratoire, on mesure l'activité $A$ d'un échantillon contenant un seul radionucléide, à différentes dates :

\begin{center}
\begin{tabular}{|c|cccccc|}
\hline
$t$ (jours) & 0 & 5 & 10 & 15 & 20 & 25 \\
\hline
$A$ ($10^{6}$ Bq) & 8,00 & 5,04 & 3,17 & 2,00 & 1,26 & 0,79 \\
\hline
\end{tabular}
\end{center}

On donne : $\ln 8{,}00 = 2,079$ ; $\ln 5{,}04 = 1,617$ ; $\ln 3{,}17 = 1,154$ ; $\ln 2{,}00 = 0,693$ ; $\ln 1{,}26 = 0,231$ ; $\ln 0{,}79 = -0,236$.
Table d'identification (périodes) : phosphore 32 : 14,3 jours ; iode 131 : 8,0 jours ; or 198 : 2,7 jours ; sodium 24 : 15,0 heures.
\begin{enumerate}
\item Rappelle la loi de décroissance de l'activité et montre que $\ln A$ est une fonction affine du temps $t$ dont on précisera la pente.
\item Recopie et complète le tableau avec les valeurs de $\ln A$ (avec $A$ en $10^{6}$ Bq).
\item Trace la courbe $\ln A = f(t)$ sur papier millimétré (échelles : 1 cm pour 2,5 jours en abscisse ; 1 cm pour 0,2 en ordonnée). Vérifie que les points sont alignés.
\item Détermine graphiquement la pente de la droite, puis la constante radioactive $\lambda$ en $\mathrm{jour^{-1}}$ et en $\mathrm{s^{-1}}$.
\item Calcule la période radioactive $t_{1/2}$ et identifie le nucléide à l'aide de la table.
\item Détermine graphiquement l'activité de l'échantillon à $t = 30$ jours et vérifie par le calcul.
\end{enumerate}
\end{exercice}



\newpage

\coursec{Corrigés des exercices}

\coursec{La radioactivité et la loi de décroissance}

\courssub{Je vérifie mes connaissances}

\begin{corrige}[Exercice 1 — QCM : premiers réflexes]
\begin{enumerate}
\item \textbf{b)} La radioactivité est un phénomène \cle{spontané} (sans cause extérieure), \cle{aléatoire} (on ne peut prévoir l'instant de désintégration d'un noyau) et \cle{inévitable} (rien ne peut l'empêcher). Elle ne dépend ni de la température, ni de la pression, ni de l'état chimique : c'est un phénomène \emph{nucléaire}, pas chimique.
\item \textbf{a)} L'émission d'une particule $\alpha$ $\ce{^{4}_{2}He}$ fait perdre 4 nucléons dont 2 protons : $A \to A-4$ et $Z \to Z-2$.
\item \textbf{b)} D'après $N = N_0\,e^{-\lambda t}$, le produit $\lambda t$ doit être sans dimension : $\lambda$ s'exprime en $\mathrm{s^{-1}}$.
\item \textbf{b)} Des isotopes ont le même $Z$ (même élément chimique) mais des nombres de neutrons $N = A - Z$ différents.
\item \textbf{b)} À chaque période, $N$ est divisé par 2 : après deux périodes, $N = N_0/4$. La décroissance exponentielle n'atteint jamais zéro.
\end{enumerate}
\end{corrige}

\begin{corrige}[Exercice 2 — Vrai ou faux ? Justifie]
\begin{enumerate}
\item \textbf{Faux.} L'activité $A = \lambda N$ décroît exponentiellement comme $N$ : $A(t) = A_0\,e^{-\lambda t}$. Elle diminue car le nombre de noyaux encore radioactifs diminue.
\item \textbf{Vrai.} Par définition $N(t_{1/2}) = N_0/2$, soit $e^{-\lambda t_{1/2}} = \dfrac{1}{2}$, d'où $\lambda\, t_{1/2} = \ln 2$ et $t_{1/2} = \dfrac{\ln 2}{\lambda}$.
\item \textbf{Faux.} L'émission $\gamma$ est la désexcitation d'un noyau fils formé en état excité : le photon $\gamma$ n'emporte ni charge ni masse ; $A$ et $Z$ sont inchangés.
\item \textbf{Vrai.} Plus $E_\ell/A$ est grande (en valeur absolue, car on la représente souvent en $-E_\ell/A$), plus il faut fournir d'énergie par nucléon pour dissocier le noyau : il est donc plus stable (zone du fer, $E_\ell/A \approx \num{8,8}\ \mathrm{MeV}$).
\item \textbf{Faux.} La désintégration d'un noyau individuel est un phénomène \cle{aléatoire} : seule la loi statistique $N = N_0 e^{-\lambda t}$ décrit le comportement d'une population de noyaux très grande.
\end{enumerate}
\end{corrige}

\begin{corrige}[Exercice 3 — Définitions et grandeurs]
\begin{enumerate}
\item \begin{itemize}
\item \cle{Nucléide} : ensemble des noyaux (ou atomes) caractérisés par un couple $(Z, A)$ donné, noté $\ce{^{A}_{Z}X}$.
\item \cle{Défaut de masse} : différence entre la somme des masses des nucléons séparés au repos et la masse du noyau : $\Delta m = Z\,m_p + (A-Z)\,m_n - m_{\text{noyau}} > 0$.
\item \cle{Énergie de liaison} $E_\ell$ : énergie qu'il faut fournir au noyau au repos pour le dissocier en ses nucléons séparés au repos.
\item \cle{Période radioactive} $t_{1/2}$ : durée au bout de laquelle la moitié des noyaux radioactifs initialement présents s'est désintégrée.
\item \cle{Activité} $A$ : nombre moyen de désintégrations par seconde dans l'échantillon : $A = \lambda N$.
\end{itemize}
\item L'unité SI de l'activité est le \cle{becquerel} (Bq) : $\SI{1}{Bq} = 1$ désintégration par seconde.
\item Relation d'Einstein : $E_\ell = \Delta m \cdot c^{2}$. Analyse dimensionnelle : $[\Delta m \cdot c^{2}] = \mathrm{kg}\cdot(\mathrm{m\cdot s^{-1}})^{2} = \mathrm{kg\cdot m^{2}\cdot s^{-2}} = \mathrm{J}$, qui est bien une énergie. La relation est homogène.
\end{enumerate}
\end{corrige}

\begin{corrige}[Exercice 4 — Calculs directs]
\begin{enumerate}
\item Conversion : $t_{1/2} = \num{5,3} \times \num{3,156e7} = \num{1,67e8}\ \mathrm{s}$. Alors
\[ \lambda = \frac{\ln 2}{t_{1/2}} = \frac{\num{0,693}}{\num{1,67e8}} \approx \num{4,1e-9}\ \mathrm{s^{-1}}. \]
\item $A = \lambda N = \num{4,0e-9} \times \num{2,0e15} = \num{8,0e6}\ \mathrm{Bq}$, soit $\SI{8,0}{MBq}$.
\item Conversion de la masse : $\Delta m = \num{0,0304} \times \num{1,66e-27} = \num{5,05e-29}\ \mathrm{kg}$.
\[ E_\ell = \Delta m\, c^{2} = \num{5,05e-29} \times (\num{3,00e8})^{2} \approx \num{4,54e-12}\ \mathrm{J}. \]
En MeV : $E_\ell = \dfrac{\num{4,54e-12}}{\num{1,60e-19}} = \num{2,84e7}\ \mathrm{eV} \approx \num{28,4}\ \mathrm{MeV}$. \\
Énergie de liaison par nucléon ($A = 4$) : $\dfrac{E_\ell}{A} = \dfrac{\num{28,4}}{4} \approx \num{7,1}\ \mathrm{MeV/nucl\'eon}$.
\end{enumerate}
\end{corrige}
\begin{attention}[Points de vigilance]
Toujours convertir la période en secondes avant de calculer $\lambda$ en $\mathrm{s^{-1}}$, et convertir le défaut de masse en kg avant d'appliquer $E = \Delta m\,c^{2}$ si le résultat est demandé en joules.
\end{attention}

\courssub{Je m'entraîne}

\begin{corrige}[Exercice 5 — Équations de désintégration et lois de Soddy]
\begin{enumerate}
\item Conservation : $A : 210 = 206 + 4$ et $Z : 84 = 82 + 2$. Particule $\ce{^{4}_{2}He}$ : radioactivité \cle{$\alpha$}.
\[ \ce{^{210}_{84}Po -> ^{206}_{82}Pb + ^{4}_{2}He} \]
\item $A : 14 = 14 + 0$ ; $Z : 6 = 7 - 1$. Noyau fils $\ce{^{14}_{7}N}$ : radioactivité \cle{$\beta^-$}.
\[ \ce{^{14}_{6}C -> ^{14}_{7}N + ^{0}_{-1}e} \]
\item $A : 22 = 22 + 0$ ; $Z : 11 = 10 + 1$. Particule $\ce{^{0}_{1}e}$ (positon) : radioactivité \cle{$\beta^+$}.
\[ \ce{^{22}_{11}Na -> ^{22}_{10}Ne + ^{0}_{1}e} \]
\item $A : 137 = 137 + 0$ ; $Z : 55 = 56 - 1$. Noyau fils $\ce{^{137m}_{56}Ba}$ (état excité) : radioactivité \cle{$\beta^-$}, suivie d'une désexcitation \cle{$\gamma$} :
\[ \ce{^{137}_{55}Cs -> ^{137m}_{56}Ba + ^{0}_{-1}e} \qquad \ce{^{137m}_{56}Ba -> ^{137}_{56}Ba + ^{0}_{0}\gamma} \]
\item $A : 226 = 222 + 4$ ; $Z : 88 = 86 + 2$. Noyau fils $\ce{^{222}_{86}Rn}$ : radioactivité \cle{$\alpha$}.
\[ \ce{^{226}_{88}Ra -> ^{222}_{86}Rn + ^{4}_{2}He} \]
\end{enumerate}
\end{corrige}
\begin{attention}[Points de vigilance]
Citer explicitement les deux lois de Soddy (conservation de $A$ et de $Z$) avant d'écrire l'équation : c'est la justification attendue. Ne pas confondre $\beta^-$ (électron, $Z$ augmente de 1) et $\beta^+$ (positon, $Z$ diminue de 1).
\end{attention}

\begin{corrige}[Exercice 6 — Énergie libérée par la désintégration $\alpha$ de l'uranium 238]
\begin{enumerate}
\item La perte de masse est la différence entre la masse du noyau père et la somme des masses des produits :
\[ \Delta m = m(\ce{^{238}U}) - \left[ m(\ce{^{234}Th}) + m(\ce{^{4}He}) \right] = \num{238,0508} - (\num{234,0436} + \num{4,0026}) = \num{0,0046}\ \mathrm{u}. \]
\item Équivalence masse-énergie :
\[ E = \Delta m \cdot c^{2} = \num{0,0046} \times \num{931,5} \approx \num{4,3}\ \mathrm{MeV}. \]
En joules : $E = \num{4,3e6} \times \num{1,60e-19} \approx \num{6,9e-13}\ \mathrm{J}$.
\item $\Delta m > 0$ : la masse des produits est inférieure à celle du noyau père. Le système perd de la masse, donc de l'énergie : celle-ci est \cle{libérée} vers l'extérieur, essentiellement sous forme d'énergie cinétique de la particule $\alpha$ (et du noyau fils par recul). C'est cette énergie qui rend la désintégration spontanée possible.
\end{enumerate}
\end{corrige}
\begin{attention}[Points de vigilance]
Le bilan se fait toujours « masse des réactifs moins masse des produits » : un $\Delta m$ positif signifie une énergie libérée. Avec $\SI{1}{u} = \SI{931,5}{MeV/c^{2}}$, le produit $\Delta m \times 931,5$ donne directement $E$ en MeV, sans passer par les joules.
\end{attention}

\begin{corrige}[Exercice 7 — Rejet d'un déchet hospitalier à l'iode 131]
\begin{enumerate}
\item $\lambda = \dfrac{\ln 2}{t_{1/2}} = \dfrac{\num{0,693}}{\num{8,0}} \approx \num{8,7e-2}\ \mathrm{jour^{-1}}$.
\item L'activité suit la même loi exponentielle que $N$ :
\[ A(t) = A_0\,e^{-\lambda t} = \num{4,0e8}\,e^{-\num{0,0866}\,t}\ \mathrm{(Bq,\ t\ en\ jours)}. \]
\item On cherche $t$ tel que $A(t) = A_{\mathrm{seuil}}$ :
\[ e^{-\lambda t} = \frac{A_{\mathrm{seuil}}}{A_0} = \frac{\num{1,0e7}}{\num{4,0e8}} = \frac{1}{40} \quad\Longrightarrow\quad t = \frac{\ln 40}{\lambda} = \frac{\num{3,69}}{\num{0,0866}} \approx \num{43}\ \mathrm{jours}. \]
En nombre de périodes : $\dfrac{t}{t_{1/2}} = \dfrac{\num{43}}{\num{8,0}} \approx \num{5,3}$ périodes. Le flacon doit donc être stocké au moins \textbf{43 jours} (environ 6 semaines) en décroissance avant rejet.
\end{enumerate}
\end{corrige}
\begin{attention}[Points de vigilance]
Ici le temps est exprimé en jours : $\lambda$ doit donc être en $\mathrm{jour^{-1}}$, cohérent avec l'unité de $t$. Ne pas mélanger jours et secondes dans le même calcul.
\end{attention}

\begin{corrige}[Exercice 8 — Courbe d'Aston et stabilité des noyaux]
\begin{enumerate}
\item $E_\ell/A$ est l'énergie qu'il faut fournir, en moyenne, pour extraire un nucléon du noyau. Plus $E_\ell/A$ est grande, plus le noyau est \cle{stable}.
\item Ordre de stabilité croissante : $\ce{^{2}H}$ ($\num{1,1}\ \mathrm{MeV}$) $<$ $\ce{^{235}U}$ ($\num{7,6}\ \mathrm{MeV}$) $<$ $\ce{^{56}Fe}$ ($\num{8,8}\ \mathrm{MeV}$). Le fer 56 est au voisinage du maximum de stabilité.
\item Sur la courbe d'Aston (tracée en $-E_\ell/A$ en fonction de $A$), les noyaux les plus stables occupent le \emph{creux} de la courbe autour du fer. La \cle{fission} d'un noyau lourd ($\ce{^{235}U}$, $A$ grand) produit des noyaux de masse moyenne, plus proches du creux, donc plus stables : la différence d'énergie de liaison est libérée. De même, la \cle{fusion} de noyaux très légers ($\ce{^{2}H}$, $A$ petit) produit un noyau plus lourd situé plus bas sur la courbe, donc plus stable : là encore, de l'énergie est libérée. Dans les deux cas, les produits sont plus liés que les réactifs.
\item Allure attendue : courbe en $-E_\ell/A$ (MeV) en fonction de $A$, décroissant rapidement depuis $\ce{^{2}H}$ ($-\num{1,1}$), passant par un minimum voisin de $-\num{8,8}$ autour de $A = 56$ (fer), puis remontant lentement vers $-\num{7,6}$ pour $\ce{^{235}U}$ ($A = 235$).
\end{enumerate}
\end{corrige}

\courssub{Je me prépare au Bac}

\begin{corrige}[Exercice 9 — Type Bac : le radium 226 de Marie Curie]
\begin{enumerate}
\item La radioactivité \cle{$\alpha$} est l'émission spontanée, par un noyau instable, d'un noyau d'hélium $\ce{^{4}_{2}He}$ (particule $\alpha$). Équation :
\[ \ce{^{226}_{88}Ra -> ^{222}_{86}Rn + ^{4}_{2}He} \]
Lois de Soddy vérifiées : conservation du nombre de masse ($226 = 222 + 4$) et conservation du nombre de charge ($88 = 86 + 2$).
\item Conversion : $t_{1/2} = \num{1600} \times \num{3,156e7} = \num{5,05e10}\ \mathrm{s}$.
\[ \lambda = \frac{\ln 2}{t_{1/2}} = \frac{\num{0,693}}{\num{5,05e10}} \approx \num{1,37e-11}\ \mathrm{s^{-1}}. \]
\item Quantité de matière : $n = \dfrac{m}{M} = \dfrac{\num{1,00}}{\num{226}} = \num{4,42e-3}\ \mathrm{mol}$.
\[ N_0 = n\,\mathcal{N}_A = \num{4,42e-3} \times \num{6,02e23} \approx \num{2,66e21}\ \text{noyaux}. \]
\item $A_0 = \lambda N_0 = \num{1,37e-11} \times \num{2,66e21} \approx \num{3,7e10}\ \mathrm{Bq}$. \\
On retrouve quasi exactement $\SI{1}{Ci}$ : le curie a justement été défini historiquement comme l'activité d'un gramme de radium 226.
\item Rapport : $\dfrac{t}{t_{1/2}} = \dfrac{\num{4800}}{\num{1600}} = 3$ périodes. L'activité est donc divisée par $2^{3} = 8$ :
\[ A = \frac{A_0}{8} = \frac{\num{3,7e10}}{8} \approx \num{4,6e9}\ \mathrm{Bq}. \]
\item On veut $A = A_0/1000$, soit $e^{-\lambda t} = 10^{-3}$ :
\[ t = \frac{\ln 1000}{\lambda} = \frac{3\ln 10}{\ln 2}\,t_{1/2} \approx \num{9,97} \times \num{1600} \approx \num{1,6e4}\ \text{ans}. \]
Il faut environ \textbf{16 000 ans}, soit près de 10 périodes : c'est pourquoi les déchets radioactifs à vie longue posent un problème de stockage durable.
\end{enumerate}
\textbf{Barème indicatif (6 pts)} : définition + équation (1 pt) ; $\lambda$ avec conversion (1 pt) ; $N_0$ (1 pt) ; $A_0$ et comparaison au curie (1 pt) ; activité à 4800 ans (1 pt) ; durée pour $A_0/1000$ (1 pt).
\end{corrige}
\begin{attention}[Points de vigilance]
\begin{itemize}
\item Justifier l'équation par les deux lois de conservation de Soddy, pas seulement l'écrire.
\item Convertir $t_{1/2}$ en secondes pour obtenir $\lambda$ en $\mathrm{s^{-1}}$.
\item Quand $t$ est un multiple entier de $t_{1/2}$, privilégier la division par $2^{n}$ : c'est plus rapide et plus sûr que l'exponentielle.
\end{itemize}
\end{attention}

\begin{corrige}[Exercice 10 — Type Bac : datation au carbone 14 d'un ossement]
\begin{enumerate}
\item Radioactivité $\beta^-$ : émission d'un électron $\ce{^{0}_{-1}e}$ avec conservation de $A$ et de $Z$ :
\[ \ce{^{14}_{6}C -> ^{14}_{7}N + ^{0}_{-1}e} \]
Vérification : $14 = 14 + 0$ et $6 = 7 - 1$.
\item De $A(t) = A_0\,e^{-\lambda t}$ avec $\lambda = \dfrac{\ln 2}{t_{1/2}}$, on tire :
\[ t = \frac{1}{\lambda}\ln\!\left(\frac{A_0}{A(t)}\right) = \frac{t_{1/2}}{\ln 2}\,\ln\!\left(\frac{A_0}{A(t)}\right). \]
\item $\dfrac{A(t)}{A_0} = \num{12,5}\ \% = \dfrac{1}{8} = \left(\dfrac{1}{2}\right)^{3}$ : l'activité a été divisée par 2 exactement \textbf{3 fois}, donc $t = 3\,t_{1/2}$.
\[ t = 3 \times \num{5730} = \num{17190}\ \text{ans} \approx \num{1,72e4}\ \text{ans}. \]
\item Rapport des activités : $\dfrac{A(t)}{A_0} = \dfrac{\num{3,4}}{\num{13,6}} = \dfrac{1}{4} = \left(\dfrac{1}{2}\right)^{2}$, soit 2 périodes :
\[ t = 2 \times \num{5730} = \num{11460}\ \text{ans} \approx \num{1,15e4}\ \text{ans}. \]
\item Après 10 périodes : $\dfrac{A}{A_0} = \left(\dfrac{1}{2}\right)^{10} = \dfrac{1}{1024} \approx \num{9,8e-4}$, soit environ $\num{0,1}\ \%$. L'activité résiduelle est alors si faible (ici environ $\num{0,013}$ désintégration par minute et par gramme) qu'elle devient comparable au bruit de fond radioactif naturel et aux incertitudes de mesure : l'âge obtenu n'est plus fiable.
\end{enumerate}
\textbf{Barème indicatif (5 pts)} : équation $\beta^-$ justifiée (1 pt) ; établissement de l'expression de $t$ (1 pt) ; âge du premier ossement (1 pt) ; âge du second ossement (1 pt) ; activité résiduelle et limite de la méthode (1 pt).
\end{corrige}
\begin{attention}[Points de vigilance]
\begin{itemize}
\item L'hypothèse fondamentale de la datation doit être citée : l'activité initiale $A_0$ de l'échantillon est égale à celle d'un organisme vivant actuel, car les échanges avec l'atmosphère cessent à la mort.
\item Reconnaître les fractions remarquables : $\num{12,5}\ \% = 1/8 = (1/2)^{3}$ et $\num{25}\ \% = 1/4 = (1/2)^{2}$ évite tout calcul de logarithme.
\end{itemize}
\end{attention}

\begin{corrige}[Exercice 11 — Type Bac : détermination expérimentale d'une période radioactive]
\begin{enumerate}
\item Loi de décroissance : $A(t) = A_0\,e^{-\lambda t}$. En prenant le logarithme népérien :
\[ \ln A = \ln A_0 - \lambda t. \]
C'est bien une fonction \cle{affine} de $t$, de pente $-\lambda$ et d'ordonnée à l'origine $\ln A_0$.
\item Tableau complété :

\begin{center}
\begin{tabular}{|c|cccccc|}
\hline
$t$ (jours) & 0 & 5 & 10 & 15 & 20 & 25 \\
\hline
$A$ ($10^{6}$ Bq) & \num{8,00} & \num{5,04} & \num{3,17} & \num{2,00} & \num{1,26} & \num{0,79} \\
\hline
$\ln A$ & \num{2,079} & \num{1,617} & \num{1,154} & \num{0,693} & \num{0,231} & \num{-0,236} \\
\hline
\end{tabular}
\end{center}

\item Les six points sont alignés sur une droite décroissante : la loi exponentielle est validée expérimentalement.
\item Pente graphique, avec les points extrêmes $(0\ ;\ \num{2,079})$ et $(25\ ;\ \num{-0,236})$ :
\[ \text{pente} = \frac{\num{-0,236} - \num{2,079}}{25 - 0} = \frac{\num{-2,315}}{25} \approx \num{-9,3e-2}\ \mathrm{jour^{-1}}. \]
Donc $\lambda \approx \num{9,3e-2}\ \mathrm{jour^{-1}}$. Conversion ($1$ jour $= \num{8,64e4}\ \mathrm{s}$) :
\[ \lambda = \frac{\num{9,3e-2}}{\num{8,64e4}} \approx \num{1,1e-6}\ \mathrm{s^{-1}}. \]
\item Période radioactive :
\[ t_{1/2} = \frac{\ln 2}{\lambda} = \frac{\num{0,693}}{\num{9,3e-2}} \approx \num{7,5}\ \mathrm{jours}. \]
La valeur la plus proche dans la table est celle de l'\cle{iode 131} ($\num{8,0}$ jours) : l'écart s'explique par les incertitudes de mesure et de lecture graphique. Le nucléide est l'iode 131.
\item Prolongement graphique : à $t = \num{30}$ jours, on lit $\ln A \approx \num{-0,70}$, d'où $A \approx e^{-\num{0,70}} \approx \num{0,50}$, soit $A \approx \num{5,0e5}\ \mathrm{Bq}$. \\
Vérification par le calcul :
\[ A(30) = \num{8,00e6}\,e^{-\num{0,0926} \times 30} = \num{8,00e6}\,e^{-\num{2,78}} \approx \num{8,00e6} \times \num{0,062} \approx \num{5,0e5}\ \mathrm{Bq}. \]
Les deux résultats concordent.
\end{enumerate}
\textbf{Barème indicatif (6 pts)} : linéarisation $\ln A = \ln A_0 - \lambda t$ (1 pt) ; tableau complété (0,5 pt) ; trace soigné et alignement constaté (1 pt) ; pente et $\lambda$ dans les deux unités (1,5 pt) ; $t_{1/2}$ et identification (1 pt) ; détermination graphique et vérification (1 pt).
\end{corrige}
\begin{attention}[Points de vigilance]
\begin{itemize}
\item La pente est \emph{négative} : c'est $-\lambda$. Ne pas oublier le signe moins en écrivant $\lambda = -\text{pente}$.
\item La pente s'exprime en $\mathrm{jour^{-1}}$ car $t$ est en jours : la conversion en $\mathrm{s^{-1}}$ est exigée séparément.
\item Pour l'identification, comparer $t_{1/2}$ à la table et commenter l'écart par les incertitudes expérimentales : c'est la démarche scientifique attendue.
\end{itemize}
\end{attention}

\newpage

\coursec{Fiche bilan}

\begin{popfigure}
\begin{tikzpicture}[mindmap, grow cyclic, every node/.style={concept, text=white, font=\small},
  root concept/.append style={concept color=popdark, font=\bfseries},
  level 1/.append style={level distance=3.4cm, sibling angle=60},
  level 1 concept/.append style={font=\footnotesize\bfseries},
  level 2/.append style={level distance=2.2cm, sibling angle=45},
  level 2 concept/.append style={font=\scriptsize}]
\node[root concept]{Radioactivité et loi de décroissance}
  child[concept color=popblue]{ node {Noyau et cohésion}
    child { node {nucléide \ce{^{A}_{Z}X}, isotopes} }
    child { node {$E_\ell = \Delta m\, c^2$} }
    child { node {courbe d'Aston} }
  }
  child[concept color=popgreen]{ node {Définition et propriétés}
    child { node {spontanée, aléatoire, inévitable} }
  }
  child[concept color=poporange]{ node {Types et lois de Soddy}
    child { node {$\alpha$, $\beta^-$, $\beta^+$, $\gamma$} }
    child { node {conservation de $A$ et $Z$} }
  }
  child[concept color=poppink]{ node {Loi de décroissance}
    child { node {$N = N_0 e^{-\lambda t}$} }
    child { node {$t_{1/2} = \dfrac{\ln 2}{\lambda}$} }
  }
  child[concept color=poppurple]{ node {Activité}
    child { node {$A = \lambda N$ (Bq)} }
    child { node {$A = A_0 e^{-\lambda t}$} }
  }
  child[concept color=popgold]{ node {Énergie et applications}
    child { node {$Q = [m_{\text{avant}} - m_{\text{après}}]c^2$} }
    child { node {datation \ce{^{14}C}, médecine} }
  };
\end{tikzpicture}
\legende{Carte mentale de la leçon : les six grandes idées à maîtriser.}
\end{popfigure}

\begin{aretenir}[L'essentiel de la leçon]
\textbf{Définitions clés}
\begin{itemize}
  \item \cle{Nucléide} \ce{^{A}_{Z}X} : noyau constitué de $Z$ protons et de $A - Z$ neutrons ; $A$ est le nombre de masse (nombre de nucléons), $Z$ le numéro atomique (nombre de charges). Des \cle{isotopes} sont des nucléides de même $Z$ mais de $A$ différents (même élément chimique, nombre de neutrons différent).
  \item \cle{Défaut de masse} : $\Delta m = Z\,m_p + (A-Z)\,m_n - m(\ce{^{A}_{Z}X}) > 0$. La masse du noyau est toujours inférieure à la somme des masses de ses nucléons séparés au repos.
  \item \cle{Énergie de liaison} : $E_\ell = \Delta m \cdot c^2$ ; c'est l'énergie qu'il faut fournir au noyau au repos pour le dissocier en ses nucléons séparés au repos. L'énergie de liaison par nucléon $\dfrac{E_\ell}{A}$ (en MeV/nucléon) mesure la stabilité : plus $\dfrac{E_\ell}{A}$ est grande, plus le noyau est stable (maximum vers le fer, $\approx 8{,}8$ MeV/nucléon).
  \item \cle{Radioactivité} : transformation spontanée, aléatoire et inévitable d'un noyau instable en un noyau plus stable, accompagnée de l'émission de rayonnements (particules $\alpha$, $\beta$ ou photons $\gamma$).
  \item \cle{Constante radioactive} $\lambda$ (en $\mathrm{s}^{-1}$) : probabilité de désintégration d'un noyau par unité de temps ; elle est caractéristique du nucléide et indépendante des conditions extérieures (température, pression, combinaison chimique).
  \item \cle{Période radioactive (demi-vie)} $t_{1/2}$ : durée au bout de laquelle la moitié des noyaux radioactifs initialement présents s'est désintégrée.
  \item \cle{Activité} $A$ d'un échantillon : nombre moyen de désintégrations par seconde ; unité SI : le becquerel ($1~\mathrm{Bq} = 1$ désintégration par seconde). Ancienne unité : le curie, $1~\mathrm{Ci} = 3{,}7 \times 10^{10}~\mathrm{Bq}$.
\end{itemize}

\textbf{Formules à connaître par cœur}
\begin{center}
\begin{tabularx}{\linewidth}{|l|X|}\hline
\rowcolor{popblueL} \textbf{Grandeur} & \textbf{Relation} \\ \hline
Loi de décroissance & $N(t) = N_0\, e^{-\lambda t}$ \\ \hline
Demi-vie & $t_{1/2} = \dfrac{\ln 2}{\lambda}$ \quad (soit $\lambda = \dfrac{\ln 2}{t_{1/2}}$) \\ \hline
Activité & $A(t) = \lambda N(t) = A_0\, e^{-\lambda t}$, avec $A_0 = \lambda N_0$ \\ \hline
Nombre de noyaux & $N_0 = \dfrac{m_0}{M}\, \mathcal{N}_A$ \quad ($m_0$ masse de l'échantillon, $M$ masse molaire) \\ \hline
Énergie de liaison & $E_\ell = \Delta m \cdot c^2$ \\ \hline
Bilan d'énergie & $Q = \left[m_{\text{avant}} - m_{\text{après}}\right] c^2$ ; \quad $1~\mathrm{u} \cdot c^2 = 931{,}5~\mathrm{MeV}$ \\ \hline
\end{tabularx}
\end{center}

\textbf{Équations de désintégration (lois de Soddy : conservation du nombre de nucléons $A$ et de la charge $Z$)}
\begin{itemize}
  \item $\alpha$ : \ce{^{A}_{Z}X -> ^{A-4}_{Z-2}Y + ^{4}_{2}He} \quad (noyaux lourds, en général $Z > 82$)
  \item $\beta^-$ : \ce{^{A}_{Z}X -> ^{A}_{Z+1}Y + ^{0}_{-1}e} \quad (noyaux avec excès de neutrons)
  \item $\beta^+$ : \ce{^{A}_{Z}X -> ^{A}_{Z-1}Y + ^{0}_{+1}e} \quad (noyaux avec excès de protons)
  \item $\gamma$ : \ce{^{A}_{Z}X^{*} -> ^{A}_{Z}X + \gamma} \quad (désexcitation d'un noyau fils formé dans un état excité ; $A$ et $Z$ inchangés)
\end{itemize}

\textbf{Méthodes express}
\begin{itemize}
  \item \cle{Datation} : $t = \dfrac{1}{\lambda}\ln\!\left(\dfrac{N_0}{N}\right) = \dfrac{t_{1/2}}{\ln 2}\ln\!\left(\dfrac{A_0}{A}\right)$.
  \item Après $n$ périodes, il reste $N = \dfrac{N_0}{2^n}$ noyaux (et $A = \dfrac{A_0}{2^n}$).
  \item Lecture graphique de $t_{1/2}$ : abscisse du point de la courbe où $N = N_0/2$ (ou $A = A_0/2$).
  \item Courbe d'Aston ($-E_\ell/A$ en fonction de $A$, ou $E_\ell/A$ selon la convention) : les noyaux les plus stables sont au voisinage du fer ($A \approx 56$) ; les noyaux lourds peuvent libérer de l'énergie par fission, les noyaux légers par fusion.
\end{itemize}
\end{aretenir}

\begin{attention}[Pièges fréquents au Bac]
\begin{itemize}
  \item Ne pas confondre $A$ (nombre de masse, sans unité) et $A(t)$ (activité, en Bq) : le contexte et les unités permettent de les distinguer.
  \item Dans $t_{1/2} = \ln 2 / \lambda$, si $t_{1/2}$ est donné en années, $\lambda$ s'exprime en $\mathrm{an}^{-1}$ : convertir en secondes seulement si l'activité est demandée en becquerels.
  \item Le bilan d'énergie s'écrit $Q = [m_{\text{avant}} - m_{\text{après}}]\,c^2$ : une désintégration spontanée dégage de l'énergie, donc $Q > 0$, ce qui impose $m_{\text{avant}} > m_{\text{après}}$. Vérifier ce signe est un excellent contrôle.
  \item Dans une équation $\beta^+$, le positon est \ce{^{0}_{+1}e} : la conservation de la charge donne $Z = (Z-1) + 1$, et non $Z+1$.
  \item L'émission $\gamma$ ne modifie ni $A$ ni $Z$ : elle accompagne souvent une désintégration $\alpha$ ou $\beta$ (désexcitation du noyau fils).
  \item Pour la datation au carbone 14, $A_0$ est l'activité de l'échantillon de matière vivante (constante au cours du temps chez l'organisme vivant) et $A$ l'activité mesurée aujourd'hui : on a toujours $A < A_0$, donc $t > 0$.
\end{itemize}
\end{attention}

\begin{aretenir}[Checklist avant le Bac]
\begin{itemize}
  \item $\square$ Je sais définir nucléide, isotope, défaut de masse et énergie de liaison.
  \item $\square$ Je sais calculer $E_\ell = \Delta m \cdot c^2$ et $E_\ell/A$, avec les bonnes conversions (u, MeV, J ; $1~\mathrm{u} \cdot c^2 = 931{,}5~\mathrm{MeV}$).
  \item $\square$ Je sais citer les trois propriétés de la radioactivité : spontanée, aléatoire, inévitable.
  \item $\square$ Je sais écrire et équilibrer une équation $\alpha$, $\beta^-$, $\beta^+$ ou $\gamma$ grâce aux lois de Soddy.
  \item $\square$ Je connais la nature des particules émises : \ce{^{4}_{2}He} (noyau d'hélium), \ce{^{0}_{-1}e} (électron), \ce{^{0}_{+1}e} (positon), photon $\gamma$.
  \item $\square$ Je sais écrire et utiliser $N = N_0 e^{-\lambda t}$ et $A = \lambda N = A_0 e^{-\lambda t}$.
  \item $\square$ Je sais passer de $\lambda$ à $t_{1/2}$ par $t_{1/2} = \ln 2 / \lambda$, avec des unités cohérentes.
  \item $\square$ Je sais déterminer graphiquement $t_{1/2}$ sur une courbe de décroissance.
  \item $\square$ Je sais calculer un âge par datation au carbone 14 : $t = \dfrac{1}{\lambda}\ln\!\left(\dfrac{A_0}{A}\right)$.
  \item $\square$ Je sais faire un bilan d'énergie $Q = \Delta m \cdot c^2$ et vérifier son signe.
  \item $\square$ Je vérifie toujours les unités (Bq, s, J ou MeV) et je donne un nombre raisonnable de chiffres significatifs.
  \item $\square$ Je sais citer des applications : datation, médecine (radiothérapie, scintigraphie), traceurs.
\end{itemize}
\end{aretenir}

\findecours

\end{document}
